\chapter{Chapter07 — Unique Games Reduction}
%%% generated by the blueprint pipeline from FABL.Chapter07.UniqueGamesReduction
\section{Overview}

Declarations of the Lean module \texttt{FABL.Chapter07.UniqueGamesReduction}.

\section{Declarations}

\begin{definition}[Dictator versus low influence tester]
\lean{FABL.IsDictatorVsNoInfluentialsTest}
\statementsource{Analysis-of-Boolean-Functions-by-Ryan-ODonnell}{pdf-ffbe9c7b3306-p191-b002}
\label{FABL.IsDictatorVsNoInfluentialsTest}
\leanok
\uses{FABL.BoundedCubeFunction, FABL.BoundedCubeFunction.ofBoolean, FABL.HasSmallInfluences, FABL.SignPredicate, FABL.SignTester, FABL.SignTester.UsesPredicates, FABL.SignTester.acceptanceProbability, FABL.dictator}
A randomized sign tester distinguishes dictators from Boolean functions with no
influential coordinate using a finite allowed predicate family if that family is finite,
$0<\alpha<\beta\le 1$, and the error function takes values in $[0,1]$ on $[0,1]$ and
tends to $0$ as its positive argument tends to $0$. For every positive input dimension,
the tester uses only predicates from that family and accepts every dictator with
probability at least $\beta$. For every $\varepsilon\in[0,1]$, it accepts each function
whose coordinate influences are all at most $\varepsilon$ with probability at most
$\alpha+\mathrm{error}(\varepsilon)$.
\end{definition}

\begin{theorem}[From No Notables to No Influentials]
\lean{FABL.IsDictatorVsNoNotablesTest.toNoInfluentials}
\label{FABL.IsDictatorVsNoNotablesTest.toNoInfluentials}
\leanok
\uses{FABL.HasSmallInfluences.hasSmallStableInfluences, FABL.IsDictatorVsNoInfluentialsTest, FABL.IsDictatorVsNoNotablesTest, FABL.SignPredicate, FABL.SignTester}
\difficulty{3}
Let $\Psi$ be a finite family of predicates on signs, let $\alpha,\beta\in\mathbb R$
satisfy $0<\alpha<\beta\le1$, and let $e:[0,1]\to[0,1]$ tend to $0$ as its argument
tends to $0$ through positive values. Suppose a family of randomized testers, one for
each positive cube dimension, uses only predicates in $\Psi$, accepts every coordinate
dictator with probability at least $\beta$, and accepts every function whose stable
influences at parameter $1-\varepsilon$ are all at most $\varepsilon$ with probability
at most $\alpha+e(\varepsilon)$ for each $\varepsilon\in[0,1]$. Then the same testers
use only predicates in $\Psi$, accept every dictator with probability at least $\beta$,
and accept every function whose influences at parameter $1$ are all at most
$\varepsilon$ with probability at most $\alpha+e(\varepsilon)$ for each
$\varepsilon\in[0,1]$.
\end{theorem}

\begin{definition}[Unique Games constraint template]
\lean{FABL.uniqueGamesTemplate}
\statementsource{Analysis-of-Boolean-Functions-by-Ryan-ODonnell}{pdf-ffbe9c7b3306-p200-b007}
\label{FABL.uniqueGamesTemplate}
\leanok
\uses{FABL.CSPTemplate}
For a natural number $q$, the Unique Games template has domain $\{0,\ldots,q-1\}$ and
predicates indexed by permutations of that domain. Each predicate has arity two and
accepts an ordered pair $(a,b)$ exactly when the permutation sends $b$ to $a$.
\end{definition}

\begin{definition}[Binary Scope Injection]
\lean{FABL.binaryInjectiveScope}
\label{FABL.binaryInjectiveScope}
\leanok
Let $V$ be a type, and let $u,v\in V$ be distinct. The map from $in 2$ to $V$ sending
$0$ to $u$ and $1$ to $v$ is an injection.
\end{definition}

\begin{theorem}[Value of the binary injection at zero]
\lean{FABL.binaryInjectiveScope_zero}
\label{FABL.binaryInjectiveScope_zero}
\leanok
\uses{FABL.binaryInjectiveScope}
\difficulty{1}
Let $u$ and $v$ be distinct elements of a type $V$. The injection from the two-element
set into $V$ that sends its first element to $u$ and its second element to $v$ takes the
value $u$ at the first element.
\end{theorem}

\begin{theorem}[Value of the binary injection at one]
\lean{FABL.binaryInjectiveScope_one}
\label{FABL.binaryInjectiveScope_one}
\leanok
\uses{FABL.binaryInjectiveScope}
\difficulty{1}
Let $u$ and $v$ be distinct elements of a type $V$. The injection from the two-element
set into $V$ that sends its first element to $u$ and its second element to $v$ takes the
value $v$ at the second element.
\end{theorem}

\begin{definition}[Finite Unique Games instance]
\lean{FABL.FiniteUniqueGame}
\statementsource{Analysis-of-Boolean-Functions-by-Ryan-ODonnell}{pdf-ffbe9c7b3306-p200-b007}
\label{FABL.FiniteUniqueGame}
\leanok
A finite Unique Games instance over a domain of size $q$ consists of a finite nonempty
set of variables and a finite nonempty set of constraint occurrences. Each occurrence
has an ordered pair of distinct endpoints and is labeled by a permutation of the domain,
and every variable is an endpoint of at least one occurrence.
\end{definition}

\begin{definition}[Constraint associated to an occurrence]
\lean{FABL.FiniteUniqueGame.constraint}
\statementsource{Analysis-of-Boolean-Functions-by-Ryan-ODonnell}{pdf-ffbe9c7b3306-p200-b007}
\label{FABL.FiniteUniqueGame.constraint}
\leanok
\uses{FABL.CSPConstraint, FABL.binaryInjectiveScope, FABL.uniqueGamesTemplate}
Given a finite Unique Games instance and one of its constraint occurrences, the
associated binary constraint uses the permutation labeling that occurrence and its
ordered pair of endpoints as its scope.
\end{definition}

\begin{theorem}[Constraint satisfaction criterion]
\lean{FABL.FiniteUniqueGame.constraint_satisfied_iff}
\label{FABL.FiniteUniqueGame.constraint_satisfied_iff}
\leanok
\uses{FABL.CSPConstraint.Satisfied, FABL.FiniteUniqueGame.constraint}
\difficulty{3}
Let $G$ be a finite unique game, let $e$ be a constraint with source $u$, target $v$,
and attached permutation $\pi_e$, and assign a label from the $q$-element label set to
every vertex. The constraint is satisfied if and only if $\pi_e$ sends the label of $v$
to the label of $u$.
\end{theorem}

\begin{definition}[CSP instance from a Unique Games instance]
\lean{FABL.FiniteUniqueGame.toCSPInstance}
\statementsource{Analysis-of-Boolean-Functions-by-Ryan-ODonnell}{
  pdf-ffbe9c7b3306-p184-b004,
  pdf-ffbe9c7b3306-p185-b004,
  pdf-ffbe9c7b3306-p200-b007
}
\label{FABL.FiniteUniqueGame.toCSPInstance}
\leanok
\uses{FABL.CSPInstance, FABL.FiniteUniqueGame.constraint, FABL.binaryInjectiveScope, FABL.uniqueGamesTemplate}
Every finite Unique Games instance determines a constraint satisfaction instance over
its Unique Games template on the same variable set: its constraints are the multiset of
constraints associated to all occurrences, retaining multiplicities. This multiset is
nonempty, and every declared variable appears in the scope of at least one constraint.
\end{definition}

\begin{definition}[Fully satisfying a Unique Games instance]
\lean{FABL.FiniteUniqueGame.FullySatisfied}
\statementsource{Analysis-of-Boolean-Functions-by-Ryan-ODonnell}{pdf-ffbe9c7b3306-p186-b001}
\label{FABL.FiniteUniqueGame.FullySatisfied}
\leanok
A labeling of the variables by elements of the domain fully satisfies a finite Unique
Games instance if, for every constraint occurrence, applying its permutation to the
label of its target gives the label of its source.
\end{definition}

\begin{definition}[Satisfiability of a Unique Games instance]
\lean{FABL.FiniteUniqueGame.Satisfiable}
\statementsource{Analysis-of-Boolean-Functions-by-Ryan-ODonnell}{pdf-ffbe9c7b3306-p186-b001}
\label{FABL.FiniteUniqueGame.Satisfiable}
\leanok
\uses{FABL.FiniteUniqueGame.FullySatisfied}
A finite Unique Games instance is satisfiable if there exists a labeling of its
variables by elements of the domain that fully satisfies every constraint occurrence.
\end{definition}

\begin{theorem}[Equivalence of full satisfaction formulations]
\lean{FABL.FiniteUniqueGame.fullySatisfied_iff_csp}
\label{FABL.FiniteUniqueGame.fullySatisfied_iff_csp}
\leanok
\uses{FABL.CSPInstance.FullySatisfied, FABL.FiniteUniqueGame.FullySatisfied, FABL.FiniteUniqueGame.constraint, FABL.FiniteUniqueGame.constraint_satisfied_iff, FABL.FiniteUniqueGame.toCSPInstance}
\difficulty{4}
Let $G$ be a finite unique game and let $\ell$ assign one of the $q$ labels to each
vertex. Then every game constraint is satisfied by $\ell$ if and only if $\ell$
satisfies every constraint in the associated constraint satisfaction instance.
\end{theorem}

\begin{theorem}[Equivalent notions of satisfiability]
\lean{FABL.FiniteUniqueGame.satisfiable_iff_csp_satisfiable}
\label{FABL.FiniteUniqueGame.satisfiable_iff_csp_satisfiable}
\leanok
\uses{FABL.CSPInstance.Satisfiable, FABL.CSPInstance.satisfiable_iff_exists_fullySatisfied, FABL.FiniteUniqueGame.Satisfiable, FABL.FiniteUniqueGame.fullySatisfied_iff_csp, FABL.FiniteUniqueGame.toCSPInstance}
\difficulty{3}
Let $q$ be a positive integer and let $G$ be a finite unique game with a nonempty set of
variables and a nonempty set of constraints. Each constraint joins two distinct
variables, carries a permutation of the $q$ labels, and every variable occurs in some
constraint. Then $G$ has a labeling satisfying every constraint if and only if the
associated constraint satisfaction instance has optimum value $1$.
\end{theorem}

\begin{definition}[Walks in a Finite Unique Games Instance]
\lean{FABL.FiniteUniqueGame.Walk}
\label{FABL.FiniteUniqueGame.Walk}
\leanok
\uses{FABL.FiniteUniqueGame}
Let $G$ be a finite Unique Games instance with a nonempty finite set of variables and a
nonempty finite set of constraint occurrences. A walk from one variable to another is
generated by the following steps: a length-zero walk at any variable; traversal of a
constraint from its source to its target or in the reverse direction; and concatenation
of two walks whose endpoints match.
\end{definition}

\begin{definition}[Reversal of a Walk]
\lean{FABL.FiniteUniqueGame.Walk.reverse}
\label{FABL.FiniteUniqueGame.Walk.reverse}
\leanok
\uses{FABL.FiniteUniqueGame.Walk}
In a finite Unique Games instance, every walk from a variable $x$ to a variable $y$
determines a walk from $y$ to $x$: a length-zero walk stays unchanged, each constraint
traversal is reversed, and the order of concatenated walks is reversed.
\end{definition}

\begin{definition}[Label Transport Along a Walk]
\lean{FABL.FiniteUniqueGame.Walk.transport}
\label{FABL.FiniteUniqueGame.Walk.transport}
\leanok
\uses{FABL.FiniteUniqueGame.Walk}
In a finite Unique Games instance, a walk from one variable to another transports labels
in $in q$. A length-zero walk leaves a label unchanged; traversing a constraint from
source to target applies the inverse of its permutation, traversing it backward applies
its permutation, and concatenated walks transport labels in the order in which the walks
are traversed.
\end{definition}

\begin{theorem}[Transport Along the Empty Walk]
\lean{FABL.FiniteUniqueGame.Walk.transport_nil}
\label{FABL.FiniteUniqueGame.Walk.transport_nil}
\leanok
\uses{FABL.FiniteUniqueGame.Walk.transport}
\difficulty{1}
Let $q$ be a natural number, let a finite unique game be given, and let $v$ be a vertex
and $a\in\{0,\ldots,q-1\}$. Transporting $a$ along the empty walk from $v$ to itself
leaves it unchanged.
\end{theorem}

\begin{theorem}[Label Transport Along a Walk]
\lean{FABL.FiniteUniqueGame.Walk.label_eq_transport}
\label{FABL.FiniteUniqueGame.Walk.label_eq_transport}
\leanok
\uses{FABL.FiniteUniqueGame.FullySatisfied, FABL.FiniteUniqueGame.Walk, FABL.FiniteUniqueGame.Walk.transport}
\difficulty{4}
Let $q$ be a natural number and let a finite unique game have a nonempty finite vertex
set and a nonempty finite set of constraint occurrences, each joining distinct vertices
and carrying a permutation of $\{0,\ldots,q-1\}$; assume every vertex occurs in some
constraint. If a labeling satisfies every constraint, then along any walk from a vertex
$u$ to a vertex $v$, the label at $v$ is the transport of the label at $u$, applying
each edge permutation in the direction prescribed by the walk.
\end{theorem}

\begin{definition}[Connectivity by Walks]
\lean{FABL.FiniteUniqueGame.Connected}
\label{FABL.FiniteUniqueGame.Connected}
\leanok
\uses{FABL.FiniteUniqueGame.Walk}
Let $G$ be a finite Unique Games instance. Two variables are connected if there exists a
walk in $G$ from the first variable to the second.
\end{definition}

\begin{definition}[Connectedness equivalence relation]
\lean{FABL.FiniteUniqueGame.connectedSetoid}
\label{FABL.FiniteUniqueGame.connectedSetoid}
\leanok
\uses{FABL.FiniteUniqueGame.Connected, FABL.FiniteUniqueGame.Walk.reverse}
Let $q$ be a nonnegative integer, and let $G$ be a finite unique game with a nonempty
finite set of variables and a nonempty finite set of constraint occurrences. Each
constraint joins two distinct variables and carries a permutation of the $q$ labels, and
every variable occurs in a constraint. Declare two variables equivalent when there is a
walk between them, where each constraint can be traversed in either direction.
\end{definition}

\begin{definition}[Connected components]
\lean{FABL.FiniteUniqueGame.ConnectedComponent}
\label{FABL.FiniteUniqueGame.ConnectedComponent}
\leanok
\uses{FABL.FiniteUniqueGame.connectedSetoid}
Let $q$ be a nonnegative integer, and let $G$ be a finite unique game with a nonempty
finite set of variables and a nonempty finite set of constraint occurrences. Each
constraint joins two distinct variables and carries a permutation of the $q$ labels, and
every variable occurs in a constraint. The connected components of $G$ are the
equivalence classes of variables under the relation that two variables are equivalent
exactly when a walk joins them, with each constraint traversable in either direction.
\end{definition}

\begin{definition}[Component of a variable]
\lean{FABL.FiniteUniqueGame.component}
\label{FABL.FiniteUniqueGame.component}
\leanok
\uses{FABL.FiniteUniqueGame.ConnectedComponent, FABL.FiniteUniqueGame.connectedSetoid}
Let $q$ be a nonnegative integer, and let $G$ be a finite unique game with a nonempty
finite set of variables and a nonempty finite set of constraint occurrences. Each
constraint joins two distinct variables and carries a permutation of the $q$ labels, and
every variable occurs in a constraint. The component of a variable is its equivalence
class under the relation that two variables are equivalent exactly when a walk joins
them, with each constraint traversable in either direction.
\end{definition}

\begin{definition}[Representative of a component]
\lean{FABL.FiniteUniqueGame.componentRoot}
\label{FABL.FiniteUniqueGame.componentRoot}
\leanok
\uses{FABL.FiniteUniqueGame.ConnectedComponent}
Let $q$ be a nonnegative integer, and let $G$ be a finite unique game with a nonempty
finite set of variables and a nonempty finite set of constraint occurrences. Each
constraint joins two distinct variables and carries a permutation of the $q$ labels, and
every variable occurs in a constraint. A representative of a connected component is a
variable chosen from that component.
\end{definition}

\begin{definition}[Path from component representative]
\lean{FABL.FiniteUniqueGame.componentPath}
\label{FABL.FiniteUniqueGame.componentPath}
\leanok
\uses{FABL.FiniteUniqueGame.Walk, FABL.FiniteUniqueGame.component, FABL.FiniteUniqueGame.componentRoot}
Let $q$ be a nonnegative integer, and let $G$ be a finite unique game with a nonempty
finite set of variables and a nonempty finite set of constraint occurrences. Each
constraint joins two distinct variables and carries a permutation of the $q$ labels, and
every variable occurs in a constraint. For each variable, the component representative
is joined to it by a walk, with each constraint traversable in either direction.
\end{definition}

\begin{definition}[Label propagated along a path]
\lean{FABL.FiniteUniqueGame.propagatedLabel}
\label{FABL.FiniteUniqueGame.propagatedLabel}
\leanok
\uses{FABL.FiniteUniqueGame.ConnectedComponent, FABL.FiniteUniqueGame.Walk.transport, FABL.FiniteUniqueGame.component, FABL.FiniteUniqueGame.componentPath}
Let $q$ be a nonnegative integer, and let $G$ be a finite unique game with a nonempty
finite set of variables and a nonempty finite set of constraint occurrences. Each
constraint joins two distinct variables and carries a permutation of the $q$ labels, and
every variable occurs in a constraint. Given a label in $	ext{Fin}(q)$ for each
connected component, assign to each variable the label obtained by transporting its
component's label from the component representative along a walk to that variable.
\end{definition}

\begin{theorem}[Propagation recovers a satisfying labeling]
\lean{FABL.FiniteUniqueGame.propagatedLabel_eq_of_fullySatisfied}
\label{FABL.FiniteUniqueGame.propagatedLabel_eq_of_fullySatisfied}
\leanok
\uses{FABL.FiniteUniqueGame.FullySatisfied, FABL.FiniteUniqueGame.Walk.label_eq_transport, FABL.FiniteUniqueGame.componentPath, FABL.FiniteUniqueGame.componentRoot, FABL.FiniteUniqueGame.propagatedLabel}
\difficulty{3}
Let $G$ be a finite unique game, and suppose a labeling of its vertices satisfies every
constraint. Assign to each component root the label it has in this labeling, then
propagate those root labels along the chosen paths. The resulting labeling is equal to
the original labeling at every vertex.
\end{theorem}

\begin{theorem}[Satisfiability via component labels]
\lean{FABL.FiniteUniqueGame.satisfiable_iff_exists_component_labels}
\label{FABL.FiniteUniqueGame.satisfiable_iff_exists_component_labels}
\leanok
\uses{FABL.FiniteUniqueGame.ConnectedComponent, FABL.FiniteUniqueGame.FullySatisfied, FABL.FiniteUniqueGame.Satisfiable, FABL.FiniteUniqueGame.componentRoot, FABL.FiniteUniqueGame.propagatedLabel, FABL.FiniteUniqueGame.propagatedLabel_eq_of_fullySatisfied}
\difficulty{4}
Let $q$ be a nonnegative integer and let $G$ be a finite unique game with a nonempty set
of variables and a nonempty set of constraints. Each constraint joins two distinct
variables, carries a permutation of the $q$ labels, and every variable occurs in some
constraint. The game has a labeling satisfying every constraint if and only if there is
a choice of one label for each connected component such that propagating those labels
along component paths gives a labeling satisfying every constraint.
\end{theorem}

\begin{theorem}[Endpoints lie in the same component]
\lean{FABL.FiniteUniqueGame.component_source_eq_target}
\label{FABL.FiniteUniqueGame.component_source_eq_target}
\leanok
\uses{FABL.FiniteUniqueGame.component}
\difficulty{3}
Let $G$ be a finite unique game, and let $e$ be any constraint with source vertex $u$
and target vertex $v$. Then $u$ and $v$ belong to the same connected component.
\end{theorem}

\begin{definition}[Valid label for a component representative]
\lean{FABL.FiniteUniqueGame.ComponentRootLabelValid}
\label{FABL.FiniteUniqueGame.ComponentRootLabelValid}
\leanok
\uses{FABL.FiniteUniqueGame.ConnectedComponent, FABL.FiniteUniqueGame.Walk.transport, FABL.FiniteUniqueGame.component, FABL.FiniteUniqueGame.componentPath}
Let $q$ be a nonnegative integer, and let $G$ be a finite unique game with a nonempty
finite set of variables and a nonempty finite set of constraint occurrences. Each
constraint joins two distinct variables and carries a permutation of the $q$ labels, and
every variable occurs in a constraint. A label in $	ext{Fin}(q)$ for a component
representative is valid if, for every constraint whose source lies in that component,
applying the constraint's permutation to the label transported from the target along its
component path gives the label transported from the source along its component path.
\end{definition}

\begin{theorem}[Satisfaction of propagated labels]
\lean{FABL.FiniteUniqueGame.fullySatisfied_propagatedLabel_iff}
\label{FABL.FiniteUniqueGame.fullySatisfied_propagatedLabel_iff}
\leanok
\uses{FABL.FiniteUniqueGame.ComponentRootLabelValid, FABL.FiniteUniqueGame.ConnectedComponent, FABL.FiniteUniqueGame.FullySatisfied, FABL.FiniteUniqueGame.component, FABL.FiniteUniqueGame.component_source_eq_target, FABL.FiniteUniqueGame.propagatedLabel}
\difficulty{4}
Let $G$ be a finite unique game, and choose a root label for each connected component.
Propagate these labels to vertices by transporting each component's root label along the
chosen paths from its root. The resulting vertex labeling satisfies every constraint if
and only if, in every component, its chosen root label satisfies the compatibility
condition for every constraint whose source lies in that component.
\end{theorem}

\begin{definition}[Valid labels for a component]
\lean{FABL.FiniteUniqueGame.validComponentRootLabels}
\label{FABL.FiniteUniqueGame.validComponentRootLabels}
\leanok
\uses{FABL.FiniteUniqueGame.ComponentRootLabelValid, FABL.FiniteUniqueGame.ConnectedComponent}
Let $q$ be a nonnegative integer, and let $G$ be a finite unique game with a nonempty
finite set of variables and a nonempty finite set of constraint occurrences. Each
constraint joins two distinct variables and carries a permutation of the $q$ labels, and
every variable occurs in a constraint. For a connected component, the valid labels are
the subset of $	ext{Fin}(q)$ consisting of labels for which every constraint whose
source lies in that component sends the label transported from its target to the label
transported from its source, using the paths from the component representative.
\end{definition}

\begin{theorem}[Membership in the valid root label set]
\lean{FABL.FiniteUniqueGame.mem_validComponentRootLabels}
\label{FABL.FiniteUniqueGame.mem_validComponentRootLabels}
\leanok
\uses{FABL.FiniteUniqueGame.ComponentRootLabelValid, FABL.FiniteUniqueGame.ConnectedComponent, FABL.FiniteUniqueGame.validComponentRootLabels}
\difficulty{1}
Let $G$ be a finite unique game, let $C$ be a connected component, and let $a$ be one of
the $q$ labels. Then $a$ belongs to the set of valid root labels for $C$ if and only if
transporting $a$ along the chosen paths from the component root satisfies the
compatibility condition for every constraint whose source lies in $C$.
\end{theorem}

\begin{theorem}[Bound on valid component root labels]
\lean{FABL.FiniteUniqueGame.card_validComponentRootLabels_le}
\label{FABL.FiniteUniqueGame.card_validComponentRootLabels_le}
\leanok
\uses{FABL.FiniteUniqueGame.ConnectedComponent, FABL.FiniteUniqueGame.validComponentRootLabels}
\difficulty{3}
Let $G$ be a finite unique game with nonempty finite vertex and constraint sets,
distinct endpoints for each constraint, and a permutation of the label set
$\\{0,\\ldots,q-1\\}$ attached to each constraint. For any connected component, the
number of labels at its chosen root that, when transported along the chosen paths from
that root, satisfy every constraint in the component is at most $q$.
\end{theorem}

\begin{definition}[Component label selection]
\lean{FABL.FiniteUniqueGame.optimalComponentLabels}
\label{FABL.FiniteUniqueGame.optimalComponentLabels}
\leanok
\uses{FABL.FiniteUniqueGame.ConnectedComponent, FABL.FiniteUniqueGame.validComponentRootLabels}
Let $G$ be a finite Unique Games instance with nonempty label set $	ext{Fin}(q)$, and
let $C$ be a connected component of its variables. If $C$ has any valid root labels,
choose one of them; otherwise choose an arbitrary label in $	ext{Fin}(q)$.
\end{definition}

\begin{theorem}[Validity of a selected component label]
\lean{FABL.FiniteUniqueGame.optimalComponentLabels_valid}
\label{FABL.FiniteUniqueGame.optimalComponentLabels_valid}
\leanok
\uses{FABL.FiniteUniqueGame.ComponentRootLabelValid, FABL.FiniteUniqueGame.ConnectedComponent, FABL.FiniteUniqueGame.mem_validComponentRootLabels, FABL.FiniteUniqueGame.optimalComponentLabels, FABL.FiniteUniqueGame.validComponentRootLabels}
\difficulty{4}
Let $G$ be a finite unique game with at least one available label, and fix a connected
component. If that component has at least one valid root label, then the selected
component label is valid.
\end{theorem}

\begin{definition}[Propagated vertex assignment]
\lean{FABL.FiniteUniqueGame.propagationAssignment}
\label{FABL.FiniteUniqueGame.propagationAssignment}
\leanok
\uses{FABL.FiniteUniqueGame.optimalComponentLabels, FABL.FiniteUniqueGame.propagatedLabel}
Let $G$ be a finite Unique Games instance with nonempty label set $	ext{Fin}(q)$. The
propagation assignment is the vertex labeling obtained by applying $G$'s label
propagation rule to its selected labels for connected components.
\end{definition}

\begin{theorem}[Satisfaction by propagated labels]
\lean{FABL.FiniteUniqueGame.propagationAssignment_fullySatisfied}
\label{FABL.FiniteUniqueGame.propagationAssignment_fullySatisfied}
\leanok
\uses{FABL.FiniteUniqueGame.FullySatisfied, FABL.FiniteUniqueGame.Satisfiable, FABL.FiniteUniqueGame.fullySatisfied_propagatedLabel_iff, FABL.FiniteUniqueGame.optimalComponentLabels, FABL.FiniteUniqueGame.optimalComponentLabels_valid, FABL.FiniteUniqueGame.propagationAssignment, FABL.FiniteUniqueGame.satisfiable_iff_exists_component_labels}
\difficulty{3}
Let $q$ be a positive integer and let $G$ be a finite unique game with a nonempty set of
variables and a nonempty set of constraints. Each constraint joins two distinct
variables, carries a permutation of the $q$ labels, and every variable occurs in some
constraint. If $G$ admits a labeling satisfying every constraint, then the labeling
obtained by propagating labels from component roots also satisfies every constraint.
\end{theorem}

\begin{theorem}[Optimal propagation has value one]
\lean{FABL.FiniteUniqueGame.propagationAssignment_one_one}
\label{FABL.FiniteUniqueGame.propagationAssignment_one_one}
\leanok
\uses{FABL.CSPInstance.optimum, FABL.CSPInstance.value, FABL.CSPInstance.value_eq_one_iff_fullySatisfied, FABL.FiniteUniqueGame.Satisfiable, FABL.FiniteUniqueGame.fullySatisfied_iff_csp, FABL.FiniteUniqueGame.propagationAssignment, FABL.FiniteUniqueGame.propagationAssignment_fullySatisfied, FABL.FiniteUniqueGame.satisfiable_iff_csp_satisfiable, FABL.FiniteUniqueGame.toCSPInstance}
\difficulty{3}
Let $q$ be a positive integer and let $G$ be a finite unique game with a nonempty set of
variables and a nonempty set of constraints. Each constraint joins two distinct
variables, carries a permutation of the $q$ labels, and every variable occurs in some
constraint. If the optimum value of the associated constraint satisfaction instance is
$1$, then the assignment obtained by propagating the selected labels from component
roots has value $1$.
\end{theorem}

\begin{definition}[Encoded family of Unique Games instances]
\lean{FABL.EncodedUniqueGamesFamily}
\statementsource{Analysis-of-Boolean-Functions-by-Ryan-ODonnell}{pdf-ffbe9c7b3306-p200-b007}
\label{FABL.EncodedUniqueGamesFamily}
\leanok
\uses{FABL.BinaryEncoding, FABL.FiniteUniqueGame}
An encoded family of Unique Games instances over a domain of size $q$ consists of a type
of semantic instance descriptions, a canonical finite binary encoding of those
descriptions, and an assignment of a finite Unique Games instance to each description.
\end{definition}

\begin{definition}[Propagation optimization encoding]
\lean{FABL.EncodedUniqueGamesFamily.PropagationOptimizationEncoding}
\label{FABL.EncodedUniqueGamesFamily.PropagationOptimizationEncoding}
\leanok
\uses{FABL.CSPInstance.optimum, FABL.CSPInstance.value, FABL.EncodedMaximizationProblem, FABL.FiniteUniqueGame.toCSPInstance}
Let $q$ be a positive number of labels, and let a family of encoded Unique Games
instances be given. A propagation optimization encoding consists of an encoded
maximization problem and, for each semantic instance and each vertex labeling of that
instance, a bit-list encoding of the labeling. Every well-formed input of the
maximization problem must encode some semantic instance, every canonical instance
encoding must be a well-formed input, and every encoded labeling must be feasible. On
each canonical instance encoding, the encoded objective value of every labeling must
equal its generic CSP value, and the encoded optimum must equal the generic CSP optimum.
\end{definition}

\begin{definition}[Propagation search machine]
\lean{FABL.EncodedUniqueGamesFamily.PropagationSearchMachine}
\label{FABL.EncodedUniqueGamesFamily.PropagationSearchMachine}
\leanok
\uses{FABL.EncodedUniqueGamesFamily.PropagationOptimizationEncoding, FABL.FiniteUniqueGame.propagationAssignment, FABL.PolytimeMap}
Let $q$ be a positive number of labels, let a family of encoded Unique Games instances
be given, and let a propagation optimization encoding for that family be fixed. A
propagation search machine consists of a polynomial-time map from bit-list inputs to
bit-list outputs such that, on every canonical instance encoding, its output is the
encoding of that instance's propagation assignment.
\end{definition}

\begin{theorem}[Propagation threshold approximation]
\lean{FABL.EncodedUniqueGamesFamily.propagationSearchMachine_isThresholdApproximation}
\label{FABL.EncodedUniqueGamesFamily.propagationSearchMachine_isThresholdApproximation}
\leanok
\uses{FABL.CSPInstance.optimum_le_one, FABL.EncodedUniqueGamesFamily.PropagationOptimizationEncoding, FABL.EncodedUniqueGamesFamily.PropagationSearchMachine, FABL.FiniteUniqueGame.propagationAssignment_one_one, FABL.FiniteUniqueGame.toCSPInstance, FABL.PolytimeMap.IsThresholdApproximation}
\difficulty{4}
Let $q>0$, and consider an encoded finite unique-games optimization problem with $q$
labels, together with a propagation optimization encoding and a polynomial-time machine
that outputs the encoding of the propagation assignment on every canonical game
instance. For every valid encoded input, the machine's output is feasible; moreover, if
the optimum is at least $1$, the output has value at least $1$.
\end{theorem}

\begin{definition}[Constraint template from a sign tester]
\lean{FABL.SignTester.cspTemplate}
\statementsource{Analysis-of-Boolean-Functions-by-Ryan-ODonnell}{
  pdf-ffbe9c7b3306-p184-b004,
  pdf-ffbe9c7b3306-p186-b001
}
\label{FABL.SignTester.cspTemplate}
\leanok
\uses{FABL.CSPTemplate, FABL.Sign, FABL.SignTester}
For a sign tester on a finite seed space, form a constraint template over the sign
domain whose predicates are the tester's seeds. The arity and acceptance rule of each
seed's predicate are those of the deterministic check selected by that seed.
\end{definition}

\begin{definition}[Template from a finite predicate family]
\lean{FABL.signPredicateSetTemplate}
\statementsource{Analysis-of-Boolean-Functions-by-Ryan-ODonnell}{
  pdf-ffbe9c7b3306-p184-b004,
  pdf-ffbe9c7b3306-p185-b003
}
\label{FABL.signPredicateSetTemplate}
\leanok
\uses{FABL.CSPTemplate, FABL.Sign, FABL.SignPredicate}
Given a finite set of sign predicates, form a constraint template over the sign domain
whose allowed predicates are that set. Each predicate retains its arity and acceptance
rule.
\end{definition}

\begin{definition}[Weighted constraint experiment]
\lean{FABL.WeightedConstraintExperiment}
\statementsource{Analysis-of-Boolean-Functions-by-Ryan-ODonnell}{
  pdf-ffbe9c7b3306-p185-b006,
  pdf-ffbe9c7b3306-p186-b001
}
\label{FABL.WeightedConstraintExperiment}
\leanok
\uses{FABL.CSPTemplate}
A weighted constraint experiment over a domain $D$, a constraint template, and a
variable set $V$ consists of a finite outcome space equipped with a probability law.
Each outcome selects an allowed predicate and an ordered tuple of variables of the
predicate's arity; repetitions in this tuple are allowed.
\end{definition}

\begin{definition}[Constraint payoff]
\lean{FABL.WeightedConstraintExperiment.payoff}
\statementsource{Analysis-of-Boolean-Functions-by-Ryan-ODonnell}{pdf-ffbe9c7b3306-p186-b001}
\label{FABL.WeightedConstraintExperiment.payoff}
\leanok
Given a weighted constraint experiment, an assignment of elements of $D$ to the
variables, and an outcome, the payoff is $1$ if the selected predicate accepts the
values assigned to the outcome's ordered query tuple, and $0$ otherwise.
\end{definition}

\begin{definition}[Expected constraint payoff]
\lean{FABL.WeightedConstraintExperiment.value}
\statementsource{Analysis-of-Boolean-Functions-by-Ryan-ODonnell}{pdf-ffbe9c7b3306-p186-b001}
\label{FABL.WeightedConstraintExperiment.value}
\leanok
\uses{FABL.WeightedConstraintExperiment.payoff, FABL.pmfExpectation}
For an assignment of elements of $D$ to the variables, the value of a weighted
constraint experiment is the expected payoff under its outcome law.
\end{definition}

\begin{theorem}[Payoff lies in the unit interval]
\lean{FABL.WeightedConstraintExperiment.payoff_mem_Icc}
\label{FABL.WeightedConstraintExperiment.payoff_mem_Icc}
\leanok
\uses{FABL.WeightedConstraintExperiment.payoff}
\difficulty{3}
In a weighted constraint experiment, fix an assignment of values to variables and an
outcome. The payoff, equal to $1$ if the outcome's predicate is satisfied by the
assigned values on its query tuple and $0$ otherwise, lies in $[0,1]$.
\end{theorem}

\begin{theorem}[Value mem icc]
\lean{FABL.WeightedConstraintExperiment.value_mem_Icc}
\statementsource{Analysis-of-Boolean-Functions-by-Ryan-ODonnell}{pdf-ffbe9c7b3306-p186-b001}
\label{FABL.WeightedConstraintExperiment.value_mem_Icc}
\leanok
\uses{FABL.WeightedConstraintExperiment.payoff_mem_Icc, FABL.WeightedConstraintExperiment.value, FABL.pmfExpectation_const, FABL.pmfExpectation_mono, FABL.pmfExpectation_nonneg}
\difficulty{3}
% dataset-pipeline: stub -- prose pending build_informal_statements.py
Statement pending regeneration.
\end{theorem}

\begin{definition}[Optimal assignment]
\lean{FABL.WeightedConstraintExperiment.optimalAssignment}
\statementsource{Analysis-of-Boolean-Functions-by-Ryan-ODonnell}{pdf-ffbe9c7b3306-p186-b001}
\label{FABL.WeightedConstraintExperiment.optimalAssignment}
\leanok
\uses{FABL.WeightedConstraintExperiment.value}
When $D$ and $V$ are finite and $D$ is nonempty, an optimal assignment is an assignment
of elements of $D$ to the variables that maximizes the experiment's value.
\end{definition}

\begin{definition}[Optimal value]
\lean{FABL.WeightedConstraintExperiment.optimum}
\statementsource{Analysis-of-Boolean-Functions-by-Ryan-ODonnell}{pdf-ffbe9c7b3306-p186-b001}
\label{FABL.WeightedConstraintExperiment.optimum}
\leanok
\uses{FABL.WeightedConstraintExperiment.optimalAssignment, FABL.WeightedConstraintExperiment.value}
When $D$ and $V$ are finite and $D$ is nonempty, the optimum is the value of an optimal
assignment.
\end{definition}

\begin{theorem}[Value bounded by the optimum]
\lean{FABL.WeightedConstraintExperiment.value_le_optimum}
\label{FABL.WeightedConstraintExperiment.value_le_optimum}
\leanok
\uses{FABL.WeightedConstraintExperiment.optimum, FABL.WeightedConstraintExperiment.value}
\difficulty{1}
Let $E$ be a weighted constraint experiment with finite variable set $V$ and finite,
nonempty domain $D$. For every assignment of elements of $D$ to the variables in $V$,
its value is at most the optimum of $E$.
\end{theorem}

\begin{definition}[Regular unique game]
\lean{FABL.RegularUniqueGame}
\statementsource{Analysis-of-Boolean-Functions-by-Ryan-ODonnell}{
  pdf-ffbe9c7b3306-p200-b007,
  pdf-ffbe9c7b3306-p201-b007
}
\label{FABL.RegularUniqueGame}
\leanok
A regular unique game on alphabet $[q]$ has a finite nonempty set of vertices. Each
vertex has a finite nonempty set of neighbor occurrences, each specifying a different
vertex and a permutation of $[q]$. There is a positive integer degree such that every
vertex has exactly that many neighbor occurrences.
\end{definition}

\begin{definition}[Finite game from a regular game]
\lean{FABL.RegularUniqueGame.toFiniteUniqueGame}
\label{FABL.RegularUniqueGame.toFiniteUniqueGame}
\leanok
\uses{FABL.FiniteUniqueGame}
Given a regular Unique Games instance, form a finite Unique Games instance with the same
variables and with one constraint occurrence for each variable together with one of its
neighbors. The source is that variable, the target is the neighbor, and the constraint
permutation is the one assigned to that neighboring pair.
\end{definition}

\begin{definition}[Expectation over constraint incidences]
\lean{FABL.RegularUniqueGame.edgeExpectation}
\label{FABL.RegularUniqueGame.edgeExpectation}
\leanok
Let $G$ be a regular Unique Games instance, and let $f$ be a real-valued function on its
variable-neighbor incidences. Its edge expectation is the uniform average of $f$ over
all such incidences.
\end{definition}

\begin{definition}[Iterated neighborhood expectation]
\lean{FABL.RegularUniqueGame.neighborhoodExpectation}
\label{FABL.RegularUniqueGame.neighborhoodExpectation}
\leanok
Let $G$ be a regular Unique Games instance, and let $f$ be a real-valued function on its
variable-neighbor incidences. Its neighborhood expectation is the average, first over
variables and then over the neighbors of each variable, of $f$ on the resulting
incidence.
\end{definition}

\begin{theorem}[Two expressions for edge expectation]
\lean{FABL.RegularUniqueGame.edgeExpectation_eq_neighborhoodExpectation}
\label{FABL.RegularUniqueGame.edgeExpectation_eq_neighborhoodExpectation}
\leanok
\uses{FABL.RegularUniqueGame.edgeExpectation, FABL.RegularUniqueGame.neighborhoodExpectation}
\difficulty{4}
Let $G$ be a regular unique game and let $f$ be a real-valued function on its directed
vertex-neighbor pairs. The expectation of $f$ over a uniformly chosen directed edge
equals the expectation obtained by first choosing a vertex uniformly and then choosing
one of its neighbors uniformly.
\end{theorem}

\begin{definition}[Long code table]
\lean{FABL.RegularUniqueGame.LongCodeTable}
\statementsource{Analysis-of-Boolean-Functions-by-Ryan-ODonnell}{
  pdf-ffbe9c7b3306-p191-b002,
  pdf-ffbe9c7b3306-p194-b003
}
\label{FABL.RegularUniqueGame.LongCodeTable}
\leanok
\uses{FABL.BooleanFunction}
For a regular unique game $G$ with alphabet size $q$, a long code table assigns to each
vertex of $G$ a Boolean function on the $q$-dimensional Boolean cube.
\end{definition}

\begin{definition}[Bounded long code table]
\lean{FABL.RegularUniqueGame.BoundedLongCodeTable}
\statementsource{Analysis-of-Boolean-Functions-by-Ryan-ODonnell}{pdf-ffbe9c7b3306-p191-b002}
\label{FABL.RegularUniqueGame.BoundedLongCodeTable}
\leanok
\uses{FABL.BoundedCubeFunction}
For a regular unique game $G$ with alphabet size $q$, a bounded long code table assigns
to each vertex of $G$ a bounded function on the $q$-dimensional Boolean cube.
\end{definition}

\begin{definition}[Bounded twisted long-code function]
\lean{FABL.RegularUniqueGame.boundedTwistedLongCode}
\label{FABL.RegularUniqueGame.boundedTwistedLongCode}
\leanok
\uses{FABL.RegularUniqueGame.BoundedLongCodeTable, FABL.permuteInput}
Let $G$ be a regular Unique Games instance with $q$ labels, and assign to each variable
a real-valued function on the sign cube $\\{-1,1\\}^q$ whose values lie in $[-1,1]$. For
a variable and one of its neighbors, the bounded twisted long-code function is the
neighbor's function evaluated after permuting its input by the permutation on that
variable-neighbor pair.
\end{definition}

\begin{definition}[Bounded averaged neighbor function]
\lean{FABL.RegularUniqueGame.boundedAveragedNeighborFunction}
\label{FABL.RegularUniqueGame.boundedAveragedNeighborFunction}
\leanok
\uses{FABL.BoundedCubeFunction, FABL.RegularUniqueGame.BoundedLongCodeTable, FABL.RegularUniqueGame.boundedTwistedLongCode, FABL.permuteInput}
Let $G$ be a regular unique game with alphabet size $q$, and let each vertex $u$ be
assigned a function $f_u:\
\{−1,1\}^q\to\mathbb R$ taking values in $[-1,1]$. For a vertex $v$, the bounded
averaged neighbor function at query $x$ is the uniform average over neighbors $w$ of $v$
of $f_w$ evaluated at the query obtained by permuting $x$ according to the constraint on
$(v,w)$.
\end{definition}

\begin{definition}[Bounded long code value]
\lean{FABL.RegularUniqueGame.boundedLongCodeValue}
\label{FABL.RegularUniqueGame.boundedLongCodeValue}
\leanok
\uses{FABL.RegularUniqueGame.BoundedLongCodeTable, FABL.RegularUniqueGame.boundedAveragedNeighborFunction, FABL.SignTester, FABL.SignTester.acceptanceProbability}
Let $G$ be a regular unique game with alphabet size $q$, let a finite randomized sign
tester of query dimension $q$ be given, and assign to each vertex of $G$ a function from
the $q$-dimensional Boolean cube to the signs. The bounded long code value is the
uniform average over vertices of the tester's acceptance probability when its queries
are answered using that vertex's bounded averaged neighbor function.
\end{definition}

\begin{definition}[Twisted long code]
\lean{FABL.RegularUniqueGame.twistedLongCode}
\label{FABL.RegularUniqueGame.twistedLongCode}
\leanok
\uses{FABL.BooleanFunction.toReal, FABL.RegularUniqueGame.LongCodeTable, FABL.permuteInput}
Let $G$ be a regular unique game with alphabet size $q$, and assign to each vertex a
function from the $q$-dimensional Boolean cube to the signs. For a vertex $v$, a
neighbor $w$ of $v$, and a query $x$, the twisted long code is the real value of the
sign returned by $w$'s function at the query obtained by permuting $x$ according to the
constraint on $(v,w)$.
\end{definition}

\begin{definition}[Averaged neighbor function]
\lean{FABL.RegularUniqueGame.averagedNeighborFunction}
\label{FABL.RegularUniqueGame.averagedNeighborFunction}
\leanok
\uses{FABL.BooleanFunction.toReal, FABL.BoundedCubeFunction, FABL.RegularUniqueGame.LongCodeTable, FABL.RegularUniqueGame.twistedLongCode, FABL.permuteInput, FABL.signValue_eq_neg_one_or_one}
Let $G$ be a regular unique game with alphabet size $q$, and assign to each vertex a
function from the $q$-dimensional Boolean cube to the signs. For a vertex $v$, its
averaged neighbor function sends each query $x$ to the uniform average, over neighbors
$w$ of $v$, of the real value of the sign returned by $w$'s function at the query
permuted according to the constraint on $(v,w)$. This function takes values in $[-1,1]$.
\end{definition}

\begin{definition}[Dictator long code]
\lean{FABL.RegularUniqueGame.dictatorLongCode}
\statementsource{Analysis-of-Boolean-Functions-by-Ryan-ODonnell}{
  pdf-ffbe9c7b3306-p175-b003,
  pdf-ffbe9c7b3306-p191-b002
}
\label{FABL.RegularUniqueGame.dictatorLongCode}
\leanok
\uses{FABL.RegularUniqueGame.LongCodeTable, FABL.dictator}
Let $G$ be a regular unique game with alphabet size $q$, and let $
ho$ assign a label in
$\\{0,\\ldots,q-1\\}$ to each vertex. The dictator long code table assigns to each
vertex the dictator function indexed by its label under $\\rho$.
\end{definition}

\begin{definition}[Long code value]
\lean{FABL.RegularUniqueGame.longCodeValue}
\label{FABL.RegularUniqueGame.longCodeValue}
\leanok
\uses{FABL.RegularUniqueGame.LongCodeTable, FABL.RegularUniqueGame.averagedNeighborFunction, FABL.SignTester, FABL.SignTester.acceptanceProbability}
Let $G$ be a regular unique game with alphabet size $q$, let a finite randomized sign
tester of query dimension $q$ be given, and assign to each vertex of $G$ a function from
the $q$-dimensional Boolean cube to the signs. The long code value is the uniform
average over vertices of the tester's acceptance probability when its queries are
answered using that vertex's averaged neighbor function.
\end{definition}

\begin{theorem}[Multilinear extension at sign inputs]
\lean{FABL.RegularUniqueGame.multilinearExtension_signValues}
\label{FABL.RegularUniqueGame.multilinearExtension_signValues}
\leanok
\uses{FABL.Sign, FABL.SignPredicate, FABL.SignPredicate.indicator, FABL.SignPredicate.multilinearExtension, FABL.fourier_expansion, FABL.monomial, FABL.signValue}
\difficulty{3}
For any sign predicate $P$ and any tuple $y$ of signs indexed by its input positions,
the multilinear extension of $P$'s acceptance indicator, evaluated at the real values of
the signs in $y$, equals the acceptance indicator of $P$ at $y$.
\end{theorem}

\begin{theorem}[Multilinear Extension and Product Expectations]
\lean{FABL.RegularUniqueGame.pmfExpectation_multilinearExtension_productLaw}
\label{FABL.RegularUniqueGame.pmfExpectation_multilinearExtension_productLaw}
\leanok
\uses{FABL.SignPredicate, FABL.SignPredicate.multilinearExtension, FABL.independentProductPMF, FABL.pmfExpectation, FABL.pmfExpectation_const_mul, FABL.pmfExpectation_const_one, FABL.pmfExpectation_independentProductPMF_prod, FABL.pmfExpectation_sum}
\difficulty{5}
Let $P$ be a predicate on sign vectors with finite arity $m$. For each position
$i\in\text{Fin}(m)$, let $A_i$ be a finite sample space with probability law $\mu_i$,
and let $g_i:A_i\to\mathbb{R}$. The expectation, under the independent product law, of
the multilinear extension of $P$ evaluated at the values $g_i$ of the sampled
coordinates equals the multilinear extension evaluated at the coordinate means
$\mathbb{E}_{\mu_i}[g_i]$.
\end{theorem}

\begin{theorem}[Marginal Expectation of a Product Law]
\lean{FABL.RegularUniqueGame.pmfExpectation_independentProductPMF_coordinate}
\label{FABL.RegularUniqueGame.pmfExpectation_independentProductPMF_coordinate}
\leanok
\uses{FABL.independentProductPMF, FABL.pmfExpectation, FABL.pmfExpectation_const_one, FABL.pmfExpectation_independentProductPMF_prod}
\difficulty{4}
Let $I$ be a finite index set, let each index $i$ have a finite sample space $A_i$ and a
probability law $\mu_i$, and fix $j\in I$ and a function $f:A_j\to\mathbb{R}$. Under the
product law with independent coordinates distributed according to the respective
$\mu_i$, the expectation of $f$ evaluated at coordinate $j$ equals its expectation under
$\mu_j$.
\end{theorem}

\begin{theorem}[Two-Coordinate Expectation Under a Product Law]
\lean{FABL.RegularUniqueGame.pmfExpectation_independentProductPMF_two_coordinates}
\label{FABL.RegularUniqueGame.pmfExpectation_independentProductPMF_two_coordinates}
\leanok
\uses{FABL.independentProductPMF, FABL.pmfExpectation, FABL.pmfExpectation_const_one, FABL.pmfExpectation_independentProductPMF_prod, FABL.pmfExpectation_mul_const, FABL.pmfExpectation_sum}
\difficulty{6}
Let $I$ and $A$ be finite types with decidable equality on $I$, and assign a probability
law $\mu_i$ on $A$ to each $i\in I$. For distinct coordinates $j,k\in I$ and any
function $F:A\times A\to\mathbb{R}$, its expectation on the pair of coordinates $(j,k)$
under the independent product law equals the iterated expectation of $F$ under $\mu_j$
and $\mu_k$.
\end{theorem}

\begin{theorem}[Randomized value of boolean]
\lean{FABL.RegularUniqueGame.randomizedValue_ofBoolean}
\statementsource{Analysis-of-Boolean-Functions-by-Ryan-ODonnell}{pdf-ffbe9c7b3306-p198-b006}
\label{FABL.RegularUniqueGame.randomizedValue_ofBoolean}
\leanok
\uses{FABL.BooleanFunction, FABL.BoundedCubeFunction.ofBoolean, FABL.RegularUniqueGame.multilinearExtension_signValues, FABL.SignPredicate.indicator, FABL.SignPredicate.randomizedValue_eq_multilinearExtension, FABL.SignQueryCheck, FABL.SignQueryCheck.randomizedValue}
\difficulty{3}
% dataset-pipeline: stub -- prose pending build_informal_statements.py
Statement pending regeneration.
\end{theorem}

\begin{definition}[Neighbor response]
\lean{FABL.RegularUniqueGame.neighborResponse}
\label{FABL.RegularUniqueGame.neighborResponse}
\leanok
\uses{FABL.RegularUniqueGame.LongCodeTable, FABL.Sign, FABL.permuteInput}
Let $G$ be a regular unique game with alphabet size $q$, and assign to each vertex a
function from the $q$-dimensional Boolean cube to the signs. Given a vertex $v$, a
neighbor $w$ of $v$, and a query $x$, the neighbor response is the sign returned by
$w$'s function on the query obtained by permuting $x$ according to the constraint on
$(v,w)$.
\end{definition}

\begin{theorem}[Expected Neighbor Response]
\lean{FABL.RegularUniqueGame.pmfExpectation_neighborResponse}
\label{FABL.RegularUniqueGame.pmfExpectation_neighborResponse}
\leanok
\uses{FABL.RegularUniqueGame.LongCodeTable, FABL.RegularUniqueGame.averagedNeighborFunction, FABL.RegularUniqueGame.neighborResponse, FABL.pmfExpectation, FABL.pmfExpectation_uniformPMF_eq_expect, FABL.signValue, FABL.uniformPMF}
\difficulty{3}
Let $G$ be a regular unique game with a finite, nonempty neighbor set at each vertex,
and let a long-code table assign a Boolean function to each vertex. For every vertex $v$
and query $x\in\{\pm1\}^q$, the uniform expectation over neighbors $w$ of the real sign
value returned by the table at the permuted query equals the averaged neighbor function
at $v$ evaluated at $x$.
\end{theorem}

\begin{definition}[Edge failure indicator]
\lean{FABL.RegularUniqueGame.edgeFailureIndicator}
\label{FABL.RegularUniqueGame.edgeFailureIndicator}
\leanok
Let $G$ be a regular unique game with alphabet size $q$, and assign each vertex a label
from the $q$-element alphabet. For a vertex $v$ and a neighbor $w$ of $v$, the edge
failure indicator is $0$ if the constraint permutation on $(v,w)$ sends $w$'s label to
$v$'s label, and is $1$ otherwise.
\end{definition}

\begin{definition}[Local edge error]
\lean{FABL.RegularUniqueGame.localEdgeError}
\label{FABL.RegularUniqueGame.localEdgeError}
\leanok
\uses{FABL.RegularUniqueGame.edgeFailureIndicator}
Let $G$ be a regular unique game with alphabet size $q$, and assign each vertex a label
from the $q$-element alphabet. At a vertex $v$, the local edge error is the uniform
average over neighbors $w$ of $v$ of the edge failure indicator for $(v,w)$.
\end{definition}

\begin{definition}[Average edge error]
\lean{FABL.RegularUniqueGame.edgeError}
\label{FABL.RegularUniqueGame.edgeError}
\leanok
\uses{FABL.RegularUniqueGame.localEdgeError}
Let $G$ be a regular unique game with alphabet size $q$, and let a labeling assign one
alphabet symbol to each vertex. Its edge error is the average, over a uniformly chosen
vertex of $G$, of the local edge error of the labeling at that vertex.
\end{definition}

\begin{theorem}[Nonnegativity of local edge error]
\lean{FABL.RegularUniqueGame.localEdgeError_nonneg}
\label{FABL.RegularUniqueGame.localEdgeError_nonneg}
\leanok
\uses{FABL.RegularUniqueGame.edgeFailureIndicator, FABL.RegularUniqueGame.localEdgeError}
\difficulty{4}
Let $G$ be a regular unique game, let $\ell$ be a labeling of its vertices by alphabet
symbols, and fix a vertex $v$. The average over neighbors of $v$ of the indicator that
the edge constraint fails is nonnegative.
\end{theorem}

\begin{theorem}[Nonnegativity of edge error]
\lean{FABL.RegularUniqueGame.edgeError_nonneg}
\label{FABL.RegularUniqueGame.edgeError_nonneg}
\leanok
\uses{FABL.RegularUniqueGame.edgeError, FABL.RegularUniqueGame.localEdgeError_nonneg}
\difficulty{1}
Let $G$ be a regular unique game with alphabet size $q$, and let a labeling assign one
of the $q$ labels to each vertex. The expected fraction of neighbors whose constraint is
violated, averaged over vertices, is nonnegative.
\end{theorem}

\begin{theorem}[Edge error and constraint value]
\lean{FABL.RegularUniqueGame.edgeError_eq_one_sub_value}
\label{FABL.RegularUniqueGame.edgeError_eq_one_sub_value}
\leanok
\uses{FABL.CSPConstraint.Satisfied, FABL.CSPConstraint.unsatisfiedIndicator, FABL.CSPInstance.unsatisfiedCost, FABL.CSPInstance.unsatisfiedCost_eq_sum, FABL.CSPInstance.value, FABL.CSPInstance.value_eq_one_sub_unsatisfiedCost_div, FABL.FiniteUniqueGame.constraint, FABL.FiniteUniqueGame.constraint_satisfied_iff, FABL.FiniteUniqueGame.toCSPInstance, FABL.RegularUniqueGame.edgeError, FABL.RegularUniqueGame.edgeExpectation, FABL.RegularUniqueGame.edgeExpectation_eq_neighborhoodExpectation, FABL.RegularUniqueGame.edgeFailureIndicator, FABL.RegularUniqueGame.neighborhoodExpectation, FABL.RegularUniqueGame.toFiniteUniqueGame}
\difficulty{6}
Let $G$ be a regular unique game and let $\ell$ assign an alphabet label to each vertex.
The average edge error of $\ell$ equals one minus the value of the associated finite
unique-game constraint instance under $\ell$.
\end{theorem}

\begin{definition}[Optimal labeling]
\lean{FABL.RegularUniqueGame.optimalLabel}
\statementsource{Analysis-of-Boolean-Functions-by-Ryan-ODonnell}{pdf-ffbe9c7b3306-p186-b001}
\label{FABL.RegularUniqueGame.optimalLabel}
\leanok
\uses{FABL.CSPInstance.optimalAssignment, FABL.FiniteUniqueGame.toCSPInstance, FABL.RegularUniqueGame.toFiniteUniqueGame}
Let $G$ be a regular unique game with alphabet size $q$, and suppose the label set
$\\{0,\\ldots,q-1\\}$ is nonempty. An optimal labeling of $G$ assigns a label from this
set to every vertex, according to an assignment maximizing the associated constraint
satisfaction objective.
\end{definition}

\begin{theorem}[Error of an optimal labeling]
\lean{FABL.RegularUniqueGame.edgeError_optimalLabel_le}
\label{FABL.RegularUniqueGame.edgeError_optimalLabel_le}
\leanok
\uses{FABL.CSPInstance.optimum, FABL.FiniteUniqueGame.toCSPInstance, FABL.RegularUniqueGame.edgeError, FABL.RegularUniqueGame.edgeError_eq_one_sub_value, FABL.RegularUniqueGame.edgeError_nonneg, FABL.RegularUniqueGame.optimalLabel, FABL.RegularUniqueGame.toFiniteUniqueGame}
\difficulty{3}
Let $G$ be a regular unique game with $q>0$ labels, and let $-\delta$ be a lower bound
for the maximum fraction of constraints satisfiable by a labeling of its vertices. Then
the expected fraction of violated constraints under an optimal labeling is at most
$\delta$.
\end{theorem}

\begin{theorem}[Dictator response under a satisfied constraint]
\lean{FABL.RegularUniqueGame.neighborResponse_dictator_of_satisfied}
\label{FABL.RegularUniqueGame.neighborResponse_dictator_of_satisfied}
\leanok
\uses{FABL.RegularUniqueGame.dictatorLongCode, FABL.RegularUniqueGame.neighborResponse, FABL.dictator, FABL.permuteInput}
\difficulty{1}
Let $G$ be a regular unique game, let $
ho$ assign an alphabet coordinate to each
vertex, and choose a vertex $v$, a neighbour $w$ of $v$, and a query $x$ in the Boolean
cube. If the permutation on this constraint sends the coordinate assigned to $w$ to the
coordinate assigned to $v$, then the response of the dictator code induced by $
ho$ at
$v$ on query $x$ and neighbour $w$ equals the dictator response at $v$ on $x$.
\end{theorem}

\begin{definition}[Operational check value]
\lean{FABL.RegularUniqueGame.operationalCheckValue}
\label{FABL.RegularUniqueGame.operationalCheckValue}
\leanok
\uses{FABL.RegularUniqueGame.LongCodeTable, FABL.RegularUniqueGame.neighborResponse, FABL.SignPredicate.indicator, FABL.SignQueryCheck, FABL.independentProductPMF, FABL.pmfExpectation, FABL.uniformPMF}
Let $G$ be a regular unique game with alphabet size $q$, let a table assign a Boolean
function to each vertex, and fix a vertex and a sign query check. Independently choose a
uniformly random neighbor of that vertex for each input position of the check; the
operational check value is the probability that the check's predicate accepts the
resulting neighbor responses from the table.
\end{definition}

\begin{theorem}[Dictator check versus operational check]
\lean{FABL.RegularUniqueGame.dictatorIndicator_le_operationalCheckValue_add_error}
\label{FABL.RegularUniqueGame.dictatorIndicator_le_operationalCheckValue_add_error}
\leanok
\uses{FABL.RegularUniqueGame.dictatorLongCode, FABL.RegularUniqueGame.edgeFailureIndicator, FABL.RegularUniqueGame.localEdgeError, FABL.RegularUniqueGame.neighborResponse, FABL.RegularUniqueGame.neighborResponse_dictator_of_satisfied, FABL.RegularUniqueGame.operationalCheckValue, FABL.RegularUniqueGame.pmfExpectation_independentProductPMF_coordinate, FABL.Sign, FABL.SignPredicate.indicator, FABL.SignQueryCheck, FABL.dictator, FABL.independentProductPMF, FABL.pmfExpectation, FABL.pmfExpectation_add, FABL.pmfExpectation_const, FABL.pmfExpectation_mono, FABL.pmfExpectation_sum, FABL.pmfExpectation_uniformPMF_eq_expect, FABL.uniformPMF}
\difficulty{6}
Let $G$ be a regular unique game with alphabet size $q$, let $\ell$ label each vertex by
an alphabet symbol, and fix a vertex and a sign-valued query check of arity $k$. The
check's indicator on the dictator queries determined by the label at that vertex is at
most the operational check value using the dictator long-code table, plus $k$ times the
vertex's local edge error.
\end{theorem}

\begin{theorem}[Operational and Randomized Check Values]
\lean{FABL.RegularUniqueGame.operationalCheckValue_eq_randomizedValue}
\label{FABL.RegularUniqueGame.operationalCheckValue_eq_randomizedValue}
\leanok
\uses{FABL.RegularUniqueGame.LongCodeTable, FABL.RegularUniqueGame.averagedNeighborFunction, FABL.RegularUniqueGame.multilinearExtension_signValues, FABL.RegularUniqueGame.neighborResponse, FABL.RegularUniqueGame.operationalCheckValue, FABL.RegularUniqueGame.pmfExpectation_multilinearExtension_productLaw, FABL.RegularUniqueGame.pmfExpectation_neighborResponse, FABL.Sign, FABL.SignPredicate.indicator, FABL.SignPredicate.multilinearExtension, FABL.SignPredicate.randomizedValue, FABL.SignPredicate.randomizedValue_eq_multilinearExtension, FABL.SignQueryCheck, FABL.SignQueryCheck.randomizedValue, FABL.independentProductPMF, FABL.pmfExpectation, FABL.signValue, FABL.uniformPMF}
\difficulty{4}
Let $G$ be a regular unique game with label set $\text{Fin}(q)$, let $T$ be a long-code
table, let $v$ be a vertex, and let $C$ be a sign-query check. The check's expected
acceptance when each query position independently chooses a uniformly random neighbor of
$v$ and returns that neighbor's permuted table response equals the check's randomized
value on the function obtained by averaging those permuted responses over the neighbors
of $v$.
\end{theorem}

\begin{definition}[Long code experiment sample]
\lean{FABL.RegularUniqueGame.LongCodeExperimentSample}
\label{FABL.RegularUniqueGame.LongCodeExperimentSample}
\leanok
\uses{FABL.SignTester}
Let $G$ be a regular unique game with alphabet size $q$, and let a sign tester have a
finite seed set. A sample consists of a vertex of $G$, a tester seed, and, for each
input position of the check selected by that seed, a neighbor of the chosen vertex.
\end{definition}

\begin{definition}[Long code experiment law]
\lean{FABL.RegularUniqueGame.longCodeExperimentLaw}
\label{FABL.RegularUniqueGame.longCodeExperimentLaw}
\leanok
\uses{FABL.RegularUniqueGame.LongCodeExperimentSample, FABL.SignTester, FABL.independentProductPMF, FABL.uniformPMF}
Let $G$ be a regular unique game with alphabet size $q$, and let a sign tester have a
finite seed set equipped with a probability law. The experiment law chooses a vertex
uniformly, chooses a seed according to that law, and then independently chooses a
uniformly random neighbor of the vertex for every input position of the check selected
by the seed; its outcome is the resulting long code experiment sample.
\end{definition}

\begin{definition}[Long code experiment responses]
\lean{FABL.RegularUniqueGame.longCodeExperimentResponses}
\label{FABL.RegularUniqueGame.longCodeExperimentResponses}
\leanok
\uses{FABL.RegularUniqueGame.LongCodeExperimentSample, FABL.RegularUniqueGame.LongCodeTable, FABL.RegularUniqueGame.neighborResponse, FABL.Sign, FABL.SignTester}
Let $G$ be a regular unique game with alphabet size $q$, let a sign tester have a finite
seed set, and let a table assign a Boolean function to each vertex. For a sample
consisting of a vertex, a seed, and the sampled neighbors, the response at each input
position is the table's neighbor response at the sampled vertex and neighbor, for the
query selected at that position by the seed's check.
\end{definition}

\begin{definition}[Long code experiment variable]
\lean{FABL.RegularUniqueGame.longCodeExperimentVariable}
\label{FABL.RegularUniqueGame.longCodeExperimentVariable}
\leanok
\uses{FABL.RegularUniqueGame.LongCodeExperimentSample, FABL.SignTester, FABL.permuteInput}
Let $G$ be a regular unique game with alphabet size $q$, let a sign tester have a finite
seed set, and fix a long code experiment sample and one input position of its selected
check. The associated variable is the pair consisting of the vertex reached from the
sampled vertex along the sampled neighbor and the query at that position after applying
the game's permutation for that edge.
\end{definition}

\begin{definition}[Long code weighted experiment]
\lean{FABL.RegularUniqueGame.longCodeWeightedExperiment}
\label{FABL.RegularUniqueGame.longCodeWeightedExperiment}
\leanok
\uses{FABL.RegularUniqueGame.LongCodeExperimentSample, FABL.RegularUniqueGame.longCodeExperimentLaw, FABL.RegularUniqueGame.longCodeExperimentVariable, FABL.Sign, FABL.SignTester, FABL.SignTester.cspTemplate, FABL.WeightedConstraintExperiment}
Let $G$ be a regular unique game with alphabet size $q$, and let a sign tester have a
finite seed set and a probability law on it. The associated weighted constraint
experiment has as outcomes the long code experiment samples, distributed according to
the experiment law; the predicate selected by an outcome is its tester seed, and its
ordered scope is the tuple of long code experiment variables at the check's input
positions.
\end{definition}

\begin{definition}[Allowed long code weighted experiment]
\lean{FABL.RegularUniqueGame.longCodeAllowedWeightedExperiment}
\label{FABL.RegularUniqueGame.longCodeAllowedWeightedExperiment}
\leanok
\uses{FABL.RegularUniqueGame.LongCodeExperimentSample, FABL.RegularUniqueGame.longCodeExperimentLaw, FABL.RegularUniqueGame.longCodeExperimentVariable, FABL.Sign, FABL.SignPredicate, FABL.SignTester, FABL.SignTester.UsesPredicates, FABL.WeightedConstraintExperiment, FABL.signPredicateSetTemplate}
Let $G$ be a regular unique game with alphabet size $q$, let a sign tester have a finite
seed set and a probability law, and let $[?]$ be a finite set of sign predicates.
Suppose every check selected by the tester uses a predicate in this set. The associated
weighted constraint experiment has as outcomes the long code experiment samples,
distributed according to the experiment law; each outcome selects its check's predicate,
and its ordered scope is the tuple of long code experiment variables at the check's
input positions.
\end{definition}

\begin{definition}[Long-code assignment]
\lean{FABL.RegularUniqueGame.longCodeAssignment}
\label{FABL.RegularUniqueGame.longCodeAssignment}
\leanok
\uses{FABL.RegularUniqueGame.LongCodeTable, FABL.Sign}
Let $G$ be a regular unique game with alphabet size $q$, and let a table assign to each
vertex a Boolean function on $\\{ -1,1 \\}^q$ with values in the sign set $\\{ -1,1
\\}$. The long-code assignment maps each pair consisting of a vertex and an input in
$\\{ -1,1 \\}^q$ to the value of that vertex's function at the input.
\end{definition}

\begin{definition}[Long-code experiment value]
\lean{FABL.RegularUniqueGame.longCodeExperimentValue}
\label{FABL.RegularUniqueGame.longCodeExperimentValue}
\leanok
\uses{FABL.RegularUniqueGame.LongCodeTable, FABL.RegularUniqueGame.longCodeExperimentLaw, FABL.RegularUniqueGame.longCodeExperimentResponses, FABL.SignPredicate.indicator, FABL.SignTester, FABL.pmfExpectation}
Let $G$ be a regular unique game with alphabet size $q$, let a finite-seed tester
specify a sign-valued predicate check for each seed, and let a long-code table be given.
The experiment value is the expected value of the check's acceptance indicator, when the
sample is drawn from the law of the long-code experiment determined by $G$, the tester,
and the table.
\end{definition}

\begin{theorem}[Weighted experiment value identity]
\lean{FABL.RegularUniqueGame.longCodeWeightedExperiment_value}
\label{FABL.RegularUniqueGame.longCodeWeightedExperiment_value}
\leanok
\uses{FABL.RegularUniqueGame.LongCodeTable, FABL.RegularUniqueGame.longCodeAssignment, FABL.RegularUniqueGame.longCodeExperimentValue, FABL.RegularUniqueGame.longCodeWeightedExperiment, FABL.SignTester, FABL.WeightedConstraintExperiment.value}
\difficulty{1}
Let $G$ be a regular unique game, let $T$ be a finite randomised tester, and let $F$
assign a sign-valued Boolean function to each vertex of $G$. The value of the weighted
constraint experiment associated with $G$ and $T$, under the assignment that gives each
pair consisting of a vertex and a Boolean-cube input the value of $F$ at that vertex and
input, equals the experiment value of $F$.
\end{theorem}

\begin{theorem}[Iterated expression for experiment value]
\lean{FABL.RegularUniqueGame.longCodeExperimentValue_eq_iteratedOperational}
\label{FABL.RegularUniqueGame.longCodeExperimentValue_eq_iteratedOperational}
\leanok
\uses{FABL.RegularUniqueGame.LongCodeTable, FABL.RegularUniqueGame.longCodeExperimentLaw, FABL.RegularUniqueGame.longCodeExperimentValue, FABL.RegularUniqueGame.operationalCheckValue, FABL.SignTester, FABL.pmfExpectation, FABL.pmfExpectation_bind, FABL.pmfExpectation_map, FABL.pmfExpectation_uniformPMF_eq_expect}
\difficulty{3}
Let $G$ be a regular unique game and let $T$ be a sign tester with a finite seed space.
For every long-code table, the value of the long-code experiment equals the average over
vertices, chosen uniformly, of the average over tester seeds according to the seed law
of the operational check value at that vertex.
\end{theorem}

\begin{theorem}[Experiment value as averaged acceptance]
\lean{FABL.RegularUniqueGame.longCodeExperimentValue_eq_longCodeValue}
\label{FABL.RegularUniqueGame.longCodeExperimentValue_eq_longCodeValue}
\leanok
\uses{FABL.RegularUniqueGame.LongCodeTable, FABL.RegularUniqueGame.longCodeExperimentValue, FABL.RegularUniqueGame.longCodeExperimentValue_eq_iteratedOperational, FABL.RegularUniqueGame.longCodeValue, FABL.RegularUniqueGame.operationalCheckValue_eq_randomizedValue, FABL.SignTester, FABL.SignTester.acceptanceProbability, FABL.pmfExpectation}
\difficulty{3}
Let $G$ be a regular unique game, let $T$ be a finite randomised tester over its
alphabet, and let $F$ assign a sign-valued Boolean function to each vertex of $G$. The
expected acceptance value of the experiment that chooses a uniform vertex, a tester
seed, and independent uniform neighbours for the seed's query positions, then evaluates
the predicate on the corresponding permuted responses from $F$, equals the uniform
average over vertices of the tester's acceptance probability on the function obtained by
averaging the permuted responses over neighbours.
\end{theorem}

\begin{theorem}[Long-code value identity]
\lean{FABL.RegularUniqueGame.longCode_value_identity}
\label{FABL.RegularUniqueGame.longCode_value_identity}
\leanok
\uses{FABL.RegularUniqueGame.LongCodeTable, FABL.RegularUniqueGame.averagedNeighborFunction, FABL.RegularUniqueGame.longCodeExperimentValue, FABL.RegularUniqueGame.longCodeExperimentValue_eq_longCodeValue, FABL.SignTester, FABL.SignTester.acceptanceProbability}
\difficulty{1}
Let $G$ be a regular unique game, let $T$ be a finite randomised tester over its
alphabet, and let $F$ assign a sign-valued Boolean function to each vertex of $G$. The
experiment value of $F$ equals the uniform average over vertices of the tester's
acceptance probability on the function obtained by averaging the permuted responses from
neighbouring vertices.
\end{theorem}

\begin{theorem}[Dictator value from low edge error]
\lean{FABL.RegularUniqueGame.longCodeExperimentValue_dictator_ge_of_edgeError}
\label{FABL.RegularUniqueGame.longCodeExperimentValue_dictator_ge_of_edgeError}
\leanok
\uses{FABL.BoundedCubeFunction.ofBoolean, FABL.RegularUniqueGame.dictatorIndicator_le_operationalCheckValue_add_error, FABL.RegularUniqueGame.dictatorLongCode, FABL.RegularUniqueGame.edgeError, FABL.RegularUniqueGame.localEdgeError, FABL.RegularUniqueGame.localEdgeError_nonneg, FABL.RegularUniqueGame.longCodeExperimentValue, FABL.RegularUniqueGame.longCodeExperimentValue_eq_iteratedOperational, FABL.RegularUniqueGame.operationalCheckValue, FABL.RegularUniqueGame.randomizedValue_ofBoolean, FABL.SignPredicate.indicator, FABL.SignTester, FABL.SignTester.acceptanceProbability, FABL.dictator, FABL.pmfExpectation, FABL.pmfExpectation_add, FABL.pmfExpectation_const, FABL.pmfExpectation_mono}
\difficulty{5}
Let $G$ be a regular unique game with alphabet size $q$, let $T$ be a sign tester, and
let $r$ be a natural number such that every predicate used by $T$ has arity at most $r$.
Suppose there are real numbers $\beta$ and $\delta$ such that, for every alphabet
coordinate, $T$ accepts the corresponding dictator function with probability at least
$\beta$. If a vertex labeling has edge error at most $\delta$, then the long-code
experiment using the dictator functions specified by that labeling has value at least
$\beta-r\delta$.
\end{theorem}

\begin{theorem}[Optimal dictator experiment guarantee]
\lean{FABL.RegularUniqueGame.longCodeExperimentValue_optimalDictator_ge}
\label{FABL.RegularUniqueGame.longCodeExperimentValue_optimalDictator_ge}
\leanok
\uses{FABL.BoundedCubeFunction.ofBoolean, FABL.CSPInstance.optimum, FABL.FiniteUniqueGame.toCSPInstance, FABL.RegularUniqueGame.dictatorLongCode, FABL.RegularUniqueGame.edgeError_optimalLabel_le, FABL.RegularUniqueGame.longCodeExperimentValue, FABL.RegularUniqueGame.longCodeExperimentValue_dictator_ge_of_edgeError, FABL.RegularUniqueGame.optimalLabel, FABL.RegularUniqueGame.toFiniteUniqueGame, FABL.SignTester, FABL.SignTester.acceptanceProbability, FABL.dictator}
\difficulty{1}
Let $G$ be a regular unique game with a nonempty alphabet, let $T$ be a finite
randomised tester, and let $r$ be a natural number such that every test selected by a
seed has arity at most $r$. Suppose that for every alphabet coordinate, the tester's
acceptance probability on the corresponding dictator function is at least $eta$, and
that the optimum value of the unique-game constraint instance associated with $G$ is at
least $1-\delta$. Then the experiment value of the dictator code induced by an optimal
labeling of $G$ is at least $\beta-r\delta$.
\end{theorem}

\begin{theorem}[Stable Influence of a Uniform Function Average]
\lean{FABL.RegularUniqueGame.stableInfluence_uniformExpectation_le}
\label{FABL.RegularUniqueGame.stableInfluence_uniformExpectation_le}
\leanok
\uses{FABL.fourierSquareFunctional, FABL.fourierSquareFunctional_uniformExpectation_le, FABL.stableInfluence}
\difficulty{3}
Let $\iota$ be a finite nonempty index set, let $\rho\in[0,1]$, and let
$f_j:\{\pm1\}^q\to\mathbb R$ be a family of functions indexed by $j\in\iota$. For every
coordinate $i$, the stable influence at parameter $\rho$ of their pointwise uniform
average is at most the uniform average of their stable influences at that coordinate.
\end{theorem}

\begin{theorem}[Stable Influence Under Coordinate Permutation]
\lean{FABL.RegularUniqueGame.stableInfluence_comp_permuteInput}
\label{FABL.RegularUniqueGame.stableInfluence_comp_permuteInput}
\leanok
\uses{FABL.fourierCoeff_comp_permuteInput, FABL.permuteFinset, FABL.permuteInput, FABL.stableInfluence}
\difficulty{3}
Let $q$ be a natural number, let $\rho\in\mathbb R$, let $\pi$ be a permutation of the
$q$ coordinates, and let $f:\{\pm1\}^q\to\mathbb R$. For every coordinate $i$, the
stable influence of the function obtained by permuting the input coordinates according
to $\pi$ at $i$ equals the stable influence of $f$ at $\pi^{-1}(i)$.
\end{theorem}

\begin{definition}[Neighbor notable coordinates]
\lean{FABL.RegularUniqueGame.neighborNotableCoordinates}
\statementsource{Analysis-of-Boolean-Functions-by-Ryan-ODonnell}{
  pdf-ffbe9c7b3306-p145-b008,
  pdf-ffbe9c7b3306-p191-b002,
  pdf-ffbe9c7b3306-p195-b001
}
\label{FABL.RegularUniqueGame.neighborNotableCoordinates}
\leanok
\uses{FABL.RegularUniqueGame.LongCodeTable, FABL.RegularUniqueGame.averagedNeighborFunction, FABL.stableInfluence}
Let $G$ be a regular unique game with alphabet size $q$, let $T$ be a long code table,
let $\\varepsilon\\in\\mathbb R$, and let $v$ be a vertex of $G$. The neighbor notable
coordinates at $v$ are the coordinates $i\\in\\{0,\\ldots,q-1\\}$ for which
$\\varepsilon$ is strictly less than the stable influence, at noise parameter
$1-\\varepsilon$, of coordinate $i$ in the averaged neighbor function of $T$ at $v$.
\end{definition}

\begin{definition}[Notable coordinates]
\lean{FABL.RegularUniqueGame.notableCoordinates}
\statementsource{Analysis-of-Boolean-Functions-by-Ryan-ODonnell}{
  pdf-ffbe9c7b3306-p057-b001,
  pdf-ffbe9c7b3306-p057-b004,
  pdf-ffbe9c7b3306-p145-b008,
  pdf-ffbe9c7b3306-p191-b002
}
\label{FABL.RegularUniqueGame.notableCoordinates}
\leanok
\uses{FABL.BooleanFunction.toReal, FABL.RegularUniqueGame.LongCodeTable, FABL.stableInfluence}
Let $G$ be a regular unique game with alphabet size $q$, let $T$ be a long code table,
let $\\varepsilon\\in\\mathbb R$, and let $v$ be a vertex of $G$. The notable
coordinates at $v$ are the coordinates $i\\in\\{0,\\ldots,q-1\\}$ for which the stable
influence, at noise parameter $1-\\varepsilon$, of coordinate $i$ in the real-valued
version of $T(v)$ is at least $\\varepsilon/2$.
\end{definition}

\begin{theorem}[Stable Influence of an Averaged Neighbor Function]
\lean{FABL.RegularUniqueGame.stableInfluence_averagedNeighborFunction_le}
\label{FABL.RegularUniqueGame.stableInfluence_averagedNeighborFunction_le}
\leanok
\uses{FABL.BooleanFunction.toReal, FABL.RegularUniqueGame.LongCodeTable, FABL.RegularUniqueGame.averagedNeighborFunction, FABL.RegularUniqueGame.stableInfluence_comp_permuteInput, FABL.RegularUniqueGame.stableInfluence_uniformExpectation_le, FABL.RegularUniqueGame.twistedLongCode, FABL.stableInfluence}
\difficulty{3}
Let $G$ be a regular unique game of dimension $q$, let a long-code table be given, and
let $\varepsilon\in[0,1]$. For every vertex $v$ and coordinate $i\in\{1,\ldots,q\}$, the
stable influence at parameter $1-\varepsilon$ of the averaged neighbor function at $v$
is at most the uniform expectation over neighbors $w$ of the stable influence, at the
same parameter, of the table function at $w$ in the coordinate corresponding to $i$
under the inverse permutation on the edge from $v$ to $w$.
\end{theorem}

\begin{definition}[Fraction of vertices with notable neighbors]
\lean{FABL.RegularUniqueGame.neighborNotableVertexFraction}
\label{FABL.RegularUniqueGame.neighborNotableVertexFraction}
\leanok
\uses{FABL.RegularUniqueGame.LongCodeTable, FABL.RegularUniqueGame.neighborNotableCoordinates}
Let $G$ be a regular unique game with alphabet size $q$, let a long-code table and a
real parameter $\\varepsilon$ be given, and choose a vertex uniformly from $G$. The
neighbor notable vertex fraction is the probability that the set of notable coordinates
at that vertex's neighbors is nonempty.
\end{definition}

\begin{theorem}[Error bounded by notable-vertex fraction]
\lean{FABL.RegularUniqueGame.error_le_neighborNotableVertexFraction}
\label{FABL.RegularUniqueGame.error_le_neighborNotableVertexFraction}
\leanok
\uses{FABL.HasSmallStableInfluences, FABL.IsDictatorVsNoNotablesTest, FABL.RegularUniqueGame.LongCodeTable, FABL.RegularUniqueGame.averagedNeighborFunction, FABL.RegularUniqueGame.longCodeValue, FABL.RegularUniqueGame.neighborNotableCoordinates, FABL.RegularUniqueGame.neighborNotableVertexFraction, FABL.SignPredicate, FABL.SignTester, FABL.SignTester.acceptanceProbability, FABL.SignTester.acceptanceProbability_mem_Icc}
\difficulty{5}
Let $G$ be a regular unique game with alphabet size $q$, let $$ be a
dictator-versus-no-notables test family with predicate set $$, parameters $,$, error
function $$, and testers $(T_n)_{n>0}$, and let $f$ be a long-code table for $G$. If
$q>0$, $0\le \le1$, and the long-code value of $f$ under the tester $T_q$ is at least
$+2()$, then $()$ is at most the fraction of vertices whose averaged neighbor
functions have a coordinate with stable influence greater than $$.
\end{theorem}

\begin{theorem}[Variance bound for bounded functions]
\lean{FABL.RegularUniqueGame.boundedCubeFunction_variance_le_one}
\label{FABL.RegularUniqueGame.boundedCubeFunction_variance_le_one}
\leanok
\uses{FABL.BoundedCubeFunction, FABL.BoundedCubeFunction.sum_sq_fourierCoeff_le_one, FABL.fourierCoeff, FABL.variance, FABL.variance_eq_sum_sq_fourierCoeff}
\difficulty{3}
Let $f:\{\pm1\}^n\to\mathbb R$ take values in $[-1,1]$. Its variance under the uniform
distribution on the cube is at most $1$.
\end{theorem}

\begin{theorem}[Stable influence bound]
\lean{FABL.RegularUniqueGame.boundedCubeFunction_stableInfluence_le_one}
\label{FABL.RegularUniqueGame.boundedCubeFunction_stableInfluence_le_one}
\leanok
\uses{FABL.BoundedCubeFunction, FABL.BoundedCubeFunction.sum_sq_fourierCoeff_le_one, FABL.fourierCoeff, FABL.stableInfluence}
\difficulty{4}
Let $f:\{\pm1\}^n\to\mathbb R$ take values in $[-1,1]$, let $0\le\rho\le1$, and let
$i\in[n]$. The stable influence of $f$ at coordinate $i$ with parameter $\rho$ is at
most $1$.
\end{theorem}

\begin{theorem}[Threshold probability lower bound]
\lean{FABL.RegularUniqueGame.half_lt_thresholdProbability_of_lt_expect}
\label{FABL.RegularUniqueGame.half_lt_thresholdProbability_of_lt_expect}
\leanok
\difficulty{4}
Let $G$ be a regular unique game, and let $I$ be a finite nonempty set. Suppose
$v_i\in[0,1]$ for every $i\in I$, and let $\varepsilon\ge 0$ satisfy
$\varepsilon<\frac{1}{|I|}\sum_{i\in I}v_i$. Then the fraction of indices $i\in I$ for
which $v_i\ge\varepsilon/2$ is strictly greater than $\varepsilon/2$.
\end{theorem}

\begin{definition}[Notable neighbor probability]
\lean{FABL.RegularUniqueGame.notableNeighborProbability}
\label{FABL.RegularUniqueGame.notableNeighborProbability}
\leanok
\uses{FABL.RegularUniqueGame.LongCodeTable, FABL.RegularUniqueGame.notableCoordinates}
Let $G$ be a regular unique game with alphabet size $q$, let a long-code table and a
real parameter $\\varepsilon$ be given, and fix a vertex $v$ and coordinate $i\\in[q]$.
The notable neighbor probability is the probability, for a uniformly chosen neighbor $w$
of $v$, that the coordinate obtained by applying the inverse edge permutation to $i$ is
notable at $w$.
\end{definition}

\begin{theorem}[Notable neighbor probability bound]
\lean{FABL.RegularUniqueGame.half_epsilon_le_notableNeighborProbability}
\label{FABL.RegularUniqueGame.half_epsilon_le_notableNeighborProbability}
\leanok
\uses{FABL.BooleanFunction.toReal, FABL.BoundedCubeFunction.ofBoolean, FABL.RegularUniqueGame.LongCodeTable, FABL.RegularUniqueGame.averagedNeighborFunction, FABL.RegularUniqueGame.boundedCubeFunction_stableInfluence_le_one, FABL.RegularUniqueGame.half_lt_thresholdProbability_of_lt_expect, FABL.RegularUniqueGame.neighborNotableCoordinates, FABL.RegularUniqueGame.notableCoordinates, FABL.RegularUniqueGame.notableNeighborProbability, FABL.RegularUniqueGame.stableInfluence_averagedNeighborFunction_le, FABL.stableInfluence, FABL.stableInfluence_nonneg}
\difficulty{5}
Let $G$ be a regular unique game, let $T$ be a table of Boolean functions on the label
cube, and let $0\le\epsilon\le1$. For any vertex $v$ and coordinate $i$ notable for the
averaged neighbor function at $v$ (its stable influence at parameter $1-\epsilon$
exceeds $\epsilon$), the probability that a uniformly chosen neighbor has the
corresponding permuted coordinate among its notable coordinates (those with stable
influence at parameter $1-\epsilon$ at least $\epsilon/2$) is at least $\epsilon/2$.
\end{theorem}

\begin{theorem}[Bound on notable neighbor coordinates]
\lean{FABL.RegularUniqueGame.card_neighborNotableCoordinates_le}
\label{FABL.RegularUniqueGame.card_neighborNotableCoordinates_le}
\leanok
\uses{FABL.RegularUniqueGame.LongCodeTable, FABL.RegularUniqueGame.averagedNeighborFunction, FABL.RegularUniqueGame.boundedCubeFunction_variance_le_one, FABL.RegularUniqueGame.neighborNotableCoordinates, FABL.card_stableInfluence_ge_le, FABL.stableInfluence}
\difficulty{4}
Let $G$ be a regular unique game with alphabet size $q$, and let $T$ assign a Boolean
function on the $q$-dimensional sign cube to each vertex. If $0<\varepsilon\le 1$, then
at any vertex the number of coordinates whose stable influence at parameter
$1-\varepsilon$ in the average of its permutation-adjusted neighbors' functions is
strictly greater than $\varepsilon$ is at most $1/\varepsilon^2$.
\end{theorem}

\begin{theorem}[Bound on notable coordinates]
\lean{FABL.RegularUniqueGame.card_notableCoordinates_le}
\proofsource{Analysis-of-Boolean-Functions-by-Ryan-ODonnell}{
  pdf-ffbe9c7b3306-p057-b004,
  pdf-ffbe9c7b3306-p057-b005
}
\label{FABL.RegularUniqueGame.card_notableCoordinates_le}
\leanok
\uses{FABL.BooleanFunction.toReal, FABL.BoundedCubeFunction.ofBoolean, FABL.RegularUniqueGame.LongCodeTable, FABL.RegularUniqueGame.boundedCubeFunction_variance_le_one, FABL.RegularUniqueGame.notableCoordinates, FABL.card_stableInfluence_ge_le, FABL.stableInfluence}
\difficulty{3}
Let $G$ be a regular unique game with alphabet size $q$, and let $T$ assign a Boolean
function on the $q$-dimensional sign cube to each vertex. If $0<\varepsilon\le 1$, then
at every vertex the number of coordinates whose stable influence at parameter
$1-\varepsilon$ is at least $\varepsilon/2$ is at most $2/\varepsilon^2$.
\end{theorem}

\begin{definition}[Decoding candidate set]
\lean{FABL.RegularUniqueGame.decodingCandidates}
\label{FABL.RegularUniqueGame.decodingCandidates}
\leanok
\uses{FABL.RegularUniqueGame.LongCodeTable, FABL.RegularUniqueGame.neighborNotableCoordinates, FABL.RegularUniqueGame.notableCoordinates}
Let $G$ be a regular unique game with alphabet size $q$, let a long-code table and a
real parameter $\\varepsilon$ be given, and fix a vertex $v$. The decoding candidate set
at $v$ is the union of the notable-coordinate set at $v$ and the neighbor
notable-coordinate set at $v$.
\end{definition}

\begin{theorem}[Decoding candidate count bound]
\lean{FABL.RegularUniqueGame.card_decodingCandidates_le}
\label{FABL.RegularUniqueGame.card_decodingCandidates_le}
\leanok
\uses{FABL.RegularUniqueGame.LongCodeTable, FABL.RegularUniqueGame.card_neighborNotableCoordinates_le, FABL.RegularUniqueGame.card_notableCoordinates_le, FABL.RegularUniqueGame.decodingCandidates, FABL.RegularUniqueGame.neighborNotableCoordinates, FABL.RegularUniqueGame.notableCoordinates}
\difficulty{3}
Let $G$ be a regular unique game with label set $in q$, and assign to each vertex a
Boolean-valued function on the Boolean cube. If $0<ε≤1$, then at every vertex the number
of decoding candidates is at most $3/ε^2$.
\end{theorem}

\begin{definition}[Finite candidate expectation]
\lean{FABL.RegularUniqueGame.finiteCandidateExpectation}
\label{FABL.RegularUniqueGame.finiteCandidateExpectation}
\leanok
Let $A$ be a finite set of candidates, let $a_0$ be a default element, and let $f$ be a
real-valued function on the candidate type. The finite candidate expectation is the
uniform average of $f$ over $A$ when $A$ is nonempty, and is $f(a_0)$ when $A$ is empty.
\end{definition}

\begin{definition}[Finite candidate distribution]
\lean{FABL.RegularUniqueGame.finiteCandidatePMF}
\label{FABL.RegularUniqueGame.finiteCandidatePMF}
\leanok
Let $A$ be a finite set of elements of a type with decidable equality, and let $a_0$ be
a default element of that type. The finite candidate distribution is uniform on $A$ when
$A$ is nonempty, and is concentrated at $a_0$ when $A$ is empty.
\end{definition}

\begin{theorem}[Expectation Under a Candidate Label Law]
\lean{FABL.RegularUniqueGame.pmfExpectation_finiteCandidatePMF}
\label{FABL.RegularUniqueGame.pmfExpectation_finiteCandidatePMF}
\leanok
\uses{FABL.RegularUniqueGame.finiteCandidateExpectation, FABL.RegularUniqueGame.finiteCandidatePMF, FABL.pmfExpectation}
\difficulty{5}
Let $\alpha$ be a finite type, let $C\subseteq\alpha$ be a finite set, choose a default
element $a_0\in\alpha$, and let $f:\alpha\to\mathbb{R}$. The expectation of $f$ under
the law that is uniform on $C$ when $C$ is nonempty and concentrated at $a_0$ otherwise
equals the uniform average of $f$ over $C$ when $C$ is nonempty, and equals $f(a_0)$
otherwise.
\end{theorem}

\begin{theorem}[Nonnegativity of candidate expectation]
\lean{FABL.RegularUniqueGame.finiteCandidateExpectation_nonneg}
\label{FABL.RegularUniqueGame.finiteCandidateExpectation_nonneg}
\leanok
\uses{FABL.RegularUniqueGame.finiteCandidateExpectation}
\difficulty{3}
For a finite candidate set, a default value, and a real-valued function on the candidate
domain, if the function is nonnegative at every candidate and at the default value, then
its finite candidate expectation is nonnegative.
\end{theorem}

\begin{theorem}[Lower bound for paired candidate expectation]
\lean{FABL.RegularUniqueGame.finiteCandidatePairExpectation_ge}
\label{FABL.RegularUniqueGame.finiteCandidatePairExpectation_ge}
\leanok
\uses{FABL.RegularUniqueGame.finiteCandidateExpectation, FABL.RegularUniqueGame.finiteCandidateExpectation_nonneg}
\difficulty{5}
Let $A$ and $B$ be finite sets of candidates, let $d$ be a default candidate, and let
$R$ be a decidable relation on pairs of candidates. If $a\in A$, $b\in B$, and $R(a,b)$,
then the iterated finite candidate expectation of the indicator of $R$ is at least
$1/(|A||B|)$.
\end{theorem}

\begin{definition}[Candidate decoding edge value]
\lean{FABL.RegularUniqueGame.candidateDecodingEdgeValue}
\label{FABL.RegularUniqueGame.candidateDecodingEdgeValue}
\leanok
\uses{FABL.RegularUniqueGame.finiteCandidateExpectation}
Let $G$ be a regular unique game with alphabet size $q$, let each vertex be assigned a
finite set of candidate labels in $[q]$, and let $a_0\\in[q]$ be a default label. For a
vertex $v$ and a neighbor $w$ of $v$, the candidate decoding edge value is the iterated
finite candidate expectation of the indicator that the edge permutation maps the target
label at $w$ to the source label at $v$, using the respective candidate sets and default
label.
\end{definition}

\begin{theorem}[Nonnegativity of edge decoding value]
\lean{FABL.RegularUniqueGame.candidateDecodingEdgeValue_nonneg}
\label{FABL.RegularUniqueGame.candidateDecodingEdgeValue_nonneg}
\leanok
\uses{FABL.RegularUniqueGame.candidateDecodingEdgeValue, FABL.RegularUniqueGame.finiteCandidateExpectation_nonneg}
\difficulty{3}
Let $G$ be a regular unique game with label set $in q$. For any family of finite
candidate-label sets, any default label, and any vertex and one of its neighbors, the
expected agreement value obtained by independently choosing labels from the two
candidate sets (using the default for an empty set) is nonnegative.
\end{theorem}

\begin{definition}[Candidate decoding label law]
\lean{FABL.RegularUniqueGame.candidateDecodingLabelLaw}
\label{FABL.RegularUniqueGame.candidateDecodingLabelLaw}
\leanok
\uses{FABL.RegularUniqueGame.finiteCandidatePMF}
Let $G$ be a regular unique game with alphabet $	ext{Fin}(q)$, let each vertex $v$ have
a finite candidate set $C_v\subseteq \text{Fin}(q)$, and choose a default label
$a_0\in\text{Fin}(q)$. For each vertex, the candidate decoding label law is the
probability mass function obtained by applying the finite-candidate law to $C_v$ and
$a_0$.
\end{definition}

\begin{definition}[Candidate decoding assignment law]
\lean{FABL.RegularUniqueGame.candidateDecodingAssignmentLaw}
\label{FABL.RegularUniqueGame.candidateDecodingAssignmentLaw}
\leanok
\uses{FABL.RegularUniqueGame.candidateDecodingLabelLaw, FABL.independentProductPMF}
Let $G$ be a regular unique game with alphabet $\text{Fin}(q)$, let each vertex $v$ have
a finite candidate set $C_v\subseteq \text{Fin}(q)$, and choose a default label
$a_0\in\text{Fin}(q)$. The candidate decoding assignment law on labelings of all
vertices is the independent product of the candidate decoding label laws at the
individual vertices.
\end{definition}

\begin{theorem}[Expected Edge Agreement Under Candidate Decoding]
\lean{FABL.RegularUniqueGame.pmfExpectation_candidateDecodingAssignmentLaw_edge}
\label{FABL.RegularUniqueGame.pmfExpectation_candidateDecodingAssignmentLaw_edge}
\leanok
\uses{FABL.RegularUniqueGame.candidateDecodingAssignmentLaw, FABL.RegularUniqueGame.candidateDecodingEdgeValue, FABL.RegularUniqueGame.candidateDecodingLabelLaw, FABL.RegularUniqueGame.pmfExpectation_finiteCandidatePMF, FABL.RegularUniqueGame.pmfExpectation_independentProductPMF_two_coordinates, FABL.independentProductPMF, FABL.pmfExpectation}
\difficulty{4}
Let $G$ be a regular unique game with label set $\text{Fin}(q)$, choose a candidate
subset of labels for each vertex, and fix a default label. Independently at each vertex,
choose a label uniformly from its candidate subset when that subset is nonempty, and
choose the default otherwise. For any vertex $v$ and neighbor $w$ of $v$, the expected
indicator that the game's permutation on $(v,w)$ sends the label at $w$ to the label at
$v$ equals the nested candidate average of that same agreement indicator, with the
default used for an empty candidate subset.
\end{theorem}

\begin{definition}[Decoding edge value]
\lean{FABL.RegularUniqueGame.decodingEdgeValue}
\label{FABL.RegularUniqueGame.decodingEdgeValue}
\leanok
\uses{FABL.RegularUniqueGame.LongCodeTable, FABL.RegularUniqueGame.candidateDecodingEdgeValue, FABL.RegularUniqueGame.decodingCandidates}
Let $G$ be a regular unique game with alphabet $\text{Fin}(q)$, let $T$ be a table
assigning a Boolean function to each vertex, let $\varepsilon\in\mathbb R$, and choose a
default label $a_0\in\text{Fin}(q)$. For a vertex $v$ and a neighbor $w$ of $v$, the
decoding edge value is the candidate decoding edge value formed from the candidates
decoded from $T$ at parameter $\varepsilon$, using $a_0$.
\end{definition}

\begin{theorem}[Nonnegativity of decoding edge value]
\lean{FABL.RegularUniqueGame.decodingEdgeValue_nonneg}
\label{FABL.RegularUniqueGame.decodingEdgeValue_nonneg}
\leanok
\uses{FABL.RegularUniqueGame.LongCodeTable, FABL.RegularUniqueGame.candidateDecodingEdgeValue_nonneg, FABL.RegularUniqueGame.decodingCandidates, FABL.RegularUniqueGame.decodingEdgeValue}
\difficulty{1}
Let $G$ be a regular unique game with alphabet size $q$, let $T$ be a long-code table,
let $\varepsilon$ be real, and choose a default label in $\{0,\ldots,q-1\}$. For every
vertex and each of its neighbors, the decoding edge value obtained by averaging the
permutation-consistency indicator over the decoding candidate sets at the two endpoints
is nonnegative.
\end{theorem}

\begin{theorem}[Neighbourhood success equals game value]
\lean{FABL.RegularUniqueGame.neighborhoodSuccessValue_eq_cspValue}
\proofsource{Analysis-of-Boolean-Functions-by-Ryan-ODonnell}{pdf-ffbe9c7b3306-p187-b001}
\label{FABL.RegularUniqueGame.neighborhoodSuccessValue_eq_cspValue}
\leanok
\uses{FABL.CSPInstance.value, FABL.FiniteUniqueGame.toCSPInstance, FABL.RegularUniqueGame.edgeError, FABL.RegularUniqueGame.edgeError_eq_one_sub_value, FABL.RegularUniqueGame.edgeFailureIndicator, FABL.RegularUniqueGame.localEdgeError, FABL.RegularUniqueGame.toFiniteUniqueGame}
\difficulty{4}
Let $G$ be a regular unique game and let $\rho$ label each vertex by an alphabet
coordinate. The average, first over a uniform vertex and then over a uniform neighbour
of that vertex, of the indicator that the constraint permutation sends the neighbour's
label to the vertex's label equals the value of the associated constraint satisfaction
instance under $\rho$.
\end{theorem}

\begin{definition}[Candidate decoding value]
\lean{FABL.RegularUniqueGame.candidateDecodingValue}
\label{FABL.RegularUniqueGame.candidateDecodingValue}
\leanok
\uses{FABL.RegularUniqueGame.candidateDecodingEdgeValue}
Let $G$ be a regular unique game with alphabet $\text{Fin}(q)$, let each vertex have a
finite candidate set, and choose a default label in $\text{Fin}(q)$. The candidate
decoding value is the expectation, first over a vertex and then over a neighbor of that
vertex, of the candidate decoding edge value.
\end{definition}

\begin{theorem}[Expected Value of Candidate Decoding]
\lean{FABL.RegularUniqueGame.pmfExpectation_candidateDecodingAssignmentLaw_value}
\label{FABL.RegularUniqueGame.pmfExpectation_candidateDecodingAssignmentLaw_value}
\leanok
\uses{FABL.CSPInstance.value, FABL.FiniteUniqueGame.toCSPInstance, FABL.RegularUniqueGame.candidateDecodingAssignmentLaw, FABL.RegularUniqueGame.candidateDecodingEdgeValue, FABL.RegularUniqueGame.candidateDecodingValue, FABL.RegularUniqueGame.neighborhoodSuccessValue_eq_cspValue, FABL.RegularUniqueGame.pmfExpectation_candidateDecodingAssignmentLaw_edge, FABL.RegularUniqueGame.toFiniteUniqueGame, FABL.fintypeExpectation_pmfExpectation_comm, FABL.pmfExpectation}
\difficulty{4}
Let $G$ be a regular unique game with label set $\text{Fin}(q)$, choose a candidate
subset of labels for each vertex, and fix a default label. Independently decode each
vertex by choosing uniformly from its candidate subset if nonempty, and otherwise using
the default. The expected value of the resulting labeling in the associated finite
unique-game constraint instance equals the average, over vertices and their neighbors,
of the edge agreement value obtained by independently averaging over the two endpoint
candidate subsets, using the default for any empty subset.
\end{theorem}

\begin{theorem}[Candidate decoding value bound]
\lean{FABL.RegularUniqueGame.candidateDecodingValue_le_optimum}
\label{FABL.RegularUniqueGame.candidateDecodingValue_le_optimum}
\leanok
\uses{FABL.CSPInstance.optimum, FABL.CSPInstance.value, FABL.CSPInstance.value_le_optimum, FABL.FiniteUniqueGame.toCSPInstance, FABL.RegularUniqueGame.candidateDecodingAssignmentLaw, FABL.RegularUniqueGame.candidateDecodingValue, FABL.RegularUniqueGame.pmfExpectation_candidateDecodingAssignmentLaw_value, FABL.RegularUniqueGame.toFiniteUniqueGame, FABL.pmfExpectation, FABL.pmfExpectation_const, FABL.pmfExpectation_mono}
\difficulty{4}
Let $G$ be a regular unique game with label set $in q$, where $q>0$. For any family of
finite candidate-label sets and any default label, the expected edge agreement from
independently choosing labels at the vertices, using the default whenever a candidate
set is empty, is at most the optimum value of the associated constraint satisfaction
instance.
\end{theorem}

\begin{definition}[Decoding value]
\lean{FABL.RegularUniqueGame.decodingValue}
\label{FABL.RegularUniqueGame.decodingValue}
\leanok
\uses{FABL.RegularUniqueGame.LongCodeTable, FABL.RegularUniqueGame.candidateDecodingValue, FABL.RegularUniqueGame.decodingCandidates}
Let $G$ be a regular unique game with alphabet $\text{Fin}(q)$, let $T$ be a table
assigning a Boolean function to each vertex, let $\varepsilon\in\mathbb R$, and choose a
default label in $\text{Fin}(q)$. The decoding value is the candidate decoding value for
the candidate sets decoded from $T$ at parameter $\varepsilon$, with that default label.
\end{definition}

\begin{theorem}[Decoding value bounded by optimum]
\lean{FABL.RegularUniqueGame.decodingValue_le_optimum}
\label{FABL.RegularUniqueGame.decodingValue_le_optimum}
\leanok
\uses{FABL.CSPInstance.optimum, FABL.FiniteUniqueGame.toCSPInstance, FABL.RegularUniqueGame.LongCodeTable, FABL.RegularUniqueGame.candidateDecodingValue_le_optimum, FABL.RegularUniqueGame.decodingCandidates, FABL.RegularUniqueGame.decodingValue, FABL.RegularUniqueGame.toFiniteUniqueGame}
\difficulty{1}
Let $G$ be a regular unique game with alphabet size $q>0$, let $T$ be a long-code table,
and let $\varepsilon$ be real. For any default label in $\{0,\ldots,q-1\}$, the decoding
value formed from the notable-coordinate candidate sets is at most the optimum value of
the finite unique-game constraint instance associated with $G$.
\end{theorem}

\begin{theorem}[Reciprocal candidate count bound]
\lean{FABL.RegularUniqueGame.epsilon_sq_div_three_le_inv_card}
\label{FABL.RegularUniqueGame.epsilon_sq_div_three_le_inv_card}
\leanok
\difficulty{4}
Let $C$ be a nonempty finite set and let $$ be a positive real number. If $|C| 
3/^2$, then $^2/3  1/|C|$.
\end{theorem}

\begin{theorem}[Lower bound for edge decoding]
\lean{FABL.RegularUniqueGame.epsilon_pow_four_div_nine_le_decodingEdgeValue}
\label{FABL.RegularUniqueGame.epsilon_pow_four_div_nine_le_decodingEdgeValue}
\leanok
\uses{FABL.RegularUniqueGame.LongCodeTable, FABL.RegularUniqueGame.candidateDecodingEdgeValue, FABL.RegularUniqueGame.card_decodingCandidates_le, FABL.RegularUniqueGame.decodingCandidates, FABL.RegularUniqueGame.decodingEdgeValue, FABL.RegularUniqueGame.epsilon_sq_div_three_le_inv_card, FABL.RegularUniqueGame.finiteCandidatePairExpectation_ge, FABL.RegularUniqueGame.neighborNotableCoordinates, FABL.RegularUniqueGame.notableCoordinates}
\difficulty{5}
Let $G$ be a regular unique game with $q$ labels, let $0<\varepsilon\le 1$, and choose a
table of Boolean functions on the Boolean cube, a default label, a vertex, one of its
neighbors, and a coordinate. Suppose the coordinate has stable influence greater than
$\varepsilon$ in the neighbor-averaged function at the vertex, and its image under the
inverse constraint permutation has stable influence at least $\varepsilon/2$ in the
neighboring table function. Then the decoding value on this directed edge is at least
$\varepsilon^4/9$.
\end{theorem}

\begin{definition}[Decoding vertex value]
\lean{FABL.RegularUniqueGame.decodingVertexValue}
\label{FABL.RegularUniqueGame.decodingVertexValue}
\leanok
\uses{FABL.RegularUniqueGame.LongCodeTable, FABL.RegularUniqueGame.decodingEdgeValue}
Let $G$ be a regular unique game with alphabet $\text{Fin}(q)$, let $T$ be a table
assigning a Boolean function to each vertex, let $\varepsilon\in\mathbb R$, and choose a
default label in $\text{Fin}(q)$. At a vertex $v$, the decoding vertex value is the
expectation over neighbors $w$ of $v$ of the decoding edge value for $(v,w)$.
\end{definition}

\begin{theorem}[Lower bound for decoding value]
\lean{FABL.RegularUniqueGame.epsilon_pow_five_div_eighteen_le_decodingVertexValue}
\label{FABL.RegularUniqueGame.epsilon_pow_five_div_eighteen_le_decodingVertexValue}
\leanok
\uses{FABL.RegularUniqueGame.LongCodeTable, FABL.RegularUniqueGame.decodingEdgeValue, FABL.RegularUniqueGame.decodingEdgeValue_nonneg, FABL.RegularUniqueGame.decodingVertexValue, FABL.RegularUniqueGame.epsilon_pow_four_div_nine_le_decodingEdgeValue, FABL.RegularUniqueGame.half_epsilon_le_notableNeighborProbability, FABL.RegularUniqueGame.neighborNotableCoordinates, FABL.RegularUniqueGame.notableCoordinates, FABL.RegularUniqueGame.notableNeighborProbability}
\difficulty{5}
Let $G$ be a regular unique game with $q$ labels, let $0<\varepsilon\le 1$, and choose a
table of Boolean functions on the Boolean cube, a default label, and a vertex. If at
least one coordinate has stable influence greater than $\varepsilon$ in the
neighbor-averaged, constraint-permuted function at that vertex, then the decoding value
at the vertex is at least $\varepsilon^5/18$.
\end{theorem}

\begin{theorem}[Decoding value lower bound]
\lean{FABL.RegularUniqueGame.error_mul_epsilon_pow_five_div_eighteen_le_decodingValue}
\label{FABL.RegularUniqueGame.error_mul_epsilon_pow_five_div_eighteen_le_decodingValue}
\leanok
\uses{FABL.IsDictatorVsNoNotablesTest, FABL.RegularUniqueGame.LongCodeTable, FABL.RegularUniqueGame.decodingEdgeValue, FABL.RegularUniqueGame.decodingEdgeValue_nonneg, FABL.RegularUniqueGame.decodingValue, FABL.RegularUniqueGame.decodingVertexValue, FABL.RegularUniqueGame.epsilon_pow_five_div_eighteen_le_decodingVertexValue, FABL.RegularUniqueGame.error_le_neighborNotableVertexFraction, FABL.RegularUniqueGame.longCodeValue, FABL.RegularUniqueGame.neighborNotableCoordinates, FABL.RegularUniqueGame.neighborNotableVertexFraction, FABL.SignPredicate, FABL.SignTester}
\difficulty{6}
Let $G$ be a regular unique game with alphabet size $q$, and let a
dictator-versus-no-notables test family have predicate set $$, parameters $,$, error
function $$, and testers $(T_n)_{n>0}$. For a long-code table $f$, suppose $q>0$,
$0<\le1$, and its long-code value under $T_q$ is at least $+2()$. Then the decoding
value obtained from the notable-coordinate candidates at threshold $$, with default
label $0$, is at least $\mathrm{error}()^5/18$.
\end{theorem}

\begin{theorem}[Optimum lower bound from decoding]
\lean{FABL.RegularUniqueGame.error_mul_epsilon_pow_five_div_eighteen_le_optimum}
\label{FABL.RegularUniqueGame.error_mul_epsilon_pow_five_div_eighteen_le_optimum}
\leanok
\uses{FABL.CSPInstance.optimum, FABL.FiniteUniqueGame.toCSPInstance, FABL.IsDictatorVsNoNotablesTest, FABL.RegularUniqueGame.LongCodeTable, FABL.RegularUniqueGame.decodingValue_le_optimum, FABL.RegularUniqueGame.error_mul_epsilon_pow_five_div_eighteen_le_decodingValue, FABL.RegularUniqueGame.longCodeValue, FABL.RegularUniqueGame.toFiniteUniqueGame, FABL.SignPredicate, FABL.SignTester}
\difficulty{1}
Let $G$ be a regular unique game with alphabet size $q$, where $q>0$, and assume the
label set is nonempty. Let a dictator-versus-no-notables test family have predicate set,
parameters $,$, error function $$, and testers $(T_n)_{n>0}$. For a long-code table
$f$, if $0<\le1$ and its long-code value under $T_q$ is at least $+2()$, then the
optimum value of the constraint satisfaction instance associated with $G$ is at least
$\mathrm{error}()^5/18$.
\end{theorem}

\begin{definition}[Coordinates outside a selected pair]
\lean{FABL.CloneRest}
\label{FABL.CloneRest}
\leanok
For a natural number $r$ and indices $i,j$ in the finite set $\\{0,\ldots,r-1\}$, let
the remaining coordinate set consist of those indices $k$ in that set for which $k\ne i$
and $k\ne j$. The indices $i$ and $j$ are allowed to be equal.
\end{definition}

\begin{definition}[Clone pair decomposition]
\lean{FABL.clonePairEquiv}
\label{FABL.clonePairEquiv}
\leanok
\uses{FABL.CloneRest}
Let $r,M$ be natural numbers, and let $i,j\in\text{Fin}(r)$ be distinct. There is a
bijection between assignments of an element of $\text{Fin}(M)$ to each coordinate in
$\text{Fin}(r)$ and data consisting of the values at $i$ and $j$, together with an
assignment of values to all coordinates different from both $i$ and $j$.
\end{definition}

\begin{theorem}[Iterated expectation on a finite product]
\lean{FABL.fintype_expect_prod}
\label{FABL.fintype_expect_prod}
\leanok
\difficulty{3}
Let $\alpha$ and $\beta$ be finite types, and let $v:\alpha\times\beta\to\mathbb R$. The
uniform average of $v$ over $\alpha\times\beta$ equals the iterated uniform average,
first over $\beta$ and then over $\alpha$: $\mathbb
E_{(a,b)\in\alpha\times\beta}v(a,b)=\mathbb E_{a\in\alpha}\,\mathbb
E_{b\in\beta}v(a,b)$.
\end{theorem}

\begin{definition}[Finite initial segment equivalence]
\lean{FABL.finInitialSegmentEquiv}
\label{FABL.finInitialSegmentEquiv}
\leanok
For natural numbers $a$ and $r$ with $a\le r$, the set of indices in $\{0,\ldots,r-1\}$
that are less than $a$ is equivalent to the set $\{0,\ldots,a-1\}$.
\end{definition}

\begin{definition}[Padded tuple equivalence]
\lean{FABL.paddedCloneTupleEquiv}
\label{FABL.paddedCloneTupleEquiv}
\leanok
\uses{FABL.finInitialSegmentEquiv}
Let $a,r,M\in\mathbb{N}$ with $a\le r$. A function from the first $r$ indices to the
first $M$ indices is equivalent to a pair consisting of its restriction to the first $a$
indices and its restriction to the remaining indices.
\end{definition}

\begin{theorem}[Restriction of a padded tuple]
\lean{FABL.paddedCloneTupleEquiv_fst_apply}
\label{FABL.paddedCloneTupleEquiv_fst_apply}
\leanok
\uses{FABL.paddedCloneTupleEquiv}
\difficulty{1}
Let $a,r,M$ be natural numbers with $a\le r$, let $c$ assign an element of
$\mathrm{Fin}(M)$ to each element of $\mathrm{Fin}(r)$, and let $i\in\mathrm{Fin}(a)$.
The first component of the tuple obtained by splitting $c$ into its first $a$
coordinates and its remaining coordinates, evaluated at $i$, is $c(i)$, where $i$ is
viewed in $\mathrm{Fin}(r)$ via the inclusion of the initial segment.
\end{theorem}

\begin{theorem}[Expectation under uniform padding]
\lean{FABL.uniformPaddedCloneExpectation}
\label{FABL.uniformPaddedCloneExpectation}
\leanok
\uses{FABL.fintype_expect_prod, FABL.paddedCloneTupleEquiv}
\difficulty{4}
Let $a,r,M$ be natural numbers with $a\le r$ and $M>0$, and let $f$ be a real-valued
function on the functions from $\{0,\ldots,a-1\}$ to $\{0,\ldots,M-1\}$. For a uniformly
chosen function $c$ from $\{0,\ldots,r-1\}$ to $\{0,\ldots,M-1\}$, the expected value of
$f$ applied to the restriction of $c$ to the first $a$ positions equals the uniform
expectation of $f$ over functions on the first $a$ positions.
\end{theorem}

\begin{theorem}[Collision probability for a uniform tuple]
\lean{FABL.uniformCloneTuple_pair_eq_probability}
\label{FABL.uniformCloneTuple_pair_eq_probability}
\leanok
\uses{FABL.CloneRest, FABL.clonePairEquiv, FABL.fintype_expect_prod}
\difficulty{5}
Let $r$ and $M$ be natural numbers with $M>0$, and let $i,j\in\{0,ldots,r-1\}$ be
distinct. If a tuple of length $r$ is chosen uniformly from ${0,ldots,M-1}^r$, the
probability that its entries at $i$ and $j$ agree is $1/M$.
\end{theorem}

\begin{definition}[Tuple collision probability]
\lean{FABL.cloneTupleCollisionProbability}
\label{FABL.cloneTupleCollisionProbability}
\leanok
For natural numbers $r$ and $M$, the collision probability is the probability that a
uniformly chosen function from the first $r$ indices to the first $M$ indices is not
injective.
\end{definition}

\begin{theorem}[Collision probability for random tuples]
\lean{FABL.cloneTupleCollisionProbability_le}
\label{FABL.cloneTupleCollisionProbability_le}
\leanok
\uses{FABL.cloneTupleCollisionProbability, FABL.uniformCloneTuple_pair_eq_probability}
\difficulty{6}
Let $r,M$ be natural numbers with $M>0$, and choose a function uniformly from all maps
from an $r$-element set to an $M$-element set. The probability that this function is not
injective is at most $r^2/M$.
\end{theorem}

\begin{definition}[Repeated query scope]
\lean{FABL.RepeatedQueryScope}
\statementsource{Analysis-of-Boolean-Functions-by-Ryan-ODonnell}{pdf-ffbe9c7b3306-p202-b003}
\label{FABL.RepeatedQueryScope}
\leanok
For a type $X$, a repeated query scope consists of a natural number $r$ of query
positions and a choice of an element of $X$ at each position; the same element may be
chosen at multiple positions.
\end{definition}

\begin{definition}[Cloned query]
\lean{FABL.RepeatedQueryScope.clonedQuery}
\statementsource{Analysis-of-Boolean-Functions-by-Ryan-ODonnell}{pdf-ffbe9c7b3306-p202-b003}
\label{FABL.RepeatedQueryScope.clonedQuery}
\leanok
Let a repeated query scope on $X$ have $r$ positions. For a natural number $M$ and an
assignment of a clone index in $\\{0,\\ldots,M-1\\}$ to each position, the cloned query
at each position is the ordered pair consisting of its queried element of $X$ and its
assigned clone index.
\end{definition}

\begin{theorem}[Injectivity of the Cloned Query Map]
\lean{FABL.RepeatedQueryScope.clonedQuery_injective_of_cloneTuple_injective}
\label{FABL.RepeatedQueryScope.clonedQuery_injective_of_cloneTuple_injective}
\leanok
\uses{FABL.RepeatedQueryScope.clonedQuery}
\difficulty{1}
Let $X$ be a type, let a repeated query scope on $X$ have arity $r$, let $M$ be a
natural number, and assign to each of its $r$ positions a clone label in
$\{1,\ldots,M\}$. If the assignment of clone labels is injective, then the map sending
each position to its original query together with its clone label is injective.
\end{theorem}

\begin{definition}[Cloned scope collision probability]
\lean{FABL.RepeatedQueryScope.clonedScopeCollisionProbability}
\statementsource{Analysis-of-Boolean-Functions-by-Ryan-ODonnell}{pdf-ffbe9c7b3306-p202-b003}
\label{FABL.RepeatedQueryScope.clonedScopeCollisionProbability}
\leanok
\uses{FABL.RepeatedQueryScope.clonedQuery}
Let $X$ be a type, let $S$ be a repeated-query scope on $X$, and let $M$ be a natural
number. The cloned-scope collision probability is the probability, for uniformly chosen
clone labels in $	ext{Fin}(M)$ at each query position, that the resulting cloned query
map is not injective.
\end{definition}

\begin{theorem}[Bound on Cloned Scope Collision Probability]
\lean{FABL.RepeatedQueryScope.clonedScopeCollisionProbability_le}
\proofsource{Analysis-of-Boolean-Functions-by-Ryan-ODonnell}{pdf-ffbe9c7b3306-p202-b003}
\label{FABL.RepeatedQueryScope.clonedScopeCollisionProbability_le}
\leanok
\uses{FABL.RepeatedQueryScope.clonedQuery, FABL.RepeatedQueryScope.clonedQuery_injective_of_cloneTuple_injective, FABL.RepeatedQueryScope.clonedScopeCollisionProbability, FABL.cloneTupleCollisionProbability, FABL.cloneTupleCollisionProbability_le}
\difficulty{4}
Let $X$ be a type with decidable equality, let a repeated query scope on $X$ have arity
$r$, and let $M$ be a positive integer. When each scope position is assigned an
independent uniformly random clone label from $\{1,\ldots,M\}$, the probability that the
resulting map from positions to the original query and clone-label pairs is not
injective is at most $r^2/M$.
\end{theorem}

\begin{theorem}[Expectation difference and failure probability]
\lean{FABL.WeightedConstraintExperiment.abs_uniformExpectation_subtype_le_failureProbability}
\label{FABL.WeightedConstraintExperiment.abs_uniformExpectation_subtype_le_failureProbability}
\leanok
\difficulty{6}
Let $\Omega$ be a nonempty finite set, let $G\subseteq\Omega$ be nonempty, and let
$f:\Omega\to\mathbb{R}$ take values in $[0,1]$. The absolute difference between the
uniform mean of $f$ on $\Omega$ and its uniform mean on $G$ is at most the fraction of
elements of $\Omega$ outside $G$.
\end{theorem}

\begin{definition}[Clone outcome]
\lean{FABL.WeightedConstraintExperiment.CloneOutcome}
\statementsource{Analysis-of-Boolean-Functions-by-Ryan-ODonnell}{pdf-ffbe9c7b3306-p198-b006}
\label{FABL.WeightedConstraintExperiment.CloneOutcome}
\leanok
Let $E$ be a weighted constraint experiment over a constraint template with predicate
family $\Psi$ and outcome space $\mathrm{Outcome}(E)$. For a natural number $M$,
a clone outcome is a pair consisting of an experiment outcome and a choice of a label in
$\text{Fin}(M)$ for each query position of the predicate selected by that outcome.
\end{definition}

\begin{definition}[Cloning a constraint experiment]
\lean{FABL.WeightedConstraintExperiment.clone}
\statementsource{Analysis-of-Boolean-Functions-by-Ryan-ODonnell}{
  pdf-ffbe9c7b3306-p198-b006,
  pdf-ffbe9c7b3306-p202-b003
}
\label{FABL.WeightedConstraintExperiment.clone}
\leanok
\uses{FABL.WeightedConstraintExperiment, FABL.WeightedConstraintExperiment.CloneOutcome, FABL.uniformPMF}
Let $E$ be a weighted constraint experiment with variable set $V$, and let $M>0$ be a
natural number. Its $M$-clone is the weighted constraint experiment on variables
$V\times\text{Fin}(M)$ obtained by first sampling an outcome from $E$ and then
independently choosing a uniformly random clone label in $\text{Fin}(M)$ for each query
position. It retains the original predicate and replaces each queried variable by the
pair of that variable and its chosen clone label.
\end{definition}

\begin{definition}[Cloned constraint payoff]
\lean{FABL.WeightedConstraintExperiment.clonePayoff}
\statementsource{Analysis-of-Boolean-Functions-by-Ryan-ODonnell}{
  pdf-ffbe9c7b3306-p186-b001,
  pdf-ffbe9c7b3306-p198-b006
}
\label{FABL.WeightedConstraintExperiment.clonePayoff}
\leanok
Let $E$ be a weighted constraint experiment, let $M$ be a natural number, let an
assignment give a value in the domain $D$ to each variable in $V\times\text{Fin}(M)$,
and fix an outcome and clone labels for its query positions. The cloned constraint
payoff is $1$ if the predicate selected by the outcome is satisfied by the assigned
values at the corresponding original variables and clone labels, and $0$ otherwise.
\end{definition}

\begin{theorem}[Clone payoff bounds]
\lean{FABL.WeightedConstraintExperiment.clonePayoff_mem_Icc}
\label{FABL.WeightedConstraintExperiment.clonePayoff_mem_Icc}
\leanok
\uses{FABL.WeightedConstraintExperiment.clonePayoff}
\difficulty{1}
In a weighted constraint experiment, fix a positive integer $M$, an assignment of values
to variable-clone pairs, an outcome, and a choice of one clone for each query position
of that outcome. The payoff, equal to $1$ when the selected predicate is satisfied by
the assigned values and $0$ otherwise, lies in $[0,1]$.
\end{theorem}

\begin{definition}[Clone collision probability]
\lean{FABL.WeightedConstraintExperiment.cloneCollisionProbability}
\statementsource{Analysis-of-Boolean-Functions-by-Ryan-ODonnell}{pdf-ffbe9c7b3306-p202-b003}
\label{FABL.WeightedConstraintExperiment.cloneCollisionProbability}
\leanok
\uses{FABL.WeightedConstraintExperiment.clone, FABL.pmfExpectation}
Let $E$ be a weighted constraint experiment and let $M>0$ be a natural number. Its clone
collision probability is the probability under the $M$-clone experiment that the scope
map from query positions to cloned variables is not injective.
\end{definition}

\begin{theorem}[Cloned value as an iterated expectation]
\lean{FABL.WeightedConstraintExperiment.clone_value_eq_iterated}
\label{FABL.WeightedConstraintExperiment.clone_value_eq_iterated}
\leanok
\uses{FABL.WeightedConstraintExperiment.clone, FABL.WeightedConstraintExperiment.clonePayoff, FABL.WeightedConstraintExperiment.value, FABL.pmfExpectation, FABL.pmfExpectation_bind, FABL.pmfExpectation_map, FABL.pmfExpectation_uniformPMF_eq_expect}
\difficulty{3}
Let $E$ be a finite weighted constraint experiment, let $M>0$, and let $a$ assign values
to variable-clone pairs. The value of the cloned experiment under $a$ equals the
expectation over an outcome drawn from the original experiment of the uniform mean, over
all clone choices for that outcome's query positions, of the cloned payoff.
\end{theorem}

\begin{definition}[Mean sign across clones]
\lean{FABL.WeightedConstraintExperiment.cloneMean}
\statementsource{Analysis-of-Boolean-Functions-by-Ryan-ODonnell}{pdf-ffbe9c7b3306-p198-b006}
\label{FABL.WeightedConstraintExperiment.cloneMean}
\leanok
\uses{FABL.Sign, FABL.signValue}
Let $E$ be a weighted constraint experiment, let $M$ be a natural number, let an
assignment give a sign to each pair of a variable in $V$ and a clone label in
$\text{Fin}(M)$, and fix a variable $v\in V$. The clone mean at $v$ is the average
numerical value of the signs assigned to $(v,j)$ as $j$ ranges uniformly over
$\text{Fin}(M)$.
\end{definition}

\begin{theorem}[Clone mean lies in the unit interval]
\lean{FABL.WeightedConstraintExperiment.cloneMean_mem_Icc}
\label{FABL.WeightedConstraintExperiment.cloneMean_mem_Icc}
\leanok
\uses{FABL.Sign, FABL.WeightedConstraintExperiment.cloneMean, FABL.signValue, FABL.signValue_eq_neg_one_or_one}
\difficulty{4}
Let $M>0$, assign to each pair consisting of a variable and one of its $M$ clones a sign
in $\{ -1,1\}$, and fix a variable. The uniform mean of its $M$ assigned signs lies in
$[-1,1]$.
\end{theorem}

\begin{theorem}[Clone value eq randomized clone mean]
\lean{FABL.WeightedConstraintExperiment.clone_value_eq_randomizedCloneMean}
\statementsource{Analysis-of-Boolean-Functions-by-Ryan-ODonnell}{pdf-ffbe9c7b3306-p198-b006}
\label{FABL.WeightedConstraintExperiment.clone_value_eq_randomizedCloneMean}
\leanok
\uses{FABL.RegularUniqueGame.multilinearExtension_signValues, FABL.RegularUniqueGame.pmfExpectation_multilinearExtension_productLaw, FABL.Sign, FABL.SignPredicate.indicator, FABL.SignPredicate.multilinearExtension, FABL.SignPredicate.randomizedValue, FABL.SignPredicate.randomizedValue_eq_multilinearExtension, FABL.SignTester, FABL.SignTester.cspTemplate, FABL.WeightedConstraintExperiment, FABL.WeightedConstraintExperiment.clone, FABL.WeightedConstraintExperiment.cloneMean, FABL.WeightedConstraintExperiment.cloneMean_mem_Icc, FABL.WeightedConstraintExperiment.clone_value_eq_iterated, FABL.WeightedConstraintExperiment.value, FABL.independentProductPMF, FABL.independentProductPMF_uniform, FABL.pmfExpectation, FABL.pmfExpectation_uniformPMF_eq_expect, FABL.signValue, FABL.uniformPMF}
\difficulty{5}
% dataset-pipeline: stub -- prose pending build_informal_statements.py
Statement pending regeneration.
\end{theorem}

\begin{definition}[Repeated not equal predicate]
\lean{FABL.WeightedConstraintExperiment.repeatedNotEqualPredicate}
\label{FABL.WeightedConstraintExperiment.repeatedNotEqualPredicate}
\leanok
\uses{FABL.SignPredicate}
For any weighted constraint experiment, define a predicate on two signs that accepts
exactly when the two signs are unequal.
\end{definition}

\begin{definition}[Repeated not equal check]
\lean{FABL.WeightedConstraintExperiment.repeatedNotEqualCheck}
\label{FABL.WeightedConstraintExperiment.repeatedNotEqualCheck}
\leanok
\uses{FABL.SignQueryCheck, FABL.WeightedConstraintExperiment.repeatedNotEqualPredicate}
For any weighted constraint experiment, define a check on one-dimensional sign vectors
using the predicate that accepts when its two inputs differ. Both queries are the
constant vector with value $1$.
\end{definition}

\begin{definition}[Repeated not equal tester]
\lean{FABL.WeightedConstraintExperiment.repeatedNotEqualTester}
\label{FABL.WeightedConstraintExperiment.repeatedNotEqualTester}
\leanok
\uses{FABL.SignTester, FABL.WeightedConstraintExperiment.repeatedNotEqualCheck}
For any weighted constraint experiment, define a tester with a single possible seed,
chosen with probability $1$, whose check is the repeated not equal check on
one-dimensional sign vectors.
\end{definition}

\begin{definition}[Repeated not equal experiment]
\lean{FABL.WeightedConstraintExperiment.repeatedNotEqualExperiment}
\label{FABL.WeightedConstraintExperiment.repeatedNotEqualExperiment}
\leanok
\uses{FABL.Sign, FABL.SignTester.cspTemplate, FABL.WeightedConstraintExperiment, FABL.WeightedConstraintExperiment.repeatedNotEqualTester}
For any weighted constraint experiment, define an experiment over signs with a single
outcome of probability $1$. Its constraint is the repeated not equal predicate, and both
query positions refer to the sole variable.
\end{definition}

\begin{theorem}[Value of the repeated not-equal experiment]
\lean{FABL.WeightedConstraintExperiment.repeatedNotEqualExperiment_value}
\label{FABL.WeightedConstraintExperiment.repeatedNotEqualExperiment_value}
\leanok
\uses{FABL.Sign, FABL.SignTester.cspTemplate, FABL.WeightedConstraintExperiment.payoff, FABL.WeightedConstraintExperiment.repeatedNotEqualCheck, FABL.WeightedConstraintExperiment.repeatedNotEqualExperiment, FABL.WeightedConstraintExperiment.repeatedNotEqualPredicate, FABL.WeightedConstraintExperiment.repeatedNotEqualTester, FABL.WeightedConstraintExperiment.value, FABL.pmfExpectation}
\difficulty{3}
For every assignment of a sign in ${\pm1}$ to the single variable, the value of the
repeated-not-equal experiment is $0$.
\end{theorem}

\begin{theorem}[Repeated not-equal experiment optimum]
\lean{FABL.WeightedConstraintExperiment.repeatedNotEqualExperiment_optimum}
\label{FABL.WeightedConstraintExperiment.repeatedNotEqualExperiment_optimum}
\leanok
\uses{FABL.WeightedConstraintExperiment.optimalAssignment, FABL.WeightedConstraintExperiment.optimum, FABL.WeightedConstraintExperiment.repeatedNotEqualExperiment, FABL.WeightedConstraintExperiment.repeatedNotEqualExperiment_value}
\difficulty{1}
The optimum of the repeated-not-equal weighted constraint experiment is $0$.
\end{theorem}

\begin{definition}[Repeated not equal clone assignment]
\lean{FABL.WeightedConstraintExperiment.repeatedNotEqualCloneAssignment}
\label{FABL.WeightedConstraintExperiment.repeatedNotEqualCloneAssignment}
\leanok
\uses{FABL.Sign}
For any weighted constraint experiment, define an assignment on pairs consisting of the
sole variable and one of two positions: it assigns $-1$ at position $0$ and $1$ at
position $1$.
\end{definition}

\begin{theorem}[Injectivity of a two-clone assignment]
\lean{FABL.WeightedConstraintExperiment.repeatedNotEqualCloneAssignment_injective}
\label{FABL.WeightedConstraintExperiment.repeatedNotEqualCloneAssignment_injective}
\leanok
\uses{FABL.WeightedConstraintExperiment.repeatedNotEqualCloneAssignment}
\difficulty{4}
Assign the sign $-1$ to clone $0$ and the sign $1$ to clone $1$ of a single variable.
The resulting map from the two clones to their assigned signs is injective.
\end{theorem}

\begin{theorem}[Value of the cloned experiment]
\lean{FABL.WeightedConstraintExperiment.repeatedNotEqual_clone_value}
\label{FABL.WeightedConstraintExperiment.repeatedNotEqual_clone_value}
\leanok
\uses{FABL.SignTester.cspTemplate, FABL.WeightedConstraintExperiment.clone, FABL.WeightedConstraintExperiment.clonePayoff, FABL.WeightedConstraintExperiment.clone_value_eq_iterated, FABL.WeightedConstraintExperiment.repeatedNotEqualCheck, FABL.WeightedConstraintExperiment.repeatedNotEqualCloneAssignment, FABL.WeightedConstraintExperiment.repeatedNotEqualCloneAssignment_injective, FABL.WeightedConstraintExperiment.repeatedNotEqualExperiment, FABL.WeightedConstraintExperiment.repeatedNotEqualPredicate, FABL.WeightedConstraintExperiment.repeatedNotEqualTester, FABL.WeightedConstraintExperiment.value, FABL.pmfExpectation, FABL.pmfExpectation_pure, FABL.uniformCloneTuple_pair_eq_probability}
\difficulty{5}
In the experiment obtained by making two independent clones of each queried variable in
the repeated-not-equal experiment, assign sign $-1$ to clone $0$ and sign $1$ to clone
$1$. The value of this assignment is $1/2$.
\end{theorem}

\begin{theorem}[Exercise7 31a repeated scope counterexample]
\lean{FABL.WeightedConstraintExperiment.exercise7_31a_repeatedScope_counterexample}
\statementsource{Analysis-of-Boolean-Functions-by-Ryan-ODonnell}{pdf-ffbe9c7b3306-p202-b003}
\label{FABL.WeightedConstraintExperiment.exercise7_31a_repeatedScope_counterexample}
\leanok
\uses{FABL.WeightedConstraintExperiment.clone, FABL.WeightedConstraintExperiment.optimum, FABL.WeightedConstraintExperiment.repeatedNotEqualCloneAssignment, FABL.WeightedConstraintExperiment.repeatedNotEqualExperiment, FABL.WeightedConstraintExperiment.repeatedNotEqualExperiment_optimum, FABL.WeightedConstraintExperiment.repeatedNotEqual_clone_value, FABL.WeightedConstraintExperiment.value_le_optimum}
\difficulty{1}
% dataset-pipeline: stub -- prose pending build_informal_statements.py
Statement pending regeneration.
\end{theorem}

\begin{definition}[Constant clone assignment]
\lean{FABL.WeightedConstraintExperiment.constantCloneAssignment}
\label{FABL.WeightedConstraintExperiment.constantCloneAssignment}
\leanok
Let $D$ and $V$ be types, let $\Psi$ be a constraint satisfaction template on $D$, and
let $E$ be a weighted constraint experiment with these data. For an assignment $a:V\to
D$ and a natural number $M$, define its clone assignment on $V\times\{0,\ldots,M-1\}$ by
sending each pair $(v,i)$ to $a(v)$.
\end{definition}

\begin{theorem}[Value under a constant clone assignment]
\lean{FABL.WeightedConstraintExperiment.clone_value_constantCloneAssignment}
\label{FABL.WeightedConstraintExperiment.clone_value_constantCloneAssignment}
\leanok
\uses{FABL.WeightedConstraintExperiment.clone, FABL.WeightedConstraintExperiment.constantCloneAssignment, FABL.WeightedConstraintExperiment.payoff, FABL.WeightedConstraintExperiment.value, FABL.pmfExpectation, FABL.pmfExpectation_bind, FABL.pmfExpectation_map, FABL.pmfExpectation_uniformPMF_eq_expect}
\difficulty{4}
Let $E$ be a finite weighted constraint experiment, let $M>0$, and let $a$ assign a
value to each variable. Extend $a$ to variable-clone pairs by giving every clone of a
variable the same value. The value of the cloned experiment under this extended
assignment equals the value of $E$ under $a$.
\end{theorem}

\begin{theorem}[Clone collision probability bound]
\lean{FABL.WeightedConstraintExperiment.cloneCollisionProbability_le}
\label{FABL.WeightedConstraintExperiment.cloneCollisionProbability_le}
\leanok
\uses{FABL.RepeatedQueryScope, FABL.RepeatedQueryScope.clonedQuery, FABL.RepeatedQueryScope.clonedScopeCollisionProbability_le, FABL.WeightedConstraintExperiment.CloneOutcome, FABL.WeightedConstraintExperiment.clone, FABL.WeightedConstraintExperiment.cloneCollisionProbability, FABL.pmfExpectation, FABL.pmfExpectation_bind, FABL.pmfExpectation_const, FABL.pmfExpectation_map, FABL.pmfExpectation_mono, FABL.pmfExpectation_uniformPMF_eq_expect, FABL.uniformPMF}
\difficulty{5}
Let $E$ be a finite weighted constraint experiment, let $r$ be a natural number, and let
$M>0$. Suppose every predicate selected by an outcome of $E$ has arity at most $r$. If
each constraint occurrence is assigned independently and uniformly one of $M$ clones of
each variable, then the probability that two positions in the resulting ordered query
tuple receive the same variable-clone pair is at most $r^2/M$.
\end{theorem}

\begin{theorem}[Value preserved by restricting predicates]
\lean{FABL.RegularUniqueGame.longCodeAllowedWeightedExperiment_clone_value}
\label{FABL.RegularUniqueGame.longCodeAllowedWeightedExperiment_clone_value}
\leanok
\uses{FABL.RegularUniqueGame.longCodeAllowedWeightedExperiment, FABL.RegularUniqueGame.longCodeExperimentLaw, FABL.RegularUniqueGame.longCodeWeightedExperiment, FABL.Sign, FABL.SignPredicate, FABL.SignTester, FABL.SignTester.UsesPredicates, FABL.WeightedConstraintExperiment.clone, FABL.WeightedConstraintExperiment.clone_value_eq_iterated, FABL.WeightedConstraintExperiment.value, FABL.pmfExpectation}
\difficulty{4}
Let $G$ be a regular unique game, let $T$ be a sign tester whose predicates all belong
to a finite set $\Psi$ of sign predicates, and let $M$ be a positive integer. For every
assignment of signs to pairs consisting of a vertex-query pair and a clone index, the
value of the cloned long-code experiment restricted to $\Psi$ equals the value of the
unrestricted cloned long-code experiment.
\end{theorem}

\begin{theorem}[Value preserved by cloning]
\lean{FABL.RegularUniqueGame.clonedLongCodeWeightedExperiment_value_constant}
\label{FABL.RegularUniqueGame.clonedLongCodeWeightedExperiment_value_constant}
\leanok
\uses{FABL.RegularUniqueGame.LongCodeTable, FABL.RegularUniqueGame.longCodeAssignment, FABL.RegularUniqueGame.longCodeExperimentValue, FABL.RegularUniqueGame.longCodeWeightedExperiment, FABL.RegularUniqueGame.longCodeWeightedExperiment_value, FABL.SignTester, FABL.WeightedConstraintExperiment.clone, FABL.WeightedConstraintExperiment.clone_value_constantCloneAssignment, FABL.WeightedConstraintExperiment.constantCloneAssignment, FABL.WeightedConstraintExperiment.value}
\difficulty{1}
Let $G$ be a regular unique game, let $T$ be a long-code table, and let a tester have a
finite seed space with a probability law and a sign-valued query check at each seed. For
every positive integer $M$, the value of the cloned long-code constraint experiment
under the assignment that gives every clone the original variable's sign equals the
original long-code experiment's acceptance value for $T$.
\end{theorem}

\begin{definition}[Clone mean long code table]
\lean{FABL.RegularUniqueGame.cloneMeanLongCodeTable}
\label{FABL.RegularUniqueGame.cloneMeanLongCodeTable}
\leanok
\uses{FABL.RegularUniqueGame.BoundedLongCodeTable, FABL.Sign, FABL.WeightedConstraintExperiment.cloneMean, FABL.WeightedConstraintExperiment.cloneMean_mem_Icc}
Let $G$ be a regular unique game with alphabet size $q$, let $M$ be a positive natural
number, and let $A:(V(G)\times\{-1,1\}^q)\times\{0,\ldots,M-1\}\to\{\pm1\}$ be an
assignment. The clone mean long code table assigns to each vertex the function on the
sign cube whose value at a query is the mean of the assignment over the $M$ clones at
that vertex and query; each such function takes values in $[-1,1]$.
\end{definition}

\begin{definition}[Bounded operational check value]
\lean{FABL.RegularUniqueGame.boundedOperationalCheckValue}
\label{FABL.RegularUniqueGame.boundedOperationalCheckValue}
\leanok
\uses{FABL.RegularUniqueGame.BoundedLongCodeTable, FABL.RegularUniqueGame.boundedTwistedLongCode, FABL.SignPredicate.randomizedValue, FABL.SignQueryCheck, FABL.independentProductPMF, FABL.permuteInput, FABL.pmfExpectation, FABL.uniformPMF}
Let $G$ be a regular unique game with alphabet size $q$, and let $T$ be a table
assigning to each vertex a function from the $q$-dimensional sign cube to $[-1,1]$. Fix
a vertex $v$ and a sign query check consisting of a Boolean predicate on sign inputs and
one cube query for each input position. The bounded operational check value is the
expected randomized value of the predicate when each position independently samples a
uniformly random neighbor of $v$, using the twisted table value at the neighbor and the
queried coordinate, together with the corresponding bounded table value at the permuted
query.
\end{definition}

\begin{theorem}[Operational and averaged check values]
\lean{FABL.RegularUniqueGame.boundedOperationalCheckValue_eq_randomizedValue}
\label{FABL.RegularUniqueGame.boundedOperationalCheckValue_eq_randomizedValue}
\leanok
\uses{FABL.RegularUniqueGame.BoundedLongCodeTable, FABL.RegularUniqueGame.boundedAveragedNeighborFunction, FABL.RegularUniqueGame.boundedOperationalCheckValue, FABL.RegularUniqueGame.boundedTwistedLongCode, FABL.RegularUniqueGame.pmfExpectation_multilinearExtension_productLaw, FABL.SignPredicate.multilinearExtension, FABL.SignPredicate.randomizedValue_eq_multilinearExtension, FABL.SignQueryCheck, FABL.SignQueryCheck.randomizedValue, FABL.independentProductPMF, FABL.pmfExpectation, FABL.pmfExpectation_uniformPMF_eq_expect, FABL.uniformPMF}
\difficulty{4}
Let $G$ be a regular unique game, let each vertex be assigned a bounded real-valued
function on its Boolean cube, and fix a vertex and a randomized sign query check. The
check’s operational value, averaging over independent uniform neighbor choices for its
query positions, equals its randomized value when applied to the function obtained by
averaging the neighbor-transformed table functions at that vertex.
\end{theorem}

\begin{theorem}[Cloned value equals bounded-table value]
\lean{FABL.RegularUniqueGame.clonedLongCode_value_eq_boundedLongCodeValue}
\label{FABL.RegularUniqueGame.clonedLongCode_value_eq_boundedLongCodeValue}
\leanok
\uses{FABL.RegularUniqueGame.boundedLongCodeValue, FABL.RegularUniqueGame.boundedOperationalCheckValue, FABL.RegularUniqueGame.boundedOperationalCheckValue_eq_randomizedValue, FABL.RegularUniqueGame.cloneMeanLongCodeTable, FABL.RegularUniqueGame.longCodeExperimentLaw, FABL.RegularUniqueGame.longCodeExperimentVariable, FABL.RegularUniqueGame.longCodeWeightedExperiment, FABL.Sign, FABL.SignPredicate.randomizedValue, FABL.SignTester, FABL.SignTester.cspTemplate, FABL.WeightedConstraintExperiment.clone, FABL.WeightedConstraintExperiment.cloneMean, FABL.WeightedConstraintExperiment.cloneMean_mem_Icc, FABL.WeightedConstraintExperiment.clone_value_eq_randomizedCloneMean, FABL.WeightedConstraintExperiment.value, FABL.pmfExpectation, FABL.pmfExpectation_bind, FABL.pmfExpectation_map, FABL.pmfExpectation_uniformPMF_eq_expect}
\difficulty{5}
Let $G$ be a regular unique game, let a tester have a finite seed space with a
probability law and a sign-valued query check at each seed, and let $M$ be a positive
integer. For any sign assignment to pairs consisting of a vertex-query pair and a clone,
the value of the cloned long-code constraint experiment equals the tester's acceptance
probability averaged over vertices when each vertex's bounded function is the mean sign
of its $M$ clones at each query.
\end{theorem}

\begin{definition}[Bounded neighbor notable coordinates]
\lean{FABL.RegularUniqueGame.boundedNeighborNotableCoordinates}
\label{FABL.RegularUniqueGame.boundedNeighborNotableCoordinates}
\leanok
\uses{FABL.RegularUniqueGame.BoundedLongCodeTable, FABL.RegularUniqueGame.boundedAveragedNeighborFunction, FABL.stableInfluence}
Let $G$ be a regular unique game with alphabet size $q$, let $T$ be a table of functions
from the $q$-dimensional sign cube to $[-1,1]$, let $\varepsilon\in\mathbb R$, and fix a
vertex $v$. The notable coordinates are those $i\in\{1,\ldots,q\}$ for which
$\varepsilon$ is strictly less than the stable influence, at noise parameter
$1-\varepsilon$, of coordinate $i$ in the averaged neighbor function of $T$ at $v$.
\end{definition}

\begin{definition}[Bounded notable coordinates]
\lean{FABL.RegularUniqueGame.boundedNotableCoordinates}
\label{FABL.RegularUniqueGame.boundedNotableCoordinates}
\leanok
\uses{FABL.RegularUniqueGame.BoundedLongCodeTable, FABL.stableInfluence}
Let $G$ be a regular unique game with alphabet size $q$, let $T$ be a table of functions
from the $q$-dimensional sign cube to $[-1,1]$, let $\varepsilon\in\mathbb R$, and fix a
vertex $v$. The notable coordinates are those $i\in\{1,\ldots,q\}$ for which the stable
influence, at noise parameter $1-\varepsilon$, of coordinate $i$ in the function $T(v)$
is at least $\varepsilon/2$.
\end{definition}

\begin{theorem}[Stable Influence of a Bounded Neighbor Average]
\lean{FABL.RegularUniqueGame.stableInfluence_boundedAveragedNeighborFunction_le}
\label{FABL.RegularUniqueGame.stableInfluence_boundedAveragedNeighborFunction_le}
\leanok
\uses{FABL.RegularUniqueGame.BoundedLongCodeTable, FABL.RegularUniqueGame.boundedAveragedNeighborFunction, FABL.RegularUniqueGame.boundedTwistedLongCode, FABL.RegularUniqueGame.stableInfluence_comp_permuteInput, FABL.RegularUniqueGame.stableInfluence_uniformExpectation_le, FABL.stableInfluence}
\difficulty{3}
Let $G$ be a regular unique game of dimension $q$, let a table of functions taking
values in $[-1,1]$ on the Boolean cube be given, and let $\varepsilon\in[0,1]$. For
every vertex $v$ and coordinate $i\in\{1,\ldots,q\}$, the stable influence at parameter
$1-\varepsilon$ of the bounded averaged neighbor function at $v$ is at most the uniform
expectation over neighbors $w$ of the stable influence, at the same parameter, of the
table function at $w$ in the coordinate corresponding to $i$ under the inverse edge
permutation.
\end{theorem}

\begin{definition}[Neighbor notable vertex fraction]
\lean{FABL.RegularUniqueGame.boundedNeighborNotableVertexFraction}
\label{FABL.RegularUniqueGame.boundedNeighborNotableVertexFraction}
\leanok
\uses{FABL.RegularUniqueGame.BoundedLongCodeTable, FABL.RegularUniqueGame.boundedNeighborNotableCoordinates}
Let $G$ be a regular unique game with alphabet size $q$, let $T$ be a table of functions
from the $q$-dimensional sign cube to $[-1,1]$, and let $\varepsilon\in\mathbb R$. The
neighbor notable vertex fraction is the expected indicator, for a uniformly selected
vertex, that the set of notable coordinates of its averaged neighbor function is
nonempty.
\end{definition}

\begin{theorem}[Excess bounded by notable-vertex fraction]
\lean{FABL.RegularUniqueGame.excess_le_boundedNeighborNotableVertexFraction}
\label{FABL.RegularUniqueGame.excess_le_boundedNeighborNotableVertexFraction}
\leanok
\uses{FABL.HasSmallStableInfluences, FABL.IsDictatorVsNoNotablesTest, FABL.RegularUniqueGame.BoundedLongCodeTable, FABL.RegularUniqueGame.boundedAveragedNeighborFunction, FABL.RegularUniqueGame.boundedLongCodeValue, FABL.RegularUniqueGame.boundedNeighborNotableCoordinates, FABL.RegularUniqueGame.boundedNeighborNotableVertexFraction, FABL.SignPredicate, FABL.SignTester, FABL.SignTester.acceptanceProbability, FABL.SignTester.acceptanceProbability_mem_Icc}
\difficulty{5}
Let $G$ be a regular unique game with alphabet size $q$, let $$ be a
dictator-versus-no-notables test family with predicate set, parameters $,$, error
function $$, and testers $(T_n)_{n>0}$, and let $f$ be a bounded long-code table for
$G$. If $q>0$, $0\le\le1$, and the value of $f$ under the averaged-neighbor tester
$T_q$ is at least $+\mathrm{error}()+$, then $$ is at most the fraction of vertices
whose averaged neighbor function has a coordinate with stable influence greater than
$$.
\end{theorem}

\begin{definition}[Notable neighbor probability]
\lean{FABL.RegularUniqueGame.boundedNotableNeighborProbability}
\label{FABL.RegularUniqueGame.boundedNotableNeighborProbability}
\leanok
\uses{FABL.RegularUniqueGame.BoundedLongCodeTable, FABL.RegularUniqueGame.boundedNotableCoordinates}
Let $G$ be a regular unique game with alphabet size $q$, let $T$ be a table of functions
from the $q$-dimensional sign cube to $[-1,1]$, let $\varepsilon\in\mathbb R$, and fix a
vertex $v$ and coordinate $i\in\{1,\ldots,q\}$. The notable neighbor probability is the
expected indicator, for a uniformly selected neighbor of $v$, that the image of $i$
under the inverse edge permutation is a notable coordinate at the neighboring vertex.
\end{definition}

\begin{theorem}[Notable neighbor probability bound]
\lean{FABL.RegularUniqueGame.half_epsilon_le_boundedNotableNeighborProbability}
\label{FABL.RegularUniqueGame.half_epsilon_le_boundedNotableNeighborProbability}
\leanok
\uses{FABL.RegularUniqueGame.BoundedLongCodeTable, FABL.RegularUniqueGame.boundedAveragedNeighborFunction, FABL.RegularUniqueGame.boundedCubeFunction_stableInfluence_le_one, FABL.RegularUniqueGame.boundedNeighborNotableCoordinates, FABL.RegularUniqueGame.boundedNotableCoordinates, FABL.RegularUniqueGame.boundedNotableNeighborProbability, FABL.RegularUniqueGame.half_lt_thresholdProbability_of_lt_expect, FABL.RegularUniqueGame.stableInfluence_boundedAveragedNeighborFunction_le, FABL.stableInfluence, FABL.stableInfluence_nonneg}
\difficulty{4}
Let $G$ be a regular unique game, let $T$ be a table of functions on the Boolean cube
taking values in $[-1,1]$, and let $0\le\epsilon\le1$. For any vertex $v$ and coordinate
$i$ notable for the averaged neighbor function at $v$ (its stable influence at parameter
$1-\epsilon$ exceeds $\epsilon$), the probability that a uniformly chosen neighbor has
the corresponding permuted coordinate among its notable coordinates (those with stable
influence at parameter $1-\epsilon$ at least $\epsilon/2$) is at least $\epsilon/2$.
\end{theorem}

\begin{theorem}[Neighbor notable coordinate bound]
\lean{FABL.RegularUniqueGame.card_boundedNeighborNotableCoordinates_le}
\label{FABL.RegularUniqueGame.card_boundedNeighborNotableCoordinates_le}
\leanok
\uses{FABL.RegularUniqueGame.BoundedLongCodeTable, FABL.RegularUniqueGame.boundedAveragedNeighborFunction, FABL.RegularUniqueGame.boundedCubeFunction_variance_le_one, FABL.RegularUniqueGame.boundedNeighborNotableCoordinates, FABL.card_stableInfluence_ge_le, FABL.stableInfluence}
\difficulty{4}
Let $G$ be a regular unique game with label set $in q$, and assign to each vertex a
function on the Boolean cube taking values in $[-1,1]$. If $0<ε≤1$, then at every vertex
the number of coordinates whose stable influence in the averaged neighbor function
exceeds $ε$ is at most $1/ε^2$.
\end{theorem}

\begin{theorem}[Notable coordinate bound]
\lean{FABL.RegularUniqueGame.card_boundedNotableCoordinates_le}
\label{FABL.RegularUniqueGame.card_boundedNotableCoordinates_le}
\leanok
\uses{FABL.RegularUniqueGame.BoundedLongCodeTable, FABL.RegularUniqueGame.boundedCubeFunction_variance_le_one, FABL.RegularUniqueGame.boundedNotableCoordinates, FABL.card_stableInfluence_ge_le, FABL.stableInfluence}
\difficulty{3}
Let $G$ be a regular unique game with label set $in q$, and assign to each vertex a
function on the Boolean cube taking values in $[-1,1]$. If $0<ε≤1$, then at every vertex
the number of coordinates whose stable influence is at least $ε/2$ is at most $2/ε^2$.
\end{theorem}

\begin{definition}[Bounded decoding candidates]
\lean{FABL.RegularUniqueGame.boundedDecodingCandidates}
\label{FABL.RegularUniqueGame.boundedDecodingCandidates}
\leanok
\uses{FABL.RegularUniqueGame.BoundedLongCodeTable, FABL.RegularUniqueGame.boundedNeighborNotableCoordinates, FABL.RegularUniqueGame.boundedNotableCoordinates}
Let $G$ be a regular unique game with alphabet size $q$, let $T$ be a table of functions
from the $q$-dimensional sign cube to $[-1,1]$, let $\varepsilon\in\mathbb R$, and fix a
vertex $v$. The decoding candidates at $v$ are the union of the notable coordinates of
$T(v)$ and the notable coordinates of its averaged neighbor function.
\end{definition}

\begin{theorem}[Bound on decoding candidate count]
\lean{FABL.RegularUniqueGame.card_boundedDecodingCandidates_le}
\label{FABL.RegularUniqueGame.card_boundedDecodingCandidates_le}
\leanok
\uses{FABL.RegularUniqueGame.BoundedLongCodeTable, FABL.RegularUniqueGame.boundedDecodingCandidates, FABL.RegularUniqueGame.boundedNeighborNotableCoordinates, FABL.RegularUniqueGame.boundedNotableCoordinates, FABL.RegularUniqueGame.card_boundedNeighborNotableCoordinates_le, FABL.RegularUniqueGame.card_boundedNotableCoordinates_le}
\difficulty{3}
Let $G$ be a regular unique game with label set $in q$, and assign to each vertex a
function on the Boolean cube taking values in $[-1,1]$. If $0<	ext{ε}	ext{≤}1$, then
at every vertex the number of bounded decoding candidates is at most $3/	ext{ε}^2$.
\end{theorem}

\begin{definition}[Bounded decoding edge value]
\lean{FABL.RegularUniqueGame.boundedDecodingEdgeValue}
\label{FABL.RegularUniqueGame.boundedDecodingEdgeValue}
\leanok
\uses{FABL.RegularUniqueGame.BoundedLongCodeTable, FABL.RegularUniqueGame.boundedDecodingCandidates, FABL.RegularUniqueGame.candidateDecodingEdgeValue}
Let $q$ be a positive integer and let $G$ be a regular unique game over $q$ symbols. Let
the table assign to each vertex a function from the Boolean cube $\\{−1,1\\}^{q}$ to
$[-1,1]$. For a real parameter $\\varepsilon$, a default symbol in
$\\{0,\\ldots,q-1\\}$, a vertex, and a neighboring edge, the bounded decoding edge value
is the candidate decoding edge value obtained from the table’s bounded decoding
candidates, using that default symbol.
\end{definition}

\begin{theorem}[Nonnegative decoding edge value]
\lean{FABL.RegularUniqueGame.boundedDecodingEdgeValue_nonneg}
\label{FABL.RegularUniqueGame.boundedDecodingEdgeValue_nonneg}
\leanok
\uses{FABL.RegularUniqueGame.BoundedLongCodeTable, FABL.RegularUniqueGame.boundedDecodingCandidates, FABL.RegularUniqueGame.boundedDecodingEdgeValue, FABL.RegularUniqueGame.candidateDecodingEdgeValue_nonneg}
\difficulty{1}
Let $G$ be a regular unique game, let each vertex be assigned a bounded real-valued
function on its Boolean cube, and choose a real threshold $\varepsilon$ and a default
label. For every vertex and each neighbor of that vertex, the ordinary decoding edge
value determined by the table, threshold, and default label is nonnegative.
\end{theorem}

\begin{definition}[Bounded decoding value]
\lean{FABL.RegularUniqueGame.boundedDecodingValue}
\label{FABL.RegularUniqueGame.boundedDecodingValue}
\leanok
\uses{FABL.RegularUniqueGame.BoundedLongCodeTable, FABL.RegularUniqueGame.boundedDecodingCandidates, FABL.RegularUniqueGame.candidateDecodingValue}
Let $q$ be a positive integer and let $G$ be a regular unique game over $q$ symbols. Let
the table assign to each vertex a function from the Boolean cube $\\{−1,1\\}^{q}$ to
$[-1,1]$. For a real parameter $\\varepsilon$ and a default symbol in
$\\{0,\\ldots,q-1\\}$, the bounded decoding value is the candidate decoding value formed
from the table’s bounded decoding candidates and that default symbol.
\end{definition}

\begin{theorem}[Decoding value bounded by game optimum]
\lean{FABL.RegularUniqueGame.boundedDecodingValue_le_optimum}
\label{FABL.RegularUniqueGame.boundedDecodingValue_le_optimum}
\leanok
\uses{FABL.CSPInstance.optimum, FABL.FiniteUniqueGame.toCSPInstance, FABL.RegularUniqueGame.BoundedLongCodeTable, FABL.RegularUniqueGame.boundedDecodingCandidates, FABL.RegularUniqueGame.boundedDecodingValue, FABL.RegularUniqueGame.candidateDecodingValue_le_optimum, FABL.RegularUniqueGame.toFiniteUniqueGame}
\difficulty{1}
Let $G$ be a regular unique game with a nonempty label set, let each vertex be assigned
a bounded real-valued function on its Boolean cube, and choose a real threshold and a
default label. The expected decoding value obtained from the influential coordinates of
the table and their neighbor averages is at most the optimum value of the constraint
satisfaction instance associated with $G$.
\end{theorem}

\begin{theorem}[Lower bound for bounded edge decoding]
\lean{FABL.RegularUniqueGame.epsilon_pow_four_div_nine_le_boundedDecodingEdgeValue}
\label{FABL.RegularUniqueGame.epsilon_pow_four_div_nine_le_boundedDecodingEdgeValue}
\leanok
\uses{FABL.RegularUniqueGame.BoundedLongCodeTable, FABL.RegularUniqueGame.boundedDecodingCandidates, FABL.RegularUniqueGame.boundedDecodingEdgeValue, FABL.RegularUniqueGame.boundedNeighborNotableCoordinates, FABL.RegularUniqueGame.boundedNotableCoordinates, FABL.RegularUniqueGame.candidateDecodingEdgeValue, FABL.RegularUniqueGame.card_boundedDecodingCandidates_le, FABL.RegularUniqueGame.epsilon_sq_div_three_le_inv_card, FABL.RegularUniqueGame.finiteCandidatePairExpectation_ge}
\difficulty{4}
Let $G$ be a regular unique game with $q$ labels, let $0<\varepsilon\le 1$, and choose a
table of functions on the Boolean cube, each bounded between $-1$ and $1$, a default
label, a vertex, one of its neighbors, and a coordinate. Suppose the coordinate has
stable influence greater than $\varepsilon$ in the neighbor-averaged function at the
vertex, and its image under the inverse constraint permutation has stable influence at
least $\varepsilon/2$ in the table function at the neighboring vertex. Then the bounded
decoding value on this directed edge is at least $\varepsilon^4/9$.
\end{theorem}

\begin{definition}[Bounded decoding vertex value]
\lean{FABL.RegularUniqueGame.boundedDecodingVertexValue}
\label{FABL.RegularUniqueGame.boundedDecodingVertexValue}
\leanok
\uses{FABL.RegularUniqueGame.BoundedLongCodeTable, FABL.RegularUniqueGame.boundedDecodingEdgeValue}
Let $q$ be a positive integer and let $G$ be a regular unique game over $q$ symbols. Let
the table assign to each vertex a function from the Boolean cube $\\{−1,1\\}^{q}$ to
$[-1,1]$. For a real parameter $\\varepsilon$, a default symbol in
$\\{0,\\ldots,q-1\\}$, and a vertex $v$, the bounded decoding vertex value is the
expectation, over a neighbor of $v$, of the bounded decoding edge value on the edge from
$v$ to that neighbor.
\end{definition}

\begin{theorem}[Lower bound for bounded decoding value]
\lean{FABL.RegularUniqueGame.epsilon_pow_five_div_eighteen_le_boundedDecodingVertexValue}
\label{FABL.RegularUniqueGame.epsilon_pow_five_div_eighteen_le_boundedDecodingVertexValue}
\leanok
\uses{FABL.RegularUniqueGame.BoundedLongCodeTable, FABL.RegularUniqueGame.boundedDecodingEdgeValue, FABL.RegularUniqueGame.boundedDecodingEdgeValue_nonneg, FABL.RegularUniqueGame.boundedDecodingVertexValue, FABL.RegularUniqueGame.boundedNeighborNotableCoordinates, FABL.RegularUniqueGame.boundedNotableCoordinates, FABL.RegularUniqueGame.boundedNotableNeighborProbability, FABL.RegularUniqueGame.epsilon_pow_four_div_nine_le_boundedDecodingEdgeValue, FABL.RegularUniqueGame.half_epsilon_le_boundedNotableNeighborProbability}
\difficulty{4}
Let $G$ be a regular unique game with $q$ labels, let $0<\varepsilon\le 1$, and choose a
table of functions on the Boolean cube, each bounded between $-1$ and $1$, a default
label, and a vertex. If at least one coordinate has stable influence greater than
$\varepsilon$ in the neighbor-averaged, constraint-permuted function at that vertex,
then the bounded decoding value at the vertex is at least $\varepsilon^5/18$.
\end{theorem}

\begin{theorem}[Excess and decoding value]
\lean{FABL.RegularUniqueGame.excess_mul_epsilon_pow_five_div_eighteen_le_boundedDecodingValue}
\label{FABL.RegularUniqueGame.excess_mul_epsilon_pow_five_div_eighteen_le_boundedDecodingValue}
\leanok
\uses{FABL.IsDictatorVsNoNotablesTest, FABL.RegularUniqueGame.BoundedLongCodeTable, FABL.RegularUniqueGame.boundedDecodingEdgeValue, FABL.RegularUniqueGame.boundedDecodingEdgeValue_nonneg, FABL.RegularUniqueGame.boundedDecodingValue, FABL.RegularUniqueGame.boundedDecodingVertexValue, FABL.RegularUniqueGame.boundedLongCodeValue, FABL.RegularUniqueGame.boundedNeighborNotableCoordinates, FABL.RegularUniqueGame.boundedNeighborNotableVertexFraction, FABL.RegularUniqueGame.candidateDecodingValue, FABL.RegularUniqueGame.epsilon_pow_five_div_eighteen_le_boundedDecodingVertexValue, FABL.RegularUniqueGame.excess_le_boundedNeighborNotableVertexFraction, FABL.SignPredicate, FABL.SignTester}
\difficulty{5}
Let $G$ be a regular unique game with label set $br^q$ in the sense of labels indexed
by $	ext{Fin}(q)$, and let a tester family satisfy the dictator-versus-no-notables
conditions with parameters $4alpha$ and $4beta$ and error function $4error$. If
$q>0$, $T$ is a table of functions on the $q$-dimensional Boolean cube taking values in
$[-1,1]$, $0<4epsilon\le 1$, and $\alpha+\operatorname{error}(\epsilon)+\tau$ is at
most the tester’s averaged acceptance value on $T$, then $\tau\epsilon^5/18$ is at most
the decoding value of $T$ at threshold $\epsilon$, using label $0$ as the default.
\end{theorem}

\begin{theorem}[Excess and game optimum]
\lean{FABL.RegularUniqueGame.excess_mul_epsilon_pow_five_div_eighteen_le_optimum}
\label{FABL.RegularUniqueGame.excess_mul_epsilon_pow_five_div_eighteen_le_optimum}
\leanok
\uses{FABL.CSPInstance.optimum, FABL.FiniteUniqueGame.toCSPInstance, FABL.IsDictatorVsNoNotablesTest, FABL.RegularUniqueGame.BoundedLongCodeTable, FABL.RegularUniqueGame.boundedDecodingValue_le_optimum, FABL.RegularUniqueGame.boundedLongCodeValue, FABL.RegularUniqueGame.excess_mul_epsilon_pow_five_div_eighteen_le_boundedDecodingValue, FABL.RegularUniqueGame.toFiniteUniqueGame, FABL.SignPredicate, FABL.SignTester}
\difficulty{1}
Let $G$ be a regular unique game with $q>0$ labels, and suppose the label set is
nonempty. Let a tester family satisfy the dictator-versus-no-notables conditions with
parameters $\alpha$ and $\beta$ and error function $\operatorname{error}$. If $T$ is a
table of functions on the $q$-dimensional Boolean cube taking values in $[-1,1]$,
$0<\epsilon\le1$, and the tester’s averaged acceptance value on $T$ is at least
$\alpha+\operatorname{error}(\epsilon)+\tau$, then $\tau\epsilon^5/18$ is at most the
optimum value of the constraint satisfaction instance associated to $G$.
\end{theorem}

\begin{theorem}[Influence of an average of functions]
\lean{FABL.RegularUniqueGame.influence_uniformExpectation_le}
\label{FABL.RegularUniqueGame.influence_uniformExpectation_le}
\leanok
\uses{FABL.RegularUniqueGame.stableInfluence_uniformExpectation_le, FABL.influence, FABL.stableInfluence_one_eq_influence}
\difficulty{1}
Let $G$ be a regular unique game with alphabet size $q$, and let $I$ be a finite
nonempty set of real-valued functions on the $q$-dimensional Boolean cube. For every
coordinate $i$, the influence of coordinate $i$ in their pointwise average is at most
the average of their influences in that coordinate.
\end{theorem}

\begin{theorem}[Influence under input permutation]
\lean{FABL.RegularUniqueGame.influence_comp_permuteInput}
\label{FABL.RegularUniqueGame.influence_comp_permuteInput}
\leanok
\uses{FABL.RegularUniqueGame.stableInfluence_comp_permuteInput, FABL.influence, FABL.permuteInput, FABL.stableInfluence_one_eq_influence}
\difficulty{1}
Let $G$ be a regular unique game with alphabet size $q$, let $\pi$ be a permutation of
the $q$ coordinates, and let $f$ be a real-valued function on the Boolean cube of
dimension $q$. For every coordinate $i$, the influence of coordinate $i$ in the function
$x\mapsto f(x\circ\pi)$ equals the influence of coordinate $\pi^{-1}(i)$ in $f$.
\end{theorem}

\begin{theorem}[Influence bound for bounded functions]
\lean{FABL.RegularUniqueGame.boundedCubeFunction_influence_le_one}
\label{FABL.RegularUniqueGame.boundedCubeFunction_influence_le_one}
\leanok
\uses{FABL.BoundedCubeFunction, FABL.RegularUniqueGame.boundedCubeFunction_stableInfluence_le_one, FABL.influence, FABL.stableInfluence_one_eq_influence}
\difficulty{1}
Let $f:\{\pm1\}^q\to\mathbb R$ take values in $[-1,1]$. For every coordinate $i\in[q]$,
its influence, defined as the uniform average of the square of the discrete derivative
in coordinate $i$, is at most $1$.
\end{theorem}

\begin{definition}[Bounded neighbor influential coordinates]
\lean{FABL.RegularUniqueGame.boundedNeighborInfluentialCoordinates}
\label{FABL.RegularUniqueGame.boundedNeighborInfluentialCoordinates}
\leanok
\uses{FABL.RegularUniqueGame.BoundedLongCodeTable, FABL.RegularUniqueGame.boundedAveragedNeighborFunction, FABL.influence}
Let $q$ be a positive integer and let $G$ be a regular unique game over $q$ symbols. Let
the table assign to each vertex a function from the Boolean cube $\\{−1,1\\}^{q}$ to
$[-1,1]$. For a real threshold $\\varepsilon$ and a vertex $v$, the bounded neighbor
influential coordinates are the coordinates $i\\in\\{0,\\ldots,q-1\\}$ for which the
influence of coordinate $i$ in the averaged neighbor function at $v$ is strictly greater
than $\\varepsilon$.
\end{definition}

\begin{definition}[Bounded influential coordinates]
\lean{FABL.RegularUniqueGame.boundedInfluentialCoordinates}
\label{FABL.RegularUniqueGame.boundedInfluentialCoordinates}
\leanok
\uses{FABL.RegularUniqueGame.BoundedLongCodeTable, FABL.influence}
Let $q$ be a positive integer and let $G$ be a regular unique game over $q$ symbols. Let
the table assign to each vertex a function from the Boolean cube $\\{−1,1\\}^{q}$ to
$[-1,1]$. For a real parameter $\\varepsilon$ and a vertex $v$, the bounded influential
coordinates are those coordinates $i\\in\\{0,\\ldots,q-1\\}$ whose influence in the
function assigned to $v$ is at least $\\varepsilon/2$.
\end{definition}

\begin{theorem}[Influence of the averaged neighbor function]
\lean{FABL.RegularUniqueGame.influence_boundedAveragedNeighborFunction_le}
\label{FABL.RegularUniqueGame.influence_boundedAveragedNeighborFunction_le}
\leanok
\uses{FABL.RegularUniqueGame.BoundedLongCodeTable, FABL.RegularUniqueGame.boundedAveragedNeighborFunction, FABL.RegularUniqueGame.boundedTwistedLongCode, FABL.RegularUniqueGame.influence_comp_permuteInput, FABL.RegularUniqueGame.influence_uniformExpectation_le, FABL.influence}
\difficulty{3}
Let $G$ be a regular unique game with alphabet size $q$, let $T$ assign to each vertex a
function on the $q$-dimensional Boolean cube taking values in $[-1,1]$, and fix a vertex
$v$ and coordinate $i\in\{1,\ldots,q\}$. The influence of coordinate $i$ in the function
obtained by averaging the twisted neighbor functions at $v$ is at most the average, over
neighbors $w$ of $v$, of the influence of coordinate $\pi_{v,w}^{-1}(i)$ in $T(w)$,
where $\pi_{v,w}$ is the edge permutation.
\end{theorem}

\begin{definition}[Bounded neighbor influential vertex fraction]
\lean{FABL.RegularUniqueGame.boundedNeighborInfluentialVertexFraction}
\label{FABL.RegularUniqueGame.boundedNeighborInfluentialVertexFraction}
\leanok
\uses{FABL.RegularUniqueGame.BoundedLongCodeTable, FABL.RegularUniqueGame.boundedNeighborInfluentialCoordinates}
Let $q$ be a positive integer and let $G$ be a regular unique game over $q$ symbols. Let
the table assign to each vertex a function from the Boolean cube $\\{−1,1\\}^{q}$ to
$[-1,1]$, and let $\\varepsilon$ be real. The bounded neighbor influential vertex
fraction is the expectation over vertices of the indicator that the vertex has at least
one bounded neighbor influential coordinate.
\end{definition}

\begin{theorem}[Excess bounded by influential-vertex fraction]
\lean{FABL.RegularUniqueGame.excess_le_boundedNeighborInfluentialVertexFraction}
\label{FABL.RegularUniqueGame.excess_le_boundedNeighborInfluentialVertexFraction}
\leanok
\uses{FABL.HasSmallInfluences, FABL.IsDictatorVsNoInfluentialsTest, FABL.RegularUniqueGame.BoundedLongCodeTable, FABL.RegularUniqueGame.boundedAveragedNeighborFunction, FABL.RegularUniqueGame.boundedLongCodeValue, FABL.RegularUniqueGame.boundedNeighborInfluentialCoordinates, FABL.RegularUniqueGame.boundedNeighborInfluentialVertexFraction, FABL.SignPredicate, FABL.SignTester, FABL.SignTester.acceptanceProbability, FABL.SignTester.acceptanceProbability_mem_Icc, FABL.hasSmallInfluences_iff}
\difficulty{5}
Let $G$ be a regular unique game with alphabet size $q$, let $$ be a
dictator-versus-no-influentials test family with predicate set, parameters $,$, error
function $$, and testers $(T_n)_{n>0}$, and let $f$ be a table of functions from the
$q$-dimensional Boolean cube to $[-1,1]$, indexed by the vertices of $G$. If $q>0$,
$0\le\le1$, and the value of $f$ under the averaged-neighbor tester $T_q$ is at least
$+\mathrm{error}()+$, then $$ is at most the fraction of vertices whose averaged
neighbor function has some coordinate with influence greater than $$.
\end{theorem}

\begin{definition}[Bounded influential neighbor probability]
\lean{FABL.RegularUniqueGame.boundedInfluentialNeighborProbability}
\label{FABL.RegularUniqueGame.boundedInfluentialNeighborProbability}
\leanok
\uses{FABL.RegularUniqueGame.BoundedLongCodeTable, FABL.RegularUniqueGame.boundedInfluentialCoordinates}
Let $q$ be a positive integer and let $G$ be a regular unique game over $q$ symbols. Let
the table assign to each vertex a function from the Boolean cube $\\{−1,1\\}^{q}$ to
$[-1,1]$. For a real parameter $\\varepsilon$, a vertex $v$, and a coordinate
$i\\in\\{0,\\ldots,q-1\\}$, the bounded influential neighbor probability is the
expectation over neighbors $w$ of $v$ of the indicator that the coordinate obtained by
applying the inverse edge permutation to $i$ is a bounded influential coordinate at $w$.
\end{definition}

\begin{theorem}[Influential neighbor probability bound]
\lean{FABL.RegularUniqueGame.half_epsilon_le_boundedInfluentialNeighborProbability}
\label{FABL.RegularUniqueGame.half_epsilon_le_boundedInfluentialNeighborProbability}
\leanok
\uses{FABL.RegularUniqueGame.BoundedLongCodeTable, FABL.RegularUniqueGame.boundedAveragedNeighborFunction, FABL.RegularUniqueGame.boundedCubeFunction_influence_le_one, FABL.RegularUniqueGame.boundedInfluentialCoordinates, FABL.RegularUniqueGame.boundedInfluentialNeighborProbability, FABL.RegularUniqueGame.boundedNeighborInfluentialCoordinates, FABL.RegularUniqueGame.half_lt_thresholdProbability_of_lt_expect, FABL.RegularUniqueGame.influence_boundedAveragedNeighborFunction_le, FABL.influence, FABL.influence_nonneg}
\difficulty{4}
Let $G$ be a regular unique game, let $T$ be a table of functions on the Boolean cube
taking values in $[-1,1]$, and let $0\le\epsilon\le1$. For any vertex $v$ and coordinate
$i$ that is notable for the averaged neighbor function at $v$ (its stable influence at
parameter $1-\epsilon$ exceeds $\epsilon$), the probability that a uniformly chosen
neighbor has the corresponding permuted coordinate among its influential coordinates
(those with influence at least $\epsilon/2$) is at least $\epsilon/2$.
\end{theorem}

\begin{definition}[Bounded ordinary decoding candidates]
\lean{FABL.RegularUniqueGame.boundedOrdinaryDecodingCandidates}
\label{FABL.RegularUniqueGame.boundedOrdinaryDecodingCandidates}
\leanok
\uses{FABL.RegularUniqueGame.BoundedLongCodeTable, FABL.RegularUniqueGame.boundedInfluentialCoordinates, FABL.RegularUniqueGame.boundedNeighborInfluentialCoordinates}
Let $q$ be a positive integer and let $G$ be a regular unique game over $q$ symbols. Let
the table assign to each vertex a function from the Boolean cube $\\{−1,1\\}^{q}$ to
$[-1,1]$. For a real parameter $\\varepsilon$ and a vertex $v$, the bounded ordinary
decoding candidates are the union of the bounded influential coordinates at $v$ and the
bounded neighbor influential coordinates at $v$.
\end{definition}

\begin{theorem}[Ordinary decoding candidate count]
\lean{FABL.RegularUniqueGame.card_boundedOrdinaryDecodingCandidates_le}
\label{FABL.RegularUniqueGame.card_boundedOrdinaryDecodingCandidates_le}
\leanok
\uses{FABL.RegularUniqueGame.BoundedLongCodeTable, FABL.RegularUniqueGame.boundedInfluentialCoordinates, FABL.RegularUniqueGame.boundedNeighborInfluentialCoordinates, FABL.RegularUniqueGame.boundedOrdinaryDecodingCandidates}
\difficulty{4}
Let $G$ be a regular unique game with label set $in q$, and assign to each vertex a
function on the Boolean cube taking values in $[-1,1]$. At every vertex, the number of
ordinary decoding candidates is at most $2q$.
\end{theorem}

\begin{theorem}[Reciprocal Candidate-Count Bound]
\lean{FABL.RegularUniqueGame.one_div_two_mul_q_le_inv_card}
\label{FABL.RegularUniqueGame.one_div_two_mul_q_le_inv_card}
\leanok
\difficulty{3}
Let $q$ be a natural number and let $C$ be a nonempty subset of the label set
$	ext{Fin}(q)$. If $|C|\le 2q$, then $\frac{1}{2q}\le \frac{1}{|C|}$.
\end{theorem}

\begin{definition}[Bounded decoding edge value]
\lean{FABL.RegularUniqueGame.boundedOrdinaryDecodingEdgeValue}
\label{FABL.RegularUniqueGame.boundedOrdinaryDecodingEdgeValue}
\leanok
\uses{FABL.RegularUniqueGame.BoundedLongCodeTable, FABL.RegularUniqueGame.boundedOrdinaryDecodingCandidates, FABL.RegularUniqueGame.candidateDecodingEdgeValue}
Let $G$ be a regular unique game with alphabet size $q$, let each vertex be assigned a
real-valued function on the sign cube $\\{-1,1\\}^q$ taking values in $[-1,1]$, and let
$\\varepsilon\\in\\mathbb R$, a default symbol in $\\mathrm{Fin}(q)$, a vertex $v$, and
a neighbor of $v$ be given. The bounded decoding edge value is the candidate-decoding
edge value at that vertex and neighbor, using the candidates produced from this table
and $\\varepsilon$ and the specified default symbol.
\end{definition}

\begin{theorem}[Nonnegative ordinary decoding edge value]
\lean{FABL.RegularUniqueGame.boundedOrdinaryDecodingEdgeValue_nonneg}
\label{FABL.RegularUniqueGame.boundedOrdinaryDecodingEdgeValue_nonneg}
\leanok
\uses{FABL.RegularUniqueGame.BoundedLongCodeTable, FABL.RegularUniqueGame.boundedOrdinaryDecodingCandidates, FABL.RegularUniqueGame.boundedOrdinaryDecodingEdgeValue, FABL.RegularUniqueGame.candidateDecodingEdgeValue_nonneg}
\difficulty{1}
Let $G$ be a regular unique game, let each vertex be assigned a bounded real-valued
function on its Boolean cube, and choose a real threshold $\varepsilon$ and a default
label. For every vertex and each neighbor of that vertex, the ordinary decoding edge
value determined by the table, threshold, and default label is nonnegative.
\end{theorem}

\begin{theorem}[Lower bound for decoding an edge]
\lean{FABL.RegularUniqueGame.one_div_four_mul_q_sq_le_boundedOrdinaryDecodingEdgeValue}
\label{FABL.RegularUniqueGame.one_div_four_mul_q_sq_le_boundedOrdinaryDecodingEdgeValue}
\leanok
\uses{FABL.RegularUniqueGame.BoundedLongCodeTable, FABL.RegularUniqueGame.boundedInfluentialCoordinates, FABL.RegularUniqueGame.boundedNeighborInfluentialCoordinates, FABL.RegularUniqueGame.boundedOrdinaryDecodingCandidates, FABL.RegularUniqueGame.boundedOrdinaryDecodingEdgeValue, FABL.RegularUniqueGame.candidateDecodingEdgeValue, FABL.RegularUniqueGame.card_boundedOrdinaryDecodingCandidates_le, FABL.RegularUniqueGame.finiteCandidatePairExpectation_ge, FABL.RegularUniqueGame.one_div_two_mul_q_le_inv_card}
\difficulty{5}
Let $G$ be a regular unique game with $q>0$, let $F$ be a table of real-valued functions
on the Boolean cube, let $\varepsilon\in\mathbb R$, and choose a default alphabet
coordinate, a vertex $v$, a neighbour $w$ of $v$, and a coordinate $i$. If $i$ is
influential for the averaged neighbour function at $v$ according to the $\varepsilon$
threshold, and the inverse image of $i$ under the constraint permutation is influential
for the function at $w$ according to the $\varepsilon/2$ threshold, then the ordinary
candidate-decoding edge value is at least $1/(4q^2)$.
\end{theorem}

\begin{definition}[Bounded decoding vertex value]
\lean{FABL.RegularUniqueGame.boundedOrdinaryDecodingVertexValue}
\label{FABL.RegularUniqueGame.boundedOrdinaryDecodingVertexValue}
\leanok
\uses{FABL.RegularUniqueGame.BoundedLongCodeTable, FABL.RegularUniqueGame.boundedOrdinaryDecodingEdgeValue}
Let $G$ be a regular unique game with alphabet size $q$, let each vertex be assigned a
real-valued function on the sign cube $\\{-1,1\\}^q$ taking values in $[-1,1]$, and let
$\\varepsilon\\in\\mathbb R$, a default symbol in $\\mathrm{Fin}(q)$, and a vertex $v$
be given. The bounded decoding vertex value is the expectation, over a neighbor of $v$,
of the bounded decoding edge value for that neighbor.
\end{definition}

\begin{theorem}[Lower bound for ordinary decoding value]
\lean{FABL.RegularUniqueGame.epsilon_div_eight_mul_q_sq_le_boundedOrdinaryDecodingVertexValue}
\label{FABL.RegularUniqueGame.epsilon_div_eight_mul_q_sq_le_boundedOrdinaryDecodingVertexValue}
\leanok
\uses{FABL.RegularUniqueGame.BoundedLongCodeTable, FABL.RegularUniqueGame.boundedInfluentialCoordinates, FABL.RegularUniqueGame.boundedInfluentialNeighborProbability, FABL.RegularUniqueGame.boundedNeighborInfluentialCoordinates, FABL.RegularUniqueGame.boundedOrdinaryDecodingEdgeValue, FABL.RegularUniqueGame.boundedOrdinaryDecodingEdgeValue_nonneg, FABL.RegularUniqueGame.boundedOrdinaryDecodingVertexValue, FABL.RegularUniqueGame.half_epsilon_le_boundedInfluentialNeighborProbability, FABL.RegularUniqueGame.one_div_four_mul_q_sq_le_boundedOrdinaryDecodingEdgeValue}
\difficulty{5}
Let $G$ be a regular unique game with $q>0$ labels, let $0<\varepsilon\le 1$, and choose
a default label and a vertex. For a table of real-valued functions on the Boolean cube,
each taking values in $[-1,1]$, suppose that at this vertex at least one coordinate has
stable influence greater than $\varepsilon$ in the neighbor-averaged,
constraint-permuted function. Then the ordinary decoding value at the vertex is at least
$\varepsilon/(8q^2)$.
\end{theorem}

\begin{definition}[Bounded ordinary decoding value]
\lean{FABL.RegularUniqueGame.boundedOrdinaryDecodingValue}
\label{FABL.RegularUniqueGame.boundedOrdinaryDecodingValue}
\leanok
\uses{FABL.RegularUniqueGame.BoundedLongCodeTable, FABL.RegularUniqueGame.boundedOrdinaryDecodingCandidates, FABL.RegularUniqueGame.candidateDecodingValue}
Let $G$ be a regular unique game with alphabet size $q$, let each vertex be assigned a
real-valued function on the sign cube $\\{-1,1\\}^q$ taking values in $[-1,1]$, and let
$\\varepsilon\\in\\mathbb R$ and a default symbol in $\\mathrm{Fin}(q)$ be given. The
bounded ordinary decoding value is the candidate-decoding value computed from the
candidates produced by this table and $\\varepsilon$, using the specified default
symbol.
\end{definition}

\begin{theorem}[Ordinary decoding value bound]
\lean{FABL.RegularUniqueGame.boundedOrdinaryDecodingValue_le_optimum}
\label{FABL.RegularUniqueGame.boundedOrdinaryDecodingValue_le_optimum}
\leanok
\uses{FABL.CSPInstance.optimum, FABL.FiniteUniqueGame.toCSPInstance, FABL.RegularUniqueGame.BoundedLongCodeTable, FABL.RegularUniqueGame.boundedOrdinaryDecodingCandidates, FABL.RegularUniqueGame.boundedOrdinaryDecodingValue, FABL.RegularUniqueGame.candidateDecodingValue_le_optimum, FABL.RegularUniqueGame.toFiniteUniqueGame}
\difficulty{1}
Let $G$ be a regular unique game with label set $in q$, where $q>0$, and let each
vertex be assigned a function on the Boolean cube taking values in $[-1,1]$. For any
threshold $	ext{ε}$ and default label, the expected edge agreement obtained by
choosing uniformly from each vertex’s ordinary decoding candidates (using the default
when a candidate set is empty) is at most the optimum value of the associated constraint
satisfaction instance.
\end{theorem}

\begin{theorem}[Ordinary decoding value lower bound]
\lean{FABL.RegularUniqueGame.excess_mul_epsilon_div_eight_mul_q_sq_le_boundedOrdinaryDecodingValue}
\label{FABL.RegularUniqueGame.excess_mul_epsilon_div_eight_mul_q_sq_le_boundedOrdinaryDecodingValue}
\leanok
\uses{FABL.IsDictatorVsNoInfluentialsTest, FABL.RegularUniqueGame.BoundedLongCodeTable, FABL.RegularUniqueGame.boundedLongCodeValue, FABL.RegularUniqueGame.boundedNeighborInfluentialCoordinates, FABL.RegularUniqueGame.boundedNeighborInfluentialVertexFraction, FABL.RegularUniqueGame.boundedOrdinaryDecodingEdgeValue, FABL.RegularUniqueGame.boundedOrdinaryDecodingEdgeValue_nonneg, FABL.RegularUniqueGame.boundedOrdinaryDecodingValue, FABL.RegularUniqueGame.boundedOrdinaryDecodingVertexValue, FABL.RegularUniqueGame.candidateDecodingValue, FABL.RegularUniqueGame.epsilon_div_eight_mul_q_sq_le_boundedOrdinaryDecodingVertexValue, FABL.RegularUniqueGame.excess_le_boundedNeighborInfluentialVertexFraction, FABL.SignPredicate, FABL.SignTester}
\difficulty{5}
Let $G$ be a regular unique game with positive alphabet size $q$, and let a
dictator-versus-no-influentials test family have predicate set $\Psi$, parameters
$\alpha,\beta$, error function $\mathrm{error}$, and testers $T_n$ for each positive
integer $n$. For any bounded long-code table $f$ on $G$, if $0<\varepsilon\le 1$ and
$\alpha+\mathrm{error}(\varepsilon)+\tau$ is at most the value of $f$ under the
averaged-neighbor tester $T_q$, then the ordinary decoding value of $f$ at threshold
$\varepsilon$, with default label $0$, is at least $\tau\varepsilon/(8q^2)$.
\end{theorem}

\begin{theorem}[Optimum bound from tester excess]
\lean{FABL.RegularUniqueGame.excess_mul_epsilon_div_eight_mul_q_sq_le_optimum_of_noInfluentials}
\label{FABL.RegularUniqueGame.excess_mul_epsilon_div_eight_mul_q_sq_le_optimum_of_noInfluentials}
\leanok
\uses{FABL.CSPInstance.optimum, FABL.FiniteUniqueGame.toCSPInstance, FABL.IsDictatorVsNoInfluentialsTest, FABL.RegularUniqueGame.BoundedLongCodeTable, FABL.RegularUniqueGame.boundedLongCodeValue, FABL.RegularUniqueGame.boundedOrdinaryDecodingValue_le_optimum, FABL.RegularUniqueGame.excess_mul_epsilon_div_eight_mul_q_sq_le_boundedOrdinaryDecodingValue, FABL.RegularUniqueGame.toFiniteUniqueGame, FABL.SignPredicate, FABL.SignTester}
\difficulty{1}
Let $q>0$, let $G$ be a regular unique game with alphabet of size $q$, and let its label
set be nonempty. Suppose a finite set of sign predicates and real parameters
$0<\alpha<\beta\le 1$ define a dictator-versus-no-influentials test family for every
positive input dimension: its error function maps $[0,1]$ into $[0,1]$ and tends to $0$
as its argument decreases to $0$, every dictator is accepted with probability at least
$\beta$, and every function whose influences are all at most $\varepsilon$ is accepted
with probability at most $\alpha+\mathrm{error}(\varepsilon)$. For any bounded long-code
table $f$ for $G$, any $0<\varepsilon\le 1$, and any $\tau\in\mathbb R$, if the
averaged-neighbor test value of $f$ is at least
$\alpha+\mathrm{error}(\varepsilon)+\tau$, then the optimum value of the constraint
satisfaction instance associated with $G$ is at least $\tau\varepsilon/(8q^2)$.
\end{theorem}

\begin{definition}[Nonempty raw constraint instance]
\lean{FABL.RawCSPInstance}
\statementsource{Analysis-of-Boolean-Functions-by-Ryan-ODonnell}{
  pdf-ffbe9c7b3306-p184-b004,
  pdf-ffbe9c7b3306-p185-b004
}
\label{FABL.RawCSPInstance}
\leanok
\uses{FABL.CSPConstraint, FABL.CSPTemplate}
Let $\Psi$ be a constraint template and $V$ a variable set. A raw constraint
satisfaction instance consists of a multiset of constraints over $\Psi$ and $V$,
required to be nonempty.
\end{definition}

\begin{definition}[Participating variables]
\lean{FABL.RawCSPInstance.ParticipatingVariable}
\statementsource{Analysis-of-Boolean-Functions-by-Ryan-ODonnell}{
  pdf-ffbe9c7b3306-p185-b004,
  pdf-ffbe9c7b3306-p185-b005
}
\label{FABL.RawCSPInstance.ParticipatingVariable}
\leanok
Let $P$ be a raw constraint satisfaction instance over a constraint template and a
variable set $V$. Its participating variables are the variables that occur in the scope
of at least one constraint in $P$.
\end{definition}

\begin{definition}[Participating constraint]
\lean{FABL.RawCSPInstance.participatingConstraint}
\statementsource{Analysis-of-Boolean-Functions-by-Ryan-ODonnell}{pdf-ffbe9c7b3306-p185-b004}
\label{FABL.RawCSPInstance.participatingConstraint}
\leanok
\uses{FABL.CSPConstraint, FABL.RawCSPInstance.ParticipatingVariable}
Given a raw constraint satisfaction instance $P$ and a constraint $C$ in its constraint
multiset, regard each variable in the scope of $C$ as the corresponding participating
variable of $P$, retaining the predicate and the ordered scope of $C$.
\end{definition}

\begin{definition}[Removing isolated variables]
\lean{FABL.RawCSPInstance.removeIsolatedVariables}
\statementsource{Analysis-of-Boolean-Functions-by-Ryan-ODonnell}{
  pdf-ffbe9c7b3306-p185-b004,
  pdf-ffbe9c7b3306-p185-b005
}
\label{FABL.RawCSPInstance.removeIsolatedVariables}
\leanok
\uses{FABL.CSPInstance, FABL.RawCSPInstance.ParticipatingVariable, FABL.RawCSPInstance.participatingConstraint}
For a raw constraint satisfaction instance $P$, form an instance whose variables are
precisely those occurring in at least one constraint scope. Its constraint multiset is
obtained by replacing each constraint of $P$ with the corresponding constraint over
these participating variables, preserving multiplicities; the resulting instance has at
least one constraint, and every variable occurs in a constraint scope.
\end{definition}

\begin{definition}[Value of a raw constraint instance]
\lean{FABL.RawCSPInstance.value}
\statementsource{Analysis-of-Boolean-Functions-by-Ryan-ODonnell}{
  pdf-ffbe9c7b3306-p185-b006,
  pdf-ffbe9c7b3306-p186-b001
}
\label{FABL.RawCSPInstance.value}
\leanok
\uses{FABL.CSPConstraint.satisfiedIndicator}
Given a raw constraint satisfaction instance $P$ and an assignment of values to its
variables, the value of the assignment is the number of constraints it satisfies,
counted with multiplicity, divided by the total number of constraints, also counted with
multiplicity.
\end{definition}

\begin{definition}[Restriction of an assignment]
\lean{FABL.RawCSPInstance.restrictAssignment}
\label{FABL.RawCSPInstance.restrictAssignment}
\leanok
\uses{FABL.RawCSPInstance.ParticipatingVariable}
Let $P$ be a constraint instance with variable set $V$ and domain $D$. Its participating
variables are those that occur in the scope of at least one constraint of $P$. Given an
assignment $a:V\\to D$, its restriction to the participating variables sends each such
variable to its value under $a$.
\end{definition}

\begin{theorem}[Satisfaction under restriction to participating variables]
\lean{FABL.RawCSPInstance.participatingConstraint_satisfiedIndicator}
\label{FABL.RawCSPInstance.participatingConstraint_satisfiedIndicator}
\leanok
\uses{FABL.CSPConstraint, FABL.CSPConstraint.satisfiedIndicator, FABL.RawCSPInstance.participatingConstraint, FABL.RawCSPInstance.restrictAssignment}
\difficulty{1}
Let $P$ be a raw constraint satisfaction instance over a domain $D$, with a finite
family of allowed predicates. For any constraint of $P$ and any assignment to all
variables, the constraint's satisfaction indicator is the same when the constraint is
viewed on the participating variables and the assignment is restricted to those
variables.
\end{theorem}

\begin{theorem}[Value is preserved by restriction]
\lean{FABL.RawCSPInstance.value_restrictAssignment}
\label{FABL.RawCSPInstance.value_restrictAssignment}
\leanok
\uses{FABL.CSPConstraint.satisfiedIndicator, FABL.CSPInstance.satisfiedCount, FABL.CSPInstance.value, FABL.RawCSPInstance.participatingConstraint_satisfiedIndicator, FABL.RawCSPInstance.removeIsolatedVariables, FABL.RawCSPInstance.restrictAssignment, FABL.RawCSPInstance.value}
\difficulty{3}
Let $P$ be a raw constraint satisfaction instance over a domain $D$, with a finite
family of allowed predicates. For every assignment to all variables, the value of its
restriction to the participating variables in the instance with isolated variables
removed equals the value of the original assignment in $P$.
\end{theorem}

\begin{definition}[Extension of an assignment]
\lean{FABL.RawCSPInstance.extendAssignment}
\label{FABL.RawCSPInstance.extendAssignment}
\leanok
\uses{FABL.RawCSPInstance.ParticipatingVariable}
Let $P$ be a constraint instance with variable set $V$, domain $D$, and at least one
element in $D$. Given an assignment from the participating variables of $P$ to $D$,
extend it to all of $V$ by retaining those assigned values on participating variables
and choosing one fixed element of $D$ for every nonparticipating variable.
\end{definition}

\begin{theorem}[Restriction of an extended assignment]
\lean{FABL.RawCSPInstance.restrictAssignment_extendAssignment}
\label{FABL.RawCSPInstance.restrictAssignment_extendAssignment}
\leanok
\uses{FABL.RawCSPInstance.ParticipatingVariable, FABL.RawCSPInstance.extendAssignment, FABL.RawCSPInstance.restrictAssignment}
\difficulty{4}
Let $P$ be a raw constraint satisfaction instance over a nonempty domain $D$. For every
assignment on the variables that occur in constraint scopes, extending it to all
variables and then restricting back to those participating variables recovers the
original assignment.
\end{theorem}

\begin{theorem}[Extending assignments without changing value]
\lean{FABL.RawCSPInstance.exists_raw_assignment_same_value}
\label{FABL.RawCSPInstance.exists_raw_assignment_same_value}
\leanok
\uses{FABL.CSPInstance.value, FABL.RawCSPInstance.ParticipatingVariable, FABL.RawCSPInstance.extendAssignment, FABL.RawCSPInstance.removeIsolatedVariables, FABL.RawCSPInstance.restrictAssignment_extendAssignment, FABL.RawCSPInstance.value, FABL.RawCSPInstance.value_restrictAssignment}
\difficulty{3}
Let $P$ be a raw constraint satisfaction instance over a nonempty domain $D$, with a
finite family of allowed predicates, and let an assignment be given on the variables
that occur in constraint scopes. There is an assignment to all variables whose value in
$P$ equals the value of the given assignment in the instance obtained by removing
isolated variables.
\end{theorem}

\begin{definition}[Optimal assignment]
\lean{FABL.RawCSPInstance.optimalAssignment}
\label{FABL.RawCSPInstance.optimalAssignment}
\leanok
\uses{FABL.RawCSPInstance.value}
Let $P$ be a constraint instance with finite, nonempty domain $D$ and finite variable
set $V$. An optimal assignment is an assignment $V\\to D$ chosen to maximize the value
of $P$.
\end{definition}

\begin{definition}[Optimum of a raw constraint instance]
\lean{FABL.RawCSPInstance.optimum}
\statementsource{Analysis-of-Boolean-Functions-by-Ryan-ODonnell}{pdf-ffbe9c7b3306-p186-b001}
\label{FABL.RawCSPInstance.optimum}
\leanok
\uses{FABL.RawCSPInstance.optimalAssignment, FABL.RawCSPInstance.value}
For a raw constraint satisfaction instance $P$ with finite, nonempty domain and a finite
set of variables, its optimum is the value of an optimal assignment to its variables.
\end{definition}

\begin{theorem}[Every assignment is bounded by the optimum]
\lean{FABL.RawCSPInstance.value_le_optimum}
\label{FABL.RawCSPInstance.value_le_optimum}
\leanok
\uses{FABL.RawCSPInstance.optimum, FABL.RawCSPInstance.value}
\difficulty{1}
Let $P$ be a raw constraint satisfaction instance over a nonempty finite domain, with a
finite family of allowed predicates, and suppose its variable set is finite. The value
of every assignment to its variables is at most the optimum value of $P$.
\end{theorem}

\begin{theorem}[Optimum after removing isolated variables]
\lean{FABL.RawCSPInstance.optimum_removeIsolatedVariables}
\label{FABL.RawCSPInstance.optimum_removeIsolatedVariables}
\leanok
\uses{FABL.CSPInstance.optimalAssignment, FABL.CSPInstance.optimum, FABL.CSPInstance.value_le_optimum, FABL.RawCSPInstance.exists_raw_assignment_same_value, FABL.RawCSPInstance.optimalAssignment, FABL.RawCSPInstance.optimum, FABL.RawCSPInstance.removeIsolatedVariables, FABL.RawCSPInstance.restrictAssignment, FABL.RawCSPInstance.value_le_optimum, FABL.RawCSPInstance.value_restrictAssignment}
\difficulty{4}
Let $P$ be a raw constraint satisfaction instance over a nonempty finite domain, with a
finite family of allowed predicates, and suppose its variable set is finite. The optimum
value is unchanged when isolated variables are removed.
\end{theorem}

\begin{definition}[Uniform refinement of a weighted experiment]
\lean{FABL.WeightedConstraintExperiment.UniformRefinement}
\statementsource{Analysis-of-Boolean-Functions-by-Ryan-ODonnell}{pdf-ffbe9c7b3306-p202-b003}
\label{FABL.WeightedConstraintExperiment.UniformRefinement}
\leanok
\uses{FABL.uniformPMF}
A uniform refinement of a weighted constraint experiment $E$ consists of a finite,
nonempty set of occurrences, an outcome in $E$ for each occurrence, and the requirement
that mapping the uniform probability distribution on the occurrences to their outcomes
gives exactly the law of $E$.
\end{definition}

\begin{definition}[Padded clone occurrences]
\lean{FABL.WeightedConstraintExperiment.UniformRefinement.PaddedCloneOccurrence}
\statementsource{Analysis-of-Boolean-Functions-by-Ryan-ODonnell}{pdf-ffbe9c7b3306-p202-b003}
\label{FABL.WeightedConstraintExperiment.UniformRefinement.PaddedCloneOccurrence}
\leanok
For a uniform refinement $R$ and natural numbers $r$ and $M$, a padded clone occurrence
is a pair consisting of an occurrence in $R$ and a function from the first $r$ positions
to the first $M$ positions.
\end{definition}

\begin{definition}[Padded clone scope]
\lean{FABL.WeightedConstraintExperiment.UniformRefinement.paddedCloneScope}
\statementsource{Analysis-of-Boolean-Functions-by-Ryan-ODonnell}{pdf-ffbe9c7b3306-p202-b003}
\label{FABL.WeightedConstraintExperiment.UniformRefinement.paddedCloneScope}
\leanok
\uses{FABL.WeightedConstraintExperiment.UniformRefinement.PaddedCloneOccurrence}
Let $E$ be a weighted constraint experiment with uniform refinement $R$, and let $r,M$
be natural numbers. Suppose every outcome of $E$ has predicate arity at most $r$. For a
padded clone occurrence, the scope sends each position of its sampled outcome's
predicate to the pair consisting of that outcome's variable at the position and the
clone index assigned to that position.
\end{definition}

\begin{definition}[Padded clone payoff]
\lean{FABL.WeightedConstraintExperiment.UniformRefinement.paddedClonePayoff}
\statementsource{Analysis-of-Boolean-Functions-by-Ryan-ODonnell}{pdf-ffbe9c7b3306-p202-b003}
\label{FABL.WeightedConstraintExperiment.UniformRefinement.paddedClonePayoff}
\leanok
\uses{FABL.WeightedConstraintExperiment.UniformRefinement.PaddedCloneOccurrence, FABL.WeightedConstraintExperiment.clonePayoff}
Let $E$ be a weighted constraint experiment with uniform refinement $R$, and let $r,M$
be natural numbers. Suppose every outcome of $E$ has predicate arity at most $r$. For an
assignment of values to pairs of a variable and a clone index, the padded clone payoff
at a sample is the payoff of its sampled outcome, using the clone indices assigned to
the positions of that outcome.
\end{definition}

\begin{definition}[Padded clone value]
\lean{FABL.WeightedConstraintExperiment.UniformRefinement.paddedCloneValue}
\statementsource{Analysis-of-Boolean-Functions-by-Ryan-ODonnell}{pdf-ffbe9c7b3306-p202-b003}
\label{FABL.WeightedConstraintExperiment.UniformRefinement.paddedCloneValue}
\leanok
\uses{FABL.WeightedConstraintExperiment.UniformRefinement.PaddedCloneOccurrence, FABL.WeightedConstraintExperiment.UniformRefinement.paddedClonePayoff}
Let $E$ be a weighted constraint experiment with constraint template $[?]$ over
variables $V$, and let $R$ be a uniform refinement of $E$. For natural numbers $r,M$,
suppose every predicate selected by an outcome has arity at most $r$, and let
$a:V\times\{0,\ldots,M-1\}\to D$ be an assignment. The padded clone value of $a$ is the
expected padded clone payoff when a refined occurrence and a map from the $r$ positions
to the $M$ clones are sampled.
\end{definition}

\begin{theorem}[Padded-clone value equals clone value]
\lean{FABL.WeightedConstraintExperiment.UniformRefinement.paddedCloneValue_eq_clone_value}
\label{FABL.WeightedConstraintExperiment.UniformRefinement.paddedCloneValue_eq_clone_value}
\leanok
\uses{FABL.WeightedConstraintExperiment.UniformRefinement.paddedCloneValue, FABL.WeightedConstraintExperiment.clone, FABL.WeightedConstraintExperiment.clonePayoff, FABL.WeightedConstraintExperiment.clone_value_eq_iterated, FABL.WeightedConstraintExperiment.value, FABL.fintype_expect_prod, FABL.pmfExpectation_map, FABL.pmfExpectation_uniformPMF_eq_expect, FABL.uniformPaddedCloneExpectation}
\difficulty{3}
Let $E$ be a weighted constraint experiment with a uniform refinement over a finite
predicate template. For natural numbers $r,M$ with $0<M$, suppose every outcome's
predicate has arity at most $r$. For every assignment of values in $D$ to
$V\times\{0,\ldots,M-1\}$, the padded-clone value equals the value of the experiment
obtained by cloning each queried variable with an independently uniform index from
$\{0,\ldots,M-1\}$.
\end{theorem}

\begin{definition}[Good padded clone occurrence]
\lean{FABL.WeightedConstraintExperiment.UniformRefinement.PaddedCloneGood}
\statementsource{Analysis-of-Boolean-Functions-by-Ryan-ODonnell}{pdf-ffbe9c7b3306-p202-b003}
\label{FABL.WeightedConstraintExperiment.UniformRefinement.PaddedCloneGood}
\leanok
\uses{FABL.WeightedConstraintExperiment.UniformRefinement.PaddedCloneOccurrence, FABL.WeightedConstraintExperiment.UniformRefinement.paddedCloneScope}
Let $E$ be a weighted constraint experiment with template $\Psi$ and variables $V$, and
let $R$ be a uniform refinement. For natural numbers $r,M$, suppose every outcome's
predicate has arity at most $r$. A padded clone occurrence is good when the map sending
each predicate position to its original variable paired with its assigned clone index is
injective.
\end{definition}

\begin{definition}[Globally filtered clone occurrence]
\lean{FABL.WeightedConstraintExperiment.UniformRefinement.GloballyFilteredCloneOccurrence}
\statementsource{Analysis-of-Boolean-Functions-by-Ryan-ODonnell}{pdf-ffbe9c7b3306-p202-b003}
\label{FABL.WeightedConstraintExperiment.UniformRefinement.GloballyFilteredCloneOccurrence}
\leanok
\uses{FABL.WeightedConstraintExperiment.UniformRefinement.PaddedCloneGood, FABL.WeightedConstraintExperiment.UniformRefinement.PaddedCloneOccurrence}
Let $E$ be a weighted constraint experiment with template $\Psi$ and variables $V$, and
let $R$ be a uniform refinement. Given natural numbers $r,M$ such that every outcome's
predicate has arity at most $r$, a globally filtered clone occurrence is a padded clone
occurrence for which the map from predicate positions to original-variable and
clone-index pairs is injective.
\end{definition}

\begin{theorem}[Exists padded clone good]
\lean{FABL.WeightedConstraintExperiment.UniformRefinement.exists_paddedCloneGood}
\statementsource{Analysis-of-Boolean-Functions-by-Ryan-ODonnell}{pdf-ffbe9c7b3306-p202-b003}
\label{FABL.WeightedConstraintExperiment.UniformRefinement.exists_paddedCloneGood}
\leanok
\uses{FABL.WeightedConstraintExperiment.UniformRefinement.PaddedCloneGood, FABL.WeightedConstraintExperiment.UniformRefinement.PaddedCloneOccurrence}
\difficulty{3}
% dataset-pipeline: stub -- prose pending build_informal_statements.py
Statement pending regeneration.
\end{theorem}

\begin{definition}[Globally filtered clone constraint]
\lean{FABL.WeightedConstraintExperiment.UniformRefinement.globallyFilteredCloneConstraint}
\statementsource{Analysis-of-Boolean-Functions-by-Ryan-ODonnell}{pdf-ffbe9c7b3306-p202-b003}
\label{FABL.WeightedConstraintExperiment.UniformRefinement.globallyFilteredCloneConstraint}
\leanok
\uses{FABL.CSPConstraint, FABL.WeightedConstraintExperiment.UniformRefinement.GloballyFilteredCloneOccurrence, FABL.WeightedConstraintExperiment.UniformRefinement.paddedCloneScope}
Let $E$ be a weighted constraint experiment with template $\Psi$ and variables $V$, and
let $R$ be a uniform refinement. Given natural numbers $r,M$ and a bound saying every
outcome's predicate has arity at most $r$, each globally filtered clone occurrence
determines a constraint over variables $V\times\mathrm{Fin}(M)$: it uses the predicate
selected by the refined outcome, and its ordered scope sends each position to the
original variable paired with that position's assigned clone index.
\end{definition}

\begin{definition}[Globally filtered clone raw instance]
\lean{FABL.WeightedConstraintExperiment.UniformRefinement.globallyFilteredCloneRawCSP}
\statementsource{Analysis-of-Boolean-Functions-by-Ryan-ODonnell}{pdf-ffbe9c7b3306-p202-b003}
\label{FABL.WeightedConstraintExperiment.UniformRefinement.globallyFilteredCloneRawCSP}
\leanok
\uses{FABL.RawCSPInstance, FABL.WeightedConstraintExperiment.UniformRefinement.GloballyFilteredCloneOccurrence, FABL.WeightedConstraintExperiment.UniformRefinement.PaddedCloneGood, FABL.WeightedConstraintExperiment.UniformRefinement.PaddedCloneOccurrence, FABL.WeightedConstraintExperiment.UniformRefinement.exists_paddedCloneGood, FABL.WeightedConstraintExperiment.UniformRefinement.globallyFilteredCloneConstraint}
Let $E$ be a weighted constraint experiment with template $\Psi$ and variables $V$, and
let $R$ be a uniform refinement. Suppose $r\le M$ and every outcome's predicate has
arity at most $r$. The globally filtered clone raw CSP instance over
$V\times\mathrm{Fin}(M)$ consists of one constraint for each globally filtered clone
occurrence, and its multiset of constraints is nonempty.
\end{definition}

\begin{theorem}[Raw CSP value as a filtered expectation]
\lean{FABL.WeightedConstraintExperiment.UniformRefinement.globallyFilteredCloneRawCSP_value_eq_expectation}
\label{FABL.WeightedConstraintExperiment.UniformRefinement.globallyFilteredCloneRawCSP_value_eq_expectation}
\leanok
\uses{FABL.CSPConstraint.satisfiedIndicator, FABL.RawCSPInstance.value, FABL.WeightedConstraintExperiment.UniformRefinement.GloballyFilteredCloneOccurrence, FABL.WeightedConstraintExperiment.UniformRefinement.globallyFilteredCloneConstraint, FABL.WeightedConstraintExperiment.UniformRefinement.globallyFilteredCloneRawCSP, FABL.WeightedConstraintExperiment.UniformRefinement.paddedClonePayoff, FABL.WeightedConstraintExperiment.UniformRefinement.paddedCloneScope, FABL.WeightedConstraintExperiment.clonePayoff}
\difficulty{5}
Let $V$ be a set on which equality is decidable, and let $E$ be a weighted constraint
experiment with a uniform refinement over a finite predicate template. For natural
numbers $r,M$ with $r\le M$, suppose every outcome's predicate has arity at most $r$.
For any assignment of values in $D$ to $V\times\{0,\ldots,M-1\}$, its value on the raw
CSP instance formed from globally filtered clone occurrences, viewed as a real number,
equals the expected padded-clone payoff when the occurrence is sampled uniformly from
the globally filtered occurrences.
\end{theorem}

\begin{theorem}[Padded-clone failure implies tuple failure]
\lean{FABL.WeightedConstraintExperiment.UniformRefinement.paddedCloneFailure_le_tupleFailure}
\label{FABL.WeightedConstraintExperiment.UniformRefinement.paddedCloneFailure_le_tupleFailure}
\leanok
\uses{FABL.WeightedConstraintExperiment.UniformRefinement.PaddedCloneGood, FABL.WeightedConstraintExperiment.UniformRefinement.PaddedCloneOccurrence}
\difficulty{4}
Let $E$ be a weighted constraint experiment with a uniform refinement over a finite
predicate template, and let $r,M$ be natural numbers. Suppose every outcome's predicate
has arity at most $r$. For each padded clone occurrence, the indicator of failure of
injectivity of its padded clone scope is at most the indicator of failure of injectivity
of its clone-index tuple $\mathrm{Fin}(r)\to\mathrm{Fin}(M)$.
\end{theorem}

\begin{theorem}[Bound on padded-clone failure probability]
\lean{FABL.WeightedConstraintExperiment.UniformRefinement.paddedCloneFailureProbability_le}
\proofsource{Analysis-of-Boolean-Functions-by-Ryan-ODonnell}{pdf-ffbe9c7b3306-p202-b003}
\label{FABL.WeightedConstraintExperiment.UniformRefinement.paddedCloneFailureProbability_le}
\leanok
\uses{FABL.WeightedConstraintExperiment.UniformRefinement.PaddedCloneGood, FABL.WeightedConstraintExperiment.UniformRefinement.PaddedCloneOccurrence, FABL.WeightedConstraintExperiment.UniformRefinement.paddedCloneFailure_le_tupleFailure, FABL.cloneTupleCollisionProbability, FABL.cloneTupleCollisionProbability_le, FABL.fintype_expect_prod}
\difficulty{4}
Let $V$ be a set on which equality is decidable, and let $E$ be a weighted constraint
experiment with a uniform refinement over a finite predicate template. For natural
numbers $r,M$ with $0<M$, suppose every outcome's predicate has arity at most $r$. The
probability that a padded clone occurrence fails to have an injective padded clone scope
is at most $r^2/M$.
\end{theorem}

\begin{theorem}[Approximation of the clone value]
\lean{FABL.WeightedConstraintExperiment.UniformRefinement.abs_paddedCloneValue_sub_globallyFilteredCloneRawCSP_value_le}
\proofsource{Analysis-of-Boolean-Functions-by-Ryan-ODonnell}{pdf-ffbe9c7b3306-p202-b003}
\label{FABL.WeightedConstraintExperiment.UniformRefinement.abs_paddedCloneValue_sub_globallyFilteredCloneRawCSP_value_le}
\leanok
\uses{FABL.RawCSPInstance.value, FABL.WeightedConstraintExperiment.UniformRefinement.PaddedCloneGood, FABL.WeightedConstraintExperiment.UniformRefinement.exists_paddedCloneGood, FABL.WeightedConstraintExperiment.UniformRefinement.globallyFilteredCloneRawCSP, FABL.WeightedConstraintExperiment.UniformRefinement.globallyFilteredCloneRawCSP_value_eq_expectation, FABL.WeightedConstraintExperiment.UniformRefinement.paddedCloneFailureProbability_le, FABL.WeightedConstraintExperiment.UniformRefinement.paddedClonePayoff, FABL.WeightedConstraintExperiment.UniformRefinement.paddedCloneValue, FABL.WeightedConstraintExperiment.abs_uniformExpectation_subtype_le_failureProbability, FABL.WeightedConstraintExperiment.clonePayoff_mem_Icc}
\difficulty{3}
Let $V$ be a set on which equality is decidable, and let $E$ be a weighted constraint
experiment with a uniform refinement over a finite predicate template. For natural
numbers $r,M$ with $0<M$ and $r\le M$, suppose every outcome's predicate has arity at
most $r$. For every assignment of values in $D$ to the pairs in
$V\times\{0,\ldots,M-1\}$, the absolute difference between its padded-clone value and
its value on the raw CSP instance formed from globally filtered clone occurrences is at
most $r^2/M$.
\end{theorem}

\begin{definition}[Optimal padded clone assignment]
\lean{FABL.WeightedConstraintExperiment.UniformRefinement.paddedCloneOptimalAssignment}
\label{FABL.WeightedConstraintExperiment.UniformRefinement.paddedCloneOptimalAssignment}
\leanok
\uses{FABL.WeightedConstraintExperiment.UniformRefinement.paddedCloneValue}
Let $E$ be a weighted constraint experiment over a CSP template with domain $D$ and
variable set $V$, and let $R$ be a uniform refinement of $E$. Suppose $D$ is finite and
nonempty and $V$ is finite. For natural numbers $r$ and $M$, assume every predicate
selected by an outcome has arity at most $r$. An optimal padded clone assignment is an
assignment $(v,i)\\mapsto d$ from $V\\times\\mathrm{Fin}(M)$ to $D$ chosen to maximize
the padded clone value determined by $E$, $R$, $r$, $M$, and this arity bound.
\end{definition}

\begin{definition}[Padded clone optimum]
\lean{FABL.WeightedConstraintExperiment.UniformRefinement.paddedCloneOptimum}
\statementsource{Analysis-of-Boolean-Functions-by-Ryan-ODonnell}{pdf-ffbe9c7b3306-p202-b003}
\label{FABL.WeightedConstraintExperiment.UniformRefinement.paddedCloneOptimum}
\leanok
\uses{FABL.WeightedConstraintExperiment.UniformRefinement.paddedCloneOptimalAssignment, FABL.WeightedConstraintExperiment.UniformRefinement.paddedCloneValue}
Let $E$ be a weighted constraint experiment with template $\Psi$ and variables $V$, and
let $R$ be a uniform refinement. Suppose the domain $D$ is finite and nonempty, the
variable type $V$ is finite, and every outcome's predicate has arity at most $r$. The
padded clone optimum is the padded clone value of the optimal assignment to
$V\times\mathrm{Fin}(M)$.
\end{definition}

\begin{theorem}[Padded-clone value is bounded by its optimum]
\lean{FABL.WeightedConstraintExperiment.UniformRefinement.paddedCloneValue_le_optimum}
\label{FABL.WeightedConstraintExperiment.UniformRefinement.paddedCloneValue_le_optimum}
\leanok
\uses{FABL.WeightedConstraintExperiment.UniformRefinement.paddedCloneOptimum, FABL.WeightedConstraintExperiment.UniformRefinement.paddedCloneValue}
\difficulty{1}
Let $D$ and $V$ be finite sets, with $D$ nonempty, and let $E$ be a weighted constraint
experiment with a uniform refinement over a finite predicate template. For natural
numbers $r,M$, suppose every outcome's predicate has arity at most $r$. For every
assignment of values in $D$ to $V\times\{0,\ldots,M-1\}$, its padded-clone value is at
most the optimum padded-clone value.
\end{theorem}

\begin{theorem}[Optimum approximation by filtered clones]
\lean{FABL.WeightedConstraintExperiment.UniformRefinement.abs_paddedCloneOptimum_sub_globallyFilteredCloneRawCSP_optimum_le}
\proofsource{Analysis-of-Boolean-Functions-by-Ryan-ODonnell}{pdf-ffbe9c7b3306-p202-b003}
\label{FABL.WeightedConstraintExperiment.UniformRefinement.abs_paddedCloneOptimum_sub_globallyFilteredCloneRawCSP_optimum_le}
\leanok
\uses{FABL.RawCSPInstance.optimalAssignment, FABL.RawCSPInstance.optimum, FABL.RawCSPInstance.value, FABL.RawCSPInstance.value_le_optimum, FABL.WeightedConstraintExperiment.UniformRefinement.abs_paddedCloneValue_sub_globallyFilteredCloneRawCSP_value_le, FABL.WeightedConstraintExperiment.UniformRefinement.globallyFilteredCloneRawCSP, FABL.WeightedConstraintExperiment.UniformRefinement.paddedCloneOptimalAssignment, FABL.WeightedConstraintExperiment.UniformRefinement.paddedCloneOptimum, FABL.WeightedConstraintExperiment.UniformRefinement.paddedCloneValue, FABL.WeightedConstraintExperiment.UniformRefinement.paddedCloneValue_le_optimum}
\difficulty{5}
Let $D$ and $V$ be finite sets, with $D$ nonempty, and let $E$ be a weighted constraint
experiment with a uniform refinement over a finite predicate template. For natural
numbers $r,M$ with $0<M$ and $r\le M$, suppose every predicate selected by an outcome
has arity at most $r$. Then the absolute difference between the optimum padded-clone
value and the optimum value of the raw CSP instance formed from globally filtered clone
occurrences is at most $r^2/M$.
\end{theorem}

\begin{definition}[Globally filtered clone instance]
\lean{FABL.WeightedConstraintExperiment.UniformRefinement.globallyFilteredCloneCSP}
\statementsource{Analysis-of-Boolean-Functions-by-Ryan-ODonnell}{
  pdf-ffbe9c7b3306-p185-b004,
  pdf-ffbe9c7b3306-p202-b003
}
\label{FABL.WeightedConstraintExperiment.UniformRefinement.globallyFilteredCloneCSP}
\leanok
\uses{FABL.RawCSPInstance.removeIsolatedVariables, FABL.WeightedConstraintExperiment.UniformRefinement.globallyFilteredCloneRawCSP}
Let $E$ be a weighted constraint experiment with template $\Psi$ and variables $V$, and
let $R$ be a uniform refinement. Suppose $r\le M$ and every outcome's predicate has
arity at most $r$. The globally filtered clone CSP is the instance obtained from the
globally filtered clone raw CSP instance by removing its isolated variables.
\end{definition}

\begin{theorem}[Removing isolated variables preserves the optimum]
\lean{FABL.WeightedConstraintExperiment.UniformRefinement.globallyFilteredCloneCSP_optimum}
\label{FABL.WeightedConstraintExperiment.UniformRefinement.globallyFilteredCloneCSP_optimum}
\leanok
\uses{FABL.CSPInstance.optimum, FABL.RawCSPInstance.optimum, FABL.RawCSPInstance.optimum_removeIsolatedVariables, FABL.WeightedConstraintExperiment.UniformRefinement.globallyFilteredCloneCSP, FABL.WeightedConstraintExperiment.UniformRefinement.globallyFilteredCloneRawCSP}
\difficulty{1}
Let $D$ and $V$ be finite sets, with $D$ nonempty, and let $E$ be a weighted constraint
experiment with a uniform refinement over a finite predicate template. For natural
numbers $r,M$ with $r\le M$, suppose every outcome's predicate has arity at most $r$.
The optimum of the CSP obtained from the globally filtered clone raw instance by
restricting its variables to those appearing in a constraint scope equals the optimum of
that raw instance.
\end{theorem}

\begin{definition}[Globally filtered long code instance]
\lean{FABL.RegularUniqueGame.globallyFilteredLongCodeCSP}
\statementsource{Analysis-of-Boolean-Functions-by-Ryan-ODonnell}{
  pdf-ffbe9c7b3306-p191-b002,
  pdf-ffbe9c7b3306-p201-b007,
  pdf-ffbe9c7b3306-p202-b003
}
\label{FABL.RegularUniqueGame.globallyFilteredLongCodeCSP}
\leanok
\uses{FABL.RegularUniqueGame.longCodeAllowedWeightedExperiment, FABL.SignPredicate, FABL.SignTester, FABL.SignTester.UsesPredicates, FABL.WeightedConstraintExperiment.UniformRefinement, FABL.WeightedConstraintExperiment.UniformRefinement.globallyFilteredCloneCSP}
Let $G$ be a regular unique game with alphabet size $q$, let $T$ be a sign tester with a
finite seed set and a probability law, and let $\Psi$ be a finite set of sign predicates
used by every check selected by $T$. Let $R$ be a uniform refinement of the resulting
allowed long code weighted experiment. Given natural numbers $r,M$ with $r\le M$, and
assuming every check selected by a seed has predicate arity at most $r$, the globally
filtered long code CSP is the globally filtered clone CSP of that experiment and
refinement.
\end{definition}

\begin{theorem}[Filtered long-code CSP optimum bound]
\lean{FABL.RegularUniqueGame.beta_sub_edgeError_sub_collision_le_globallyFilteredLongCodeCSP_optimum}
\label{FABL.RegularUniqueGame.beta_sub_edgeError_sub_collision_le_globallyFilteredLongCodeCSP_optimum}
\leanok
\uses{FABL.BoundedCubeFunction.ofBoolean, FABL.CSPInstance.optimum, FABL.FiniteUniqueGame.toCSPInstance, FABL.RawCSPInstance.optimum, FABL.RegularUniqueGame.clonedLongCodeWeightedExperiment_value_constant, FABL.RegularUniqueGame.dictatorLongCode, FABL.RegularUniqueGame.globallyFilteredLongCodeCSP, FABL.RegularUniqueGame.longCodeAllowedWeightedExperiment, FABL.RegularUniqueGame.longCodeAllowedWeightedExperiment_clone_value, FABL.RegularUniqueGame.longCodeAssignment, FABL.RegularUniqueGame.longCodeExperimentValue, FABL.RegularUniqueGame.longCodeExperimentValue_optimalDictator_ge, FABL.RegularUniqueGame.optimalLabel, FABL.RegularUniqueGame.toFiniteUniqueGame, FABL.SignPredicate, FABL.SignTester, FABL.SignTester.UsesPredicates, FABL.SignTester.acceptanceProbability, FABL.WeightedConstraintExperiment.UniformRefinement, FABL.WeightedConstraintExperiment.UniformRefinement.abs_paddedCloneOptimum_sub_globallyFilteredCloneRawCSP_optimum_le, FABL.WeightedConstraintExperiment.UniformRefinement.globallyFilteredCloneCSP_optimum, FABL.WeightedConstraintExperiment.UniformRefinement.globallyFilteredCloneRawCSP, FABL.WeightedConstraintExperiment.UniformRefinement.paddedCloneOptimum, FABL.WeightedConstraintExperiment.UniformRefinement.paddedCloneValue, FABL.WeightedConstraintExperiment.UniformRefinement.paddedCloneValue_eq_clone_value, FABL.WeightedConstraintExperiment.UniformRefinement.paddedCloneValue_le_optimum, FABL.WeightedConstraintExperiment.constantCloneAssignment, FABL.dictator, FABL.signPredicateSetTemplate}
\difficulty{5}
Let $G$ be a regular unique game with label set $[q]$, where $[q]$ is nonempty, and let
a finite predicate family $\Psi$ contain every predicate used by a randomized tester.
Choose a uniform refinement $R$ of the tester’s induced weighted constraint experiment,
and integers $r,M$ with $0<M$ and $r\le M$, such that every tester predicate has arity
at most $r$. If every dictator check is accepted with probability at least $\beta$, and
the optimum value of the CSP associated with $G$ is at least $1-\delta$, then the
optimum value of the globally filtered long-code CSP built from $G$, the tester, $R$,
and the clone parameters $r,M$ is at least $\beta-r\delta-r^2/M$.
\end{theorem}

\begin{theorem}[Globally filtered long code csp optimum le of bounded soundness]
\lean{FABL.RegularUniqueGame.globallyFilteredLongCodeCSP_optimum_le_of_boundedSoundness}
\statementsource{Analysis-of-Boolean-Functions-by-Ryan-ODonnell}{
  pdf-ffbe9c7b3306-p191-b002,
  pdf-ffbe9c7b3306-p192-b003,
  pdf-ffbe9c7b3306-p201-b007
}
\label{FABL.RegularUniqueGame.globallyFilteredLongCodeCSP_optimum_le_of_boundedSoundness}
\leanok
\uses{FABL.CSPInstance.optimum, FABL.FiniteUniqueGame.toCSPInstance, FABL.RawCSPInstance.optimalAssignment, FABL.RawCSPInstance.optimum, FABL.RawCSPInstance.value, FABL.RegularUniqueGame.BoundedLongCodeTable, FABL.RegularUniqueGame.boundedLongCodeValue, FABL.RegularUniqueGame.cloneMeanLongCodeTable, FABL.RegularUniqueGame.clonedLongCode_value_eq_boundedLongCodeValue, FABL.RegularUniqueGame.globallyFilteredLongCodeCSP, FABL.RegularUniqueGame.longCodeAllowedWeightedExperiment, FABL.RegularUniqueGame.longCodeAllowedWeightedExperiment_clone_value, FABL.RegularUniqueGame.toFiniteUniqueGame, FABL.SignPredicate, FABL.SignTester, FABL.SignTester.UsesPredicates, FABL.WeightedConstraintExperiment.UniformRefinement, FABL.WeightedConstraintExperiment.UniformRefinement.abs_paddedCloneValue_sub_globallyFilteredCloneRawCSP_value_le, FABL.WeightedConstraintExperiment.UniformRefinement.globallyFilteredCloneCSP_optimum, FABL.WeightedConstraintExperiment.UniformRefinement.globallyFilteredCloneRawCSP, FABL.WeightedConstraintExperiment.UniformRefinement.paddedCloneValue, FABL.WeightedConstraintExperiment.UniformRefinement.paddedCloneValue_eq_clone_value, FABL.WeightedConstraintExperiment.clone, FABL.WeightedConstraintExperiment.value, FABL.signPredicateSetTemplate}
\difficulty{5}
% dataset-pipeline: stub -- prose pending build_informal_statements.py
Statement pending regeneration.
\end{theorem}

\begin{theorem}[Optimum bound for filtered long-code instance]
\lean{FABL.RegularUniqueGame.globallyFilteredLongCodeCSP_optimum_le_alpha_add_error_add_excess_add_collision}
\label{FABL.RegularUniqueGame.globallyFilteredLongCodeCSP_optimum_le_alpha_add_error_add_excess_add_collision}
\leanok
\uses{FABL.CSPInstance.optimum, FABL.FiniteUniqueGame.toCSPInstance, FABL.IsDictatorVsNoNotablesTest, FABL.RegularUniqueGame.excess_mul_epsilon_pow_five_div_eighteen_le_optimum, FABL.RegularUniqueGame.globallyFilteredLongCodeCSP, FABL.RegularUniqueGame.globallyFilteredLongCodeCSP_optimum_le_of_boundedSoundness, FABL.RegularUniqueGame.longCodeAllowedWeightedExperiment, FABL.RegularUniqueGame.toFiniteUniqueGame, FABL.SignPredicate, FABL.SignTester, FABL.WeightedConstraintExperiment.UniformRefinement}
\difficulty{3}
Let $G$ be a regular unique game with a nonempty finite label set, and let a tester
family satisfy the dictator-versus-no-notables conditions with parameters $\alpha$ and
$\beta$ and error function $\operatorname{error}$. Choose a uniform refinement of the
tester’s allowed weighted constraint experiment, and let $r,M$ be natural numbers with
$0<M$ and $r\le M$, such that every tester predicate has arity at most $r$. If
$0<\epsilon\le1$ and the optimum of the constraint satisfaction instance associated to
$G$ is less than $\tau\epsilon^5/18$, then the optimum of the globally filtered
long-code constraint satisfaction instance is at most
$\alpha+\operatorname{error}(\epsilon)+\tau+r^2/M$.
\end{theorem}

\begin{theorem}[Globally filtered long code csp optimum le of no influentials]
\lean{FABL.RegularUniqueGame.globallyFilteredLongCodeCSP_optimum_le_of_noInfluentials}
\statementsource{Analysis-of-Boolean-Functions-by-Ryan-ODonnell}{
  pdf-ffbe9c7b3306-p191-b002,
  pdf-ffbe9c7b3306-p192-b003,
  pdf-ffbe9c7b3306-p201-b007
}
\label{FABL.RegularUniqueGame.globallyFilteredLongCodeCSP_optimum_le_of_noInfluentials}
\leanok
\uses{FABL.CSPInstance.optimum, FABL.FiniteUniqueGame.toCSPInstance, FABL.IsDictatorVsNoInfluentialsTest, FABL.RegularUniqueGame.excess_mul_epsilon_div_eight_mul_q_sq_le_optimum_of_noInfluentials, FABL.RegularUniqueGame.globallyFilteredLongCodeCSP, FABL.RegularUniqueGame.globallyFilteredLongCodeCSP_optimum_le_of_boundedSoundness, FABL.RegularUniqueGame.longCodeAllowedWeightedExperiment, FABL.RegularUniqueGame.toFiniteUniqueGame, FABL.SignPredicate, FABL.SignTester, FABL.WeightedConstraintExperiment.UniformRefinement}
\difficulty{3}
% dataset-pipeline: stub -- prose pending build_informal_statements.py
Statement pending regeneration.
\end{theorem}

\begin{definition}[Finite Long-Code Test Boundary]
\lean{FABL.LongCodeTestBoundary}
\statementsource{Analysis-of-Boolean-Functions-by-Ryan-ODonnell}{pdf-ffbe9c7b3306-p191-b002}
\label{FABL.LongCodeTestBoundary}
\leanok
\uses{FABL.SignPredicate, FABL.SignTester, FABL.SignTester.UsesPredicates}
Let $q$ be a natural number, let $\Psi$ be a set of sign predicates, and let a sign
tester with alphabet size $q$ be given. A long-code test boundary consists of the
assumptions that $\Psi$ is finite and that every check selected by the tester uses a
predicate in $\Psi$.
\end{definition}

\begin{theorem}[Predicate Boundary of a Notables Test]
\lean{FABL.IsDictatorVsNoNotablesTest.longCodeTestBoundary}
\proofsource{Analysis-of-Boolean-Functions-by-Ryan-ODonnell}{pdf-ffbe9c7b3306-p191-b002}
\label{FABL.IsDictatorVsNoNotablesTest.longCodeTestBoundary}
\leanok
\uses{FABL.IsDictatorVsNoNotablesTest, FABL.LongCodeTestBoundary, FABL.SignPredicate, FABL.SignTester}
\difficulty{1}
Let $\Psi$ be a finite family of predicates on signs, let $\alpha,\beta\in\mathbb R$
satisfy $0<\alpha<\beta\le 1$, and let $e:[0,1]\to[0,1]$ tend to $0$ as its argument
tends to $0$ through positive values. Suppose a family of randomized testers, one for
each positive cube dimension, uses only predicates in $\Psi$; accepts every coordinate
dictator with probability at least $\beta$; and accepts every function whose stable
influences at parameter $1-\varepsilon$ are all at most $\varepsilon$ with probability
at most $\alpha+e(\varepsilon)$ for each $\varepsilon\in[0,1]$. For every positive
dimension, the tester's allowed predicate family is finite and every check selected by
any seed uses a predicate in that family.
\end{theorem}

\begin{theorem}[Predicate Boundary of a Dictator Test]
\lean{FABL.IsDictatorVsNoInfluentialsTest.longCodeTestBoundary}
\proofsource{Analysis-of-Boolean-Functions-by-Ryan-ODonnell}{pdf-ffbe9c7b3306-p191-b002}
\label{FABL.IsDictatorVsNoInfluentialsTest.longCodeTestBoundary}
\leanok
\uses{FABL.IsDictatorVsNoInfluentialsTest, FABL.LongCodeTestBoundary, FABL.SignPredicate, FABL.SignTester}
\difficulty{1}
Let $\Psi$ be a finite family of predicates on signs, let $\alpha,\beta\in\mathbb R$
satisfy $0<\alpha<\beta\le 1$, and let $e:\mathbb R\to\mathbb R$ take $[0,1]$ into
$[0,1]$ and tend to $0$ as its argument tends to $0$ through positive values. Suppose a
family of randomized testers, one for each positive cube dimension, uses only predicates
in $\Psi$; accepts every coordinate dictator with probability at least $\beta$; and
accepts every function with all stable influences at most $\varepsilon$ with probability
at most $\alpha+e(\varepsilon)$ for each $\varepsilon\in[0,1]$. For every positive
dimension, the tester's allowed predicate family is finite and every check selected by
any seed uses a predicate in that family.
\end{theorem}

\begin{definition}[Encoded Regular Unique Games Family]
\lean{FABL.EncodedRegularUniqueGamesFamily}
\statementsource{Analysis-of-Boolean-Functions-by-Ryan-ODonnell}{
  pdf-ffbe9c7b3306-p191-b002,
  pdf-ffbe9c7b3306-p192-b003,
  pdf-ffbe9c7b3306-p201-b007
}
\label{FABL.EncodedRegularUniqueGamesFamily}
\leanok
\uses{FABL.BinaryEncoding, FABL.RegularUniqueGame}
For a natural number $q$, an encoded regular Unique Games family consists of a type of
semantic instances, a binary encoding of those instances, and an assignment of a regular
Unique Games instance with alphabet size $q$ to each semantic instance. It also
specifies, for each semantic instance, a decidable equality relation on the vertices of
its associated game.
\end{definition}

\begin{definition}[Optimization Encoding]
\lean{FABL.EncodedRegularUniqueGamesFamily.OptimizationEncoding}
\statementsource{Analysis-of-Boolean-Functions-by-Ryan-ODonnell}{
  pdf-ffbe9c7b3306-p184-b004,
  pdf-ffbe9c7b3306-p186-b001
}
\label{FABL.EncodedRegularUniqueGamesFamily.OptimizationEncoding}
\leanok
\uses{FABL.CSPInstance.optimum, FABL.EncodedMaximizationProblem, FABL.FiniteUniqueGame.toCSPInstance, FABL.RegularUniqueGame.toFiniteUniqueGame}
Let $q$ be a natural number such that $\mathrm{Fin}(q)$ is nonempty, and let an encoded
regular Unique Games family with alphabet size $q$ be given. An optimization encoding
consists of an encoded maximization problem such that its valid inputs are exactly the
canonical encodings of semantic instances in the family, and such that for every
semantic instance its encoded optimum equals the optimum of the associated finite Unique
Games constraint satisfaction problem.
\end{definition}

\begin{definition}[Filtered Long-Code Optimum]
\lean{FABL.EncodedRegularUniqueGamesFamily.filteredLongCodeOptimum}
\statementsource{Analysis-of-Boolean-Functions-by-Ryan-ODonnell}{pdf-ffbe9c7b3306-p202-b003}
\label{FABL.EncodedRegularUniqueGamesFamily.filteredLongCodeOptimum}
\leanok
\uses{FABL.CSPInstance.optimum, FABL.LongCodeTestBoundary, FABL.RegularUniqueGame.globallyFilteredLongCodeCSP, FABL.RegularUniqueGame.longCodeAllowedWeightedExperiment, FABL.SignPredicate, FABL.SignTester, FABL.WeightedConstraintExperiment.UniformRefinement}
Let $q$ be a natural number, let an encoded regular Unique Games family with alphabet
size $q$ be given, and let a finite long-code test boundary for a sign tester with
alphabet size $q$ be given. For a semantic instance, a uniform refinement of its allowed
long-code weighted experiment, and natural numbers $r,M$ with $r\le M$, assume every
tester check has arity at most $r$. The filtered long-code optimum is the rational
optimum of the globally filtered long-code constraint satisfaction problem formed from
these data.
\end{definition}

\begin{theorem}[Equality of Filtered Long-Code Optima]
\lean{FABL.EncodedRegularUniqueGamesFamily.filteredLongCodeOptimum_eq}
\label{FABL.EncodedRegularUniqueGamesFamily.filteredLongCodeOptimum_eq}
\leanok
\uses{FABL.CSPInstance.optimum, FABL.EncodedRegularUniqueGamesFamily.filteredLongCodeOptimum, FABL.LongCodeTestBoundary, FABL.RegularUniqueGame.globallyFilteredLongCodeCSP, FABL.RegularUniqueGame.longCodeAllowedWeightedExperiment, FABL.SignPredicate, FABL.SignTester, FABL.WeightedConstraintExperiment.UniformRefinement}
\difficulty{1}
Let a family of encoded regular Unique Games, a finite set of allowed predicates, a
tester using only those predicates, and a semantic instance from the family be given.
Let $R$ be a uniform refinement of the associated weighted long-code experiment. For
natural numbers $r$ and $M$ with $r\le M$, suppose every tester predicate has arity at
most $r$. Then the filtered long-code optimum defined for this instance equals the
optimum of the corresponding globally filtered long-code constraint satisfaction
instance.
\end{theorem}

\begin{definition}[Long-Code Finite Transform]
\lean{FABL.EncodedRegularUniqueGamesFamily.LongCodeFiniteTransform}
\statementsource{Analysis-of-Boolean-Functions-by-Ryan-ODonnell}{
  pdf-ffbe9c7b3306-p191-b002,
  pdf-ffbe9c7b3306-p201-b007
}
\label{FABL.EncodedRegularUniqueGamesFamily.LongCodeFiniteTransform}
\leanok
\uses{FABL.LongCodeTestBoundary, FABL.RegularUniqueGame.longCodeAllowedWeightedExperiment, FABL.SignPredicate, FABL.SignTester, FABL.uniformPMF}
Let $q$ be a natural number, let an encoded regular Unique Games family with alphabet
size $q$ be given, and let a finite long-code test boundary for a sign tester with
alphabet size $q$ be given. A long-code finite transform specifies natural numbers $r$
and $M$ with $0<M$ and $r\le M$, such that every tester check has arity at most $r$. For
every semantic instance it also specifies a finite, nonempty set of refined occurrences,
a map from those occurrences to outcomes of the associated allowed long-code weighted
experiment, and the requirement that the uniform law on refined occurrences maps to that
experiment's law.
\end{definition}

\begin{definition}[Uniform Refinement]
\lean{FABL.EncodedRegularUniqueGamesFamily.LongCodeFiniteTransform.uniformRefinement}
\statementsource{Analysis-of-Boolean-Functions-by-Ryan-ODonnell}{pdf-ffbe9c7b3306-p202-b003}
\label{FABL.EncodedRegularUniqueGamesFamily.LongCodeFiniteTransform.uniformRefinement}
\leanok
\uses{FABL.EncodedRegularUniqueGamesFamily.LongCodeFiniteTransform, FABL.LongCodeTestBoundary, FABL.RegularUniqueGame.longCodeAllowedWeightedExperiment, FABL.SignPredicate, FABL.SignTester, FABL.WeightedConstraintExperiment.UniformRefinement}
Let $q$ be a natural number, let an encoded regular Unique Games family with alphabet
size $q$ be given, and let a long-code finite transform for a finite test boundary be
given. For each semantic instance, the transform's refined occurrence set, its map to
outcomes of that instance's allowed long-code weighted experiment, and the uniform-law
pushforward identity define a uniform refinement of that experiment.
\end{definition}

\begin{definition}[Long code output codec]
\lean{FABL.EncodedRegularUniqueGamesFamily.LongCodeOutputCodec}
\label{FABL.EncodedRegularUniqueGamesFamily.LongCodeOutputCodec}
\leanok
\uses{FABL.EncodedMaximizationProblem, FABL.EncodedRegularUniqueGamesFamily.LongCodeFiniteTransform, FABL.EncodedRegularUniqueGamesFamily.LongCodeFiniteTransform.uniformRefinement, FABL.EncodedRegularUniqueGamesFamily.filteredLongCodeOptimum, FABL.LongCodeTestBoundary, FABL.SignPredicate, FABL.SignTester}
Let a family of encoded regular unique games, a finite allowed family of sign
predicates, a tester using only predicates in that family, and a finite transform for
this boundary be given. A long code output codec consists of an encoded maximization
problem and a map sending each semantic instance in the game family to a Boolean list.
Every such list is a valid instance of the encoded problem, and the optimum of that
encoded instance equals the filtered long code optimum of the semantic instance under
the uniform refinement, with the transform's clone count and arity bound.
\end{definition}

\begin{definition}[Polynomial-time long-code transform]
\lean{FABL.EncodedRegularUniqueGamesFamily.LongCodeTransformMachine}
\label{FABL.EncodedRegularUniqueGamesFamily.LongCodeTransformMachine}
\leanok
\uses{FABL.EncodedRegularUniqueGamesFamily.LongCodeFiniteTransform, FABL.EncodedRegularUniqueGamesFamily.LongCodeOutputCodec, FABL.LongCodeTestBoundary, FABL.PolytimeMap, FABL.SignPredicate, FABL.SignTester}
Let a family of regular unique games, a finite family of allowed sign predicates, a
tester using only predicates from that family, and a finite long-code transform be
given, together with a binary codec for the transform’s finite CSP outputs. A
polynomial-time map from binary strings to binary strings is required to send the
canonical encoding of each semantic instance to exactly the codec’s encoding of that
instance’s output.
\end{definition}

\begin{definition}[Promise Reduction Gap Hypothesis]
\lean{FABL.EncodedRegularUniqueGamesFamily.GapHypothesis}
\statementsource{Analysis-of-Boolean-Functions-by-Ryan-ODonnell}{
  pdf-ffbe9c7b3306-p200-b007,
  pdf-ffbe9c7b3306-p201-b007
}
\label{FABL.EncodedRegularUniqueGamesFamily.GapHypothesis}
\leanok
\uses{FABL.BinaryPromiseProblem, FABL.EncodedMaximizationProblem.gapPromiseReal, FABL.EncodedRegularUniqueGamesFamily.OptimizationEncoding, FABL.PromiseKarpReducible}
Let $q$ be a natural number with $\mathrm{Fin}(q)$ nonempty, let an encoded regular
Unique Games family with alphabet size $q$ be given, and let an optimization encoding of
that family and a binary promise problem be given. A gap hypothesis consists of real
thresholds $s$ and $c$ with $s<c$, together with a promise-preserving Karp reduction
from the given promise problem to the real-valued Unique Games gap problem determined by
the encoding and these thresholds.
\end{definition}

\begin{theorem}[Long-code gap reduction]
\lean{FABL.EncodedRegularUniqueGamesFamily.longCodeFiniteTransform_reducesGap_of_boundedSoundness}
\proofsource{Analysis-of-Boolean-Functions-by-Ryan-ODonnell}{
  pdf-ffbe9c7b3306-p191-b002,
  pdf-ffbe9c7b3306-p192-b003,
  pdf-ffbe9c7b3306-p201-b007
}
\label{FABL.EncodedRegularUniqueGamesFamily.longCodeFiniteTransform_reducesGap_of_boundedSoundness}
\leanok
\uses{FABL.BoundedCubeFunction.ofBoolean, FABL.CSPInstance.optimum, FABL.EncodedMaximizationProblem.gapPromiseReal, FABL.EncodedRegularUniqueGamesFamily.LongCodeFiniteTransform, FABL.EncodedRegularUniqueGamesFamily.LongCodeFiniteTransform.uniformRefinement, FABL.EncodedRegularUniqueGamesFamily.LongCodeOutputCodec, FABL.EncodedRegularUniqueGamesFamily.LongCodeTransformMachine, FABL.EncodedRegularUniqueGamesFamily.OptimizationEncoding, FABL.EncodedRegularUniqueGamesFamily.filteredLongCodeOptimum_eq, FABL.FiniteUniqueGame.toCSPInstance, FABL.LongCodeTestBoundary, FABL.PromiseKarpReducible, FABL.RegularUniqueGame.BoundedLongCodeTable, FABL.RegularUniqueGame.beta_sub_edgeError_sub_collision_le_globallyFilteredLongCodeCSP_optimum, FABL.RegularUniqueGame.boundedLongCodeValue, FABL.RegularUniqueGame.globallyFilteredLongCodeCSP_optimum_le_of_boundedSoundness, FABL.RegularUniqueGame.toFiniteUniqueGame, FABL.SignPredicate, FABL.SignTester, FABL.SignTester.acceptanceProbability, FABL.dictator}
\difficulty{5}
Let $q>0$, let a finite family of encoded regular unique games with $q$ labels be
equipped with a finite family of sign predicates and a tester whose checks all use
predicates from that family. Suppose $eta$ is a lower bound on the tester's acceptance
probability on every coordinate dictator, and let an optimization encoding, a finite
long-code transform, an output codec, and a polynomial-time machine satisfy their stated
encoding and output properties. For real parameters
$	ext{ugSoundness}<	ext{ugCompleteness}$ and
$	ext{targetSoundness}<	ext{targetCompleteness}$, suppose
$1-	ext{δUG}	ext{ugCompleteness}$, $	ext{ugSoundness}<	ext{decodedLower}$,
$	ext{targetCompleteness}eta-r	ext{δUG}-r^2/M$, and $	ext{targetSoundness}$, where
the last two inequalities are respectively
$	ext{targetCompleteness}\le\beta-r\,\text{δUG}-r^2/M$ and
$\alpha+\mathrm{error}(\varepsilon)+\tau+r^2/M\le\text{targetSoundness}$. Assume also
that, for every semantic game instance and bounded long-code table, a long-code value at
least $\alpha+\mathrm{error}(\varepsilon)+\tau$ implies that the finite unique game's
optimum is at least $\text{decodedLower}$. Then the machine gives a polynomial-time Karp
reduction from the encoded unique-games promise problem with completeness threshold
$\text{ugCompleteness}$ and soundness threshold $\text{ugSoundness}$ to the encoded
target maximization promise problem with completeness threshold
$\text{targetCompleteness}$ and soundness threshold $\text{targetSoundness}$.
\end{theorem}

\begin{theorem}[Long code finite transform reduces gap]
\lean{FABL.EncodedRegularUniqueGamesFamily.longCodeFiniteTransform_reducesGap}
\statementsource{Analysis-of-Boolean-Functions-by-Ryan-ODonnell}{
  pdf-ffbe9c7b3306-p192-b003,
  pdf-ffbe9c7b3306-p201-b007
}
\label{FABL.EncodedRegularUniqueGamesFamily.longCodeFiniteTransform_reducesGap}
\leanok
\uses{FABL.EncodedMaximizationProblem.gapPromiseReal, FABL.EncodedRegularUniqueGamesFamily.LongCodeFiniteTransform, FABL.EncodedRegularUniqueGamesFamily.LongCodeOutputCodec, FABL.EncodedRegularUniqueGamesFamily.LongCodeTransformMachine, FABL.EncodedRegularUniqueGamesFamily.OptimizationEncoding, FABL.EncodedRegularUniqueGamesFamily.longCodeFiniteTransform_reducesGap_of_boundedSoundness, FABL.IsDictatorVsNoNotablesTest, FABL.IsDictatorVsNoNotablesTest.longCodeTestBoundary, FABL.PromiseKarpReducible, FABL.RegularUniqueGame.excess_mul_epsilon_pow_five_div_eighteen_le_optimum, FABL.SignPredicate, FABL.SignTester}
\difficulty{3}
% dataset-pipeline: stub -- prose pending build_informal_statements.py
Statement pending regeneration.
\end{theorem}

\begin{theorem}[Long code finite transform reduces gap of no influentials]
\lean{FABL.EncodedRegularUniqueGamesFamily.longCodeFiniteTransform_reducesGap_of_noInfluentials}
\statementsource{Analysis-of-Boolean-Functions-by-Ryan-ODonnell}{
  pdf-ffbe9c7b3306-p192-b003,
  pdf-ffbe9c7b3306-p201-b007
}
\label{FABL.EncodedRegularUniqueGamesFamily.longCodeFiniteTransform_reducesGap_of_noInfluentials}
\leanok
\uses{FABL.EncodedMaximizationProblem.gapPromiseReal, FABL.EncodedRegularUniqueGamesFamily.LongCodeFiniteTransform, FABL.EncodedRegularUniqueGamesFamily.LongCodeOutputCodec, FABL.EncodedRegularUniqueGamesFamily.LongCodeTransformMachine, FABL.EncodedRegularUniqueGamesFamily.OptimizationEncoding, FABL.EncodedRegularUniqueGamesFamily.longCodeFiniteTransform_reducesGap_of_boundedSoundness, FABL.IsDictatorVsNoInfluentialsTest, FABL.IsDictatorVsNoInfluentialsTest.longCodeTestBoundary, FABL.PromiseKarpReducible, FABL.RegularUniqueGame.excess_mul_epsilon_div_eight_mul_q_sq_le_optimum_of_noInfluentials, FABL.SignPredicate, FABL.SignTester}
\difficulty{3}
% dataset-pipeline: stub -- prose pending build_informal_statements.py
Statement pending regeneration.
\end{theorem}

\begin{theorem}[Conditional Unique Games Hardness]
\lean{FABL.EncodedRegularUniqueGamesFamily.conditional_uniqueGames_hardness}
\proofsource{Analysis-of-Boolean-Functions-by-Ryan-ODonnell}{pdf-ffbe9c7b3306-p192-b003}
\label{FABL.EncodedRegularUniqueGamesFamily.conditional_uniqueGames_hardness}
\leanok
\uses{FABL.BinaryPromiseProblem, FABL.EncodedMaximizationProblem.gapPromiseReal, FABL.EncodedRegularUniqueGamesFamily.GapHypothesis, FABL.EncodedRegularUniqueGamesFamily.LongCodeFiniteTransform, FABL.EncodedRegularUniqueGamesFamily.LongCodeOutputCodec, FABL.EncodedRegularUniqueGamesFamily.LongCodeTransformMachine, FABL.EncodedRegularUniqueGamesFamily.OptimizationEncoding, FABL.EncodedRegularUniqueGamesFamily.longCodeFiniteTransform_reducesGap, FABL.IsDictatorVsNoNotablesTest, FABL.IsDictatorVsNoNotablesTest.longCodeTestBoundary, FABL.PromiseKarpReducible, FABL.PromiseKarpReducible.trans, FABL.SignPredicate, FABL.SignTester}
\difficulty{1}
Let a family of encoded regular Unique Games, a source binary promise problem, and a
finite allowed family of sign predicates be given. Suppose a tester uses only allowed
predicates, has completeness at least $\beta$ on every dictator, and accepts every
function with no $(\varepsilon,\varepsilon)$-notable coordinate with probability at most
$\alpha+\operatorname{error}(\varepsilon)$; assume $0<\alpha<\beta\le1$, the error takes
values in $[0,1]$ on $[0,1]$ and tends to zero as $\varepsilon$ tends to zero through
positive values. Let $q>0$, and suppose the family has an optimization encoding and a
Unique Games gap reduction from the source promise problem, with thresholds
$s_{UG}<c_{UG}$. Suppose a finite long-code transformation for this tester has $M>0$
clones, $r\le M$, and the required arity bounds; suppose also that its output has an
encoding and a polynomial-time machine implementing that encoding on canonical inputs.
For real target thresholds $s_T<c_T$ and parameters $\delta_{UG},\varepsilon,\tau$ with
$0<\varepsilon\le1$, assume $1-\delta_{UG}\le c_{UG}$, $s_{UG}<\tau\varepsilon^5/18$,
$c_T\le\beta-r\delta_{UG}-r^2/M$, and
$\alpha+\operatorname{error}(\varepsilon)+\tau+r^2/M\le s_T$. Then the source promise
problem is polynomial-time Karp reducible to the target encoded maximization problem's
gap promise with soundness threshold $s_T$ and completeness threshold $c_T$.
\end{theorem}

\begin{theorem}[Conditional unique games hardness of no influentials]
\lean{FABL.EncodedRegularUniqueGamesFamily.conditional_uniqueGames_hardness_of_noInfluentials}
\statementsource{Analysis-of-Boolean-Functions-by-Ryan-ODonnell}{
  pdf-ffbe9c7b3306-p191-b002,
  pdf-ffbe9c7b3306-p192-b003
}
\label{FABL.EncodedRegularUniqueGamesFamily.conditional_uniqueGames_hardness_of_noInfluentials}
\leanok
\uses{FABL.BinaryPromiseProblem, FABL.EncodedMaximizationProblem.gapPromiseReal, FABL.EncodedRegularUniqueGamesFamily.GapHypothesis, FABL.EncodedRegularUniqueGamesFamily.LongCodeFiniteTransform, FABL.EncodedRegularUniqueGamesFamily.LongCodeOutputCodec, FABL.EncodedRegularUniqueGamesFamily.LongCodeTransformMachine, FABL.EncodedRegularUniqueGamesFamily.OptimizationEncoding, FABL.EncodedRegularUniqueGamesFamily.longCodeFiniteTransform_reducesGap_of_noInfluentials, FABL.IsDictatorVsNoInfluentialsTest, FABL.IsDictatorVsNoInfluentialsTest.longCodeTestBoundary, FABL.PromiseKarpReducible, FABL.PromiseKarpReducible.trans, FABL.SignPredicate, FABL.SignTester}
\difficulty{1}
% dataset-pipeline: stub -- prose pending build_informal_statements.py
Statement pending regeneration.
\end{theorem}

\begin{theorem}[Max e3 sat conditional unique games hardness]
\lean{FABL.maxE3Sat_conditional_uniqueGames_hardness}
\statementsource{Analysis-of-Boolean-Functions-by-Ryan-ODonnell}{
  pdf-ffbe9c7b3306-p192-b005,
  pdf-ffbe9c7b3306-p193-b002
}
\label{FABL.maxE3Sat_conditional_uniqueGames_hardness}
\leanok
\uses{FABL.BinaryPromiseProblem, FABL.EncodedMaximizationProblem.gapPromiseReal, FABL.EncodedRegularUniqueGamesFamily, FABL.EncodedRegularUniqueGamesFamily.GapHypothesis, FABL.EncodedRegularUniqueGamesFamily.LongCodeFiniteTransform, FABL.EncodedRegularUniqueGamesFamily.LongCodeOutputCodec, FABL.EncodedRegularUniqueGamesFamily.LongCodeTransformMachine, FABL.EncodedRegularUniqueGamesFamily.OptimizationEncoding, FABL.EncodedRegularUniqueGamesFamily.conditional_uniqueGames_hardness, FABL.IsDictatorVsNoNotablesTest.longCodeTestBoundary, FABL.PromiseKarpReducible, FABL.hastadThreeSatDegree, FABL.hastadThreeSatOXRError, FABL.hastadThreeSat_dictatorVsNoNotables, FABL.oxrE3SatReplacementError}
\difficulty{4}
% dataset-pipeline: stub -- prose pending build_informal_statements.py
Statement pending regeneration.
\end{theorem}

\begin{theorem}[Max e3 lin conditional unique games hardness]
\lean{FABL.maxE3Lin_conditional_uniqueGames_hardness}
\statementsource{Analysis-of-Boolean-Functions-by-Ryan-ODonnell}{pdf-ffbe9c7b3306-p193-b002}
\label{FABL.maxE3Lin_conditional_uniqueGames_hardness}
\leanok
\uses{FABL.BinaryPromiseProblem, FABL.EncodedMaximizationProblem.gapPromiseReal, FABL.EncodedRegularUniqueGamesFamily, FABL.EncodedRegularUniqueGamesFamily.GapHypothesis, FABL.EncodedRegularUniqueGamesFamily.LongCodeFiniteTransform, FABL.EncodedRegularUniqueGamesFamily.LongCodeOutputCodec, FABL.EncodedRegularUniqueGamesFamily.LongCodeTransformMachine, FABL.EncodedRegularUniqueGamesFamily.OptimizationEncoding, FABL.EncodedRegularUniqueGamesFamily.conditional_uniqueGames_hardness, FABL.IsDictatorVsNoNotablesTest.longCodeTestBoundary, FABL.PromiseKarpReducible, FABL.hastadSoundnessError, FABL.hastad_dictatorVsNoNotables}
\difficulty{3}
% dataset-pipeline: stub -- prose pending build_informal_statements.py
Statement pending regeneration.
\end{theorem}
